% Modernised reading — fonds Alexandre Grothendieck, shelfmark 161-2
% (pages 1-111, the whole folder).
%
% This is an INTERPRETATION, not a transcription. It re-reads the pages in
% today's notation and vocabulary, states the hypotheses the manuscript leaves
% implicit, and drops the critical apparatus so that the mathematics reads
% straight through. Everything the brackets and underlines carried — a doubtful
% reading, a notation he used differently, a missing passage — moves into a
% footnote.
%
% It is held to being mathematically correct as it stands, not merely faithful
% to the page. Where the manuscript is loose or mistaken, the true statement is
% what appears here, and the footnote says what the page has. In any dispute of
% substance the transcription governs: whoever cites Grothendieck must cite
% batch-0{1,...,6}.fr.tex.
%
% Pass: Fable 5.1 (claude-fable-5-1), 2026-09-05 — first pass, unchecked against
% the pages by a human. The folder was read whole, all six batches (pages
% 1-111) together, and this is a SINGLE document for the whole folder. The
% notation fixed for the document: Δ (or Cat_λ) for the 2-category of
% categories with limits of the prescribed types; T for a theory (« type »,
% « espèce de structure »), T(C) its models in C, R (= R_T, B_T) the small
% category representing it (« catégorie modelante », his « modelaire »,
% « modulaire »), S = T(Ens) its set-models, Σ = R° ⊂ S the representable
% ones, B = R̂ the classifying topos; in the ring notebook R is the category
% of affine schemes of finite type over Z, which is the same R for the theory
% of rings, and X_P the classifying subtopos of a ring property P. Substantive
% departures from the pages, each footnoted where it occurs: the well-poweredness
% hypothesis supplied to the SAFT-type representability claims (pp. 71-72, 74);
% a reflective subcategory is closed under limits, against p. 73; the six
% subtopoi of p. 24 are not a classification; the exponent n restored in the
% cones D_n of p. 42; the source of the descent equivalence on p. 34; the
% « i ∈ Z » and « ξ ∈ C_λ » slips of pp. 105-106. Announced and never
% established: « Pb 1 » and « Pb 2 » are named from p. 3 on and never stated;
% the operation K of the envelope theorem is enumerated (pp. 4-8, 27-28) and
% never written; Σ ≃ R° (pp. 68, 106) is asserted; X_car ≃ R̂_car (p. 49) is
% asserted; § 9 « Relation avec analyseurs de Lazard » (p. 89) has no text;
% the fair copy ①-⑩ stops mid-sentence on p. 106.
\documentclass[11pt,a4paper]{article}
% Le préambule commun à toutes les transcriptions du fonds Grothendieck.
%
% Il tient en une idée : **l'incertitude se compose, elle ne se résout pas.**
% Une transcription qui lisse un mot douteux a détruit la seule chose pour
% laquelle on la faisait — un lecteur doit pouvoir savoir, à chaque endroit,
% ce qui était sur le feuillet et ce qui a été ajouté. Les six macros qui
% suivent sont donc l'appareil critique, et non de la décoration.
%
% Les mêmes commandes sont reconnues par `scripts/render.mjs`, qui produit la
% vue de lecture du site. Une macro ajoutée ici doit l'être aussi là-bas,
% faute de quoi le rendu échouera bruyamment — ce qui est voulu : un rendu
% silencieusement faux est pire qu'un rendu absent.

\ProvidesPackage{grothendieck}[2026/08/08 Transcriptions du fonds Grothendieck]

\usepackage{amsmath,amssymb,amsthm}
\usepackage{mathrsfs}
\usepackage{stmaryrd}
% Les diagrammes commutatifs sont partout dans ce fonds : c'est la langue
% même de la géométrie algébrique de Grothendieck, et une transcription qui
% les rendrait en prose ne transcrirait rien.
\usepackage{tikz-cd}
% Français par défaut — c'est la langue des feuillets — mais l'anglais est
% chargé aussi : la traduction et le résumé sont des documents anglais, et la
% typographie française (espaces avant les deux-points) y serait une faute.
% Ces documents basculent par \selectlanguage{english} en tête de corps.
\usepackage[english,french]{babel}
\usepackage{fontspec}
\usepackage{geometry}
\geometry{a4paper,margin=25mm,marginparwidth=28mm}
\usepackage{marginnote}
\usepackage{xcolor}
% ulem, not soul. `soul`'s \st and \ul are built on a character-by-character
% scan that XeLaTeX's font machinery breaks: the document fails with an
% undefined control sequence rather than degrading. ulem does the same two jobs
% and is Unicode-engine safe. `normalem` stops it hijacking \emph, which the
% transcriptions use for Grothendieck's own emphasis.
\usepackage[normalem]{ulem}
\usepackage{hyperref}
\hypersetup{colorlinks=true,linkcolor=black,urlcolor=black,pdfborder={0 0 0}}

% --- Métadonnées du lot -----------------------------------------------------
% Reprises telles quelles par le script de rendu, qui les affiche en tête de
% la vue de lecture. Elles fixent ce que le fichier prétend couvrir : sans
% elles, une transcription flotte sans point d'attache dans un fonds de seize
% mille feuillets.

\newcommand{\folder}[1]{\gdef\grothFolder{#1}}
\newcommand{\batch}[1]{\gdef\grothBatch{#1}}
\newcommand{\leaves}[2]{\gdef\grothFirst{#1}\gdef\grothLast{#2}}
\newcommand{\foldertitle}[1]{\gdef\grothFoldertitle{#1}}
\gdef\grothFolder{?}\gdef\grothBatch{?}\gdef\grothFirst{?}\gdef\grothLast{?}\gdef\grothFoldertitle{}

% --- Le résumé --------------------------------------------------------------
%
% Il ouvre la lecture modernisée, sous le titre « Résumé », et il est composé
% en petit. Ce n'est pas un ornement : c'est le seul passage du document qui
% ne soit pas des mathématiques, et un lecteur doit pouvoir le reconnaître
% avant de l'avoir lu — puis le sauter s'il connaît déjà le sujet, ou s'y
% arrêter s'il ne le connaît pas. Le filet qui le suit marque la reprise du
% corps.

\newenvironment{resume}
  {\small\setlength{\parskip}{0.45em}}
  {\par\medskip\hrule\medskip}

% --- Mots-clés ---------------------------------------------------------------
%
% En anglais, en clôture du résumé de la lecture modernisée : le vocabulaire
% moderne sous lequel chercher ce que les feuillets construisent. Le manifeste
% du site (`npm run manifest`) les extrait pour étiqueter la cote entière —
% cette ligne est la source unique des tags, il n'y a pas de fichier à côté.
\newcommand{\keywords}[1]{\par\smallskip\noindent
  {\footnotesize\textsc{Keywords} --- \textit{#1}}\par}

% --- L'appareil critique ----------------------------------------------------

% Le numéro de feuillet, en marge. C'est l'ancre qui permet de régler un
% désaccord sur une formule en regardant *un* feuillet plutôt qu'une chemise
% de six cents. Le site s'en sert aussi pour tourner les pages du fac-similé
% quand on fait défiler la transcription.
\newcommand{\leaf}[1]{%
  \marginnote{\footnotesize\color{gray!70}\textsf{#1}}%
  \label{leaf:#1}\ignorespaces}

% Illisible : on ne devine pas. Un mot inventé qui se lit comme les autres est
% le pire résultat possible.
\newcommand{\ill}{\textcolor{red!60!black}{[\,\ldots\,]}}

% Lecture douteuse : le mot est donné, et signalé comme incertain.
\newcommand{\uncertain}[1]{\uline{#1}}

% Ajout de l'éditeur — une lettre manquante, un mot sous-entendu.
\newcommand{\add}[1]{\textcolor{blue!50!black}{[#1]}}

% Biffé par Grothendieck lui-même. Ce qu'il a rayé fait partie du document :
% une rature dit souvent où la pensée a changé de direction.
\newcommand{\struck}[1]{\sout{#1}}

% Note du transcripteur, sur l'état matériel du feuillet.
\newcommand{\note}[1]{\par\smallskip{\small\color{gray!60!black}\textit{[#1]}}\par\smallskip}

% Note marginale de Grothendieck, distincte du corps du texte.
\newcommand{\marginal}[1]{\par\smallskip\noindent\hspace{1em}%
  {\small\color{gray!60!black}#1}\par\smallskip}

% --- Titre ------------------------------------------------------------------

\AtBeginDocument{%
  \begin{center}
    {\large\bfseries Fonds Alexandre Grothendieck --- cote n\textsuperscript{o}~\grothFolder}\\[2pt]
    {\itshape\grothFoldertitle}\\[4pt]
    {\small feuillets \grothFirst--\grothLast\ (lot \grothBatch)}\\[2pt]
    {\footnotesize\color{gray!60!black}Transcription établie d'après le fac-similé numérisé
     par l'Université de Montpellier.\\
     Les lectures incertaines sont soulignées, les illisibles notées [\,\ldots\,],
     les ajouts entre crochets.}
  \end{center}
  \vspace{1em}}

\endinput


\folder{161-2}
\pages{1}{111}
\dating{[vers 1963-1973]}
\watermark{Édition de démonstration\\interprétation personnelle de l'œuvre}
\foldertitle{Algèbre universelle [ou catégories] : notes manuscrites (s.d.) — lecture modernisée du dossier entier}

\begin{document}

\section*{Résumé}

\begin{resume}
Une chemise dont l'onglet porte, de sa main, « Algèbre universelle », et
que l'inventaire hésite à ranger sous « catégories » : cent onze pages sans
date, où l'on reconnaît quatre ou cinq campagnes d'écriture — des feuillets
lignés de brouillon, un cahier à spirale quadrillé, des feuillets jaunis
d'une autre plume, et pour finir une copie au net paginée de ① à ⑩. Ce ne
sont pas des notes de lecture ; c'est une théorie en train de se faire, et
qui se refait quatre fois.

De quoi s'agit-il ? L'algèbre universelle, au sens classique, étudie ce que
les structures algébriques — groupes, anneaux, modules, treillis — ont en
commun : chacune est donnée par des opérations et des axiomes, et l'on peut
raisonner sur « une structure quelconque » sans dire laquelle. La question
de ces pages est la version que ce mot prend chez quelqu'un qui vient de
passer dix ans à faire de la géométrie \emph{dans} des catégories : que
veut dire « une structure algébrique dans une catégorie $C$ » — un groupe
dans les faisceaux sur un espace, un anneau dans un topos —, et qu'est-ce
qu'une \emph{espèce} de structure, prise indépendamment de la catégorie où
on la réalise ?

L'idée qui arrive est d'une simplicité qui la rend difficile à voir. Une
espèce de structure est une \emph{recette} : elle demande certaines
limites — des produits pour écrire une loi de composition, des égalisateurs
pour écrire un axiome — et c'est tout. La recette elle-même est une petite
catégorie $R$, celle des « constructions » que l'on peut faire à partir
d'une structure de l'espèce ; un plat, c'est-à-dire une structure dans $C$,
est un foncteur $R \to C$ qui respecte les limites en jeu. Toute l'algèbre
universelle tient alors dans un dictionnaire, que le dossier écrit
littéralement sous forme de tables (pages 90 et 91) : aux espèces de
structures correspondent les catégories $R$, aux sous-espèces les
localisations de $R$, aux structures « plus riches » les foncteurs qui
engendrent, aux limites d'espèces les limites inductives de catégories.

Deux faits font de ce dictionnaire autre chose qu'un changement de
vocabulaire, et les pages les découvrent l'un après l'autre. Le premier :
\emph{les modèles déterminent la théorie}. Si l'on connaît la catégorie $S$
de toutes les structures ensemblistes de l'espèce — tous les groupes, avec
leurs homomorphismes —, on retrouve la recette $R$ à l'intérieur de $S$,
comme la sous-catégorie des structures « de présentation finie », celles
que l'on peut décrire par un nombre fini de générateurs et de relations
(pages 77 et 78). La théorie n'est donc pas un choix de présentation : elle
est une propriété de ses modèles. Le second : \emph{il existe un modèle
universel}. Pour chaque espèce il y a un topos — une catégorie de faisceaux
sur un site convenable — qui contient une structure de l'espèce, et telle
que toute structure de l'espèce, dans n'importe quel topos, s'en déduise par
une unique image inverse. C'est le \emph{topos classifiant} de l'espèce, et
le cahier à spirale en fait la liste pour les anneaux : le topos des
préfaisceaux sur les schémas affines classifie les anneaux, le topos de
Zariski les anneaux locaux, le topos étale les anneaux strictement
locaux, et ainsi de suite pour vingt-deux propriétés, réduit, intègre,
corps, corps algébriquement clos, de caractéristique donnée. L'anneau
« universel » y est le faisceau structural, et l'on peut lire l'ensemble du
cahier comme la réponse à une seule question : \emph{quel est le corps
universel} — et pourquoi n'est-il pas un anneau universel comme les
autres, parce que « être un corps » demande une disjonction, « nul ou
inversible », que les limites seules ne savent pas dire (pages 92 à 95).

Ces pages ne citent presque personne — Lazard pour ses « analyseurs »,
Lawvere une fois, Schanuel pour une idée, et SGA 4 pour tout le reste. Ce
qu'elles construisent porte aujourd'hui d'autres noms : les théories de
Lawvere et les esquisses d'Ehresmann pour les recettes, les catégories
localement présentables et la dualité de Gabriel--Ulmer pour « les modèles
déterminent la théorie », les topos classifiants et la logique géométrique
pour le modèle universel. Toutes ces théories datent des années 1963 à
1977, c'est-à-dire de la fourchette que l'inventaire propose pour le
dossier ; la lecture qui suit dit, à chaque endroit, ce que la page
établit, ce qu'elle affirme sans preuve, et où la littérature l'a
retrouvé.

Ce qui tient lieu de conclusion, page 80, est une liste des cas à traiter,
et la phrase : « Je m'attends à des ennuis (de sortie de l'Univers) ». Le
mot revient : le « grain de sel » de la page 84, « on n'a pas encore prouvé
que cela marche » de la page 111. La copie au net s'arrête au milieu d'une
phrase. C'est un dossier où l'on voit, à la lettre, une notion arriver et
se stabiliser — quatre énoncés successifs du même théorème d'enveloppe,
pages 3, 15, 26 et 97 —, et où la prudence sur la taille des objets n'est
pas une manie mais le lieu exact où le problème résiste.

\keywords{universal algebra, sketch, essentially algebraic theory, Lawvere theory, locally presentable category, accessible category, Gabriel–Ulmer duality, dense functor, free completion, transfinite construction, classifying topos, geometric theory, cartesian theory, coherent topos, enough points, subtopos, Zariski topos, local ring, von Neumann regular ring, generic field, fibre product of topoi, monadicity, descent, orthogonality}
\end{resume}

\pagerange{3}{3}
\section*{Le fil du dossier, et les conventions}

Le dossier n'est pas relié dans l'ordre de sa composition, mais ses
campagnes se distinguent au papier et à la plume, et l'argument, lui, est
un. Voici ses stations, dans l'ordre des feuillets.

\begin{itemize}
\item \textbf{Pages 3 à 15.} Le \emph{théorème de la $\Delta$-enveloppe} :
une petite catégorie munie de cônes se complète librement en une catégorie
ayant des limites de types prescrits ; démonstration par récurrence
transfinie ; extension aux conditions d'exactitude ; une typologie de
celles-ci. Repris pages 26 à 31 sous sa forme la plus générale.
\item \textbf{Pages 17 à 24.} Le cas des \emph{topos} : une théorie à
limites finies a pour modèle universel un topos de préfaisceaux, ses
sous-théories sont les sous-topos, et un objet d'une 2-catégorie de « bons
topos » représente une théorie constructible si et seulement s'il est un
topos de Grothendieck (Théorème 4). Produits fibrés et limites transfinies
de topos. Le topos classifiant d'un objet.
\item \textbf{Pages 33 et 34.} Deux feuilles de travail : théories de
Lawvere, monadicité, descente.
\item \textbf{Pages 36 à 50.} Le \emph{cahier des anneaux} : vingt-deux
propriétés d'anneaux et le sous-topos du classifiant qui les classifie ;
la question des points ; le recollement de deux sous-topos.
\item \textbf{Pages 51 et 52.} Quelles propriétés de modules et d'algèbres
sont « typiques » (géométriques).
\item \textbf{Pages 54 à 61.} L'outillage : \emph{catégories accessibles}
d'après SGA 4 I § 9, foncteurs strictement générateurs, et la
caractérisation des catégories localement présentables.
\item \textbf{Pages 62 à 68.} « Compléments » : théories sur $\Delta$,
constructibilité, changement de $\Delta$, limites dans $T(C)$, et le
passage aux préfaisceaux $R^{\circ} \subset S \subset \widehat{R^{\circ}}$.
\item \textbf{Pages 69 à 80.} « Exemples » : les théories définies par des
limites — toutes les limites, les limites d'un type $d'$, les limites
finies — et pour ces dernières la dualité complète entre théories et
catégories de modèles. La liste des cas à traiter.
\item \textbf{Pages 81 à 87.} L'adjonction entre théories et catégories,
la catégorie des constructions $B_{T}$, sa construction explicite pas à
pas, et la réciproque : toute catégorie à générateurs est classifiante.
\item \textbf{Pages 89 à 91.} Les \emph{dictionnaires}.
\item \textbf{Pages 92 à 95.} Le \emph{corps universel}.
\item \textbf{Pages 97 à 106.} La \emph{copie au net} ① à ⑩ :
$\lambda$-types, représentabilité, règles de formation, dualité, et le
cas des limites projectives.
\item \textbf{Pages 108 à 111.} Le classifiant depuis le topos : image
inverse de la structure universelle, image inverse d'un champ, et le topos
modulaire relatif.
\end{itemize}

\textbf{Conventions, valables pour tout le document.} $\Delta$, ou
$\mathrm{Cat}_{\lambda}$ dans la copie au net, est la 2-catégorie des
catégories possédant les limites de types prescrits $d = (d', d'')$ (resp.
$\lambda = (\overleftarrow{\lambda}, \overrightarrow{\lambda})$), avec les
foncteurs qui y commutent. Une \emph{théorie} (le manuscrit dit aussi
\emph{type}, \emph{espèce de structure}) est un 2-foncteur $T : \Delta \to
\mathrm{Cat}$, $T(C)$ étant la catégorie des structures de l'espèce dans
$C$ ; elle est \emph{représentable} par $R \in \mathrm{Ob}\,\Delta$ si $T(C)
\simeq \underline{\mathrm{Hom}}_{\Delta}(R, C)$, et $R$ — que le manuscrit
note $R$, $R_{T}$, $R^{\lambda}_{T}$, $B_{T}$, et nomme « catégorie des
constructions », « modelaire », « modulaire », « classifiante » selon les
feuillets — est appelé ici sa \emph{catégorie modelante}. Elle est
\emph{algébrique} (\emph{constructible}) si elle s'obtient de la structure
vide sur des objets de base par des ajouts de flèches, des égalités et des
inversions, éventuellement transfinis. $S = T(\mathrm{Ens})$ est sa catégorie
de modèles ensemblistes, $\Sigma = R^{\circ} \subset S$ la sous-catégorie
des modèles représentables, et, pour une théorie à limites finies, $B =
\widehat{R}$ son topos classifiant. Le \emph{topos} est toujours un topos
de Grothendieck, sauf dans « bon topos » (page 17), où la condition de
taille est omise. Dans le cahier des anneaux, $R$ désigne la catégorie des
schémas affines de type fini sur $\mathbb{Z}$ : c'est la catégorie
modelante de la théorie des anneaux, et il n'y a donc pas conflit de
notation ; $A$ y est l'anneau universel, $X_{P}$ le sous-topos classifiant
la propriété $P$, $\widetilde{R}_{\mathrm{top}}$ les faisceaux pour une
topologie. « Pb 1 » et « Pb 2 » désignent la $\Delta$-enveloppe d'une
catégorie et la localisation dans $\Delta$ ; ils sont nommés dès la page 3
et jamais énoncés, et le sens retenu est justifié à cet endroit. Le
2-foncteur $R \mapsto T_{R}$ est contravariant : un morphisme de théories
$T' \to T$ est un $\Delta$-foncteur $R \to R'$, et le foncteur d'oubli
$T'(\mathrm{Ens}) \to T(\mathrm{Ens})$ va dans le sens de $T' \to T$.

\textbf{Ce que le dossier annonce et n'établit pas.} L'opération $K$ de la
construction transfinie est énumérée par ses propriétés (pages 4 à 8, 27
et 28) et jamais écrite ; le renvoi (*) de la page 31 s'arrête après trois
lignes. L'équivalence $\Sigma \simeq R^{\circ}$ entre les modèles
« représentables en toute construction » et la catégorie modelante est
affirmée (pages 68 et 106) et démontrée seulement pour les limites finies
(page 77). L'identification $X_{\mathrm{car}} \simeq \widehat{R_{\mathrm{car}}}$
(page 49) est posée sans preuve. Le § 9, « Relation avec analyseurs de
Lazard », n'a qu'un titre. Le point (4) de la page 68 et le crochet de la
page 66 sont abandonnés après quelques mots, et la copie au net s'arrête
page 106 au milieu de sa dernière phrase. Enfin les « six sous-topos » de
la page 24 ne sont pas tous les sous-topos du classifiant d'un objet, et
la note le dit.
\pagerange{3}{15}
\section*{Le théorème de la $\Delta$-enveloppe (pages 3 à 15, reprise pages 26 à 31)}

\pagerange{3}{3}
\subsection*{Le cadre}

On fixe un couple $d = (d', d'')$ de petits ensembles de \emph{types de
diagrammes} : chaque élément de $d'$ (resp. $d''$) est une petite catégorie
d'indices, et l'on demandera les limites projectives (resp. inductives)
indexées par ces catégories. On note $\Delta = \Delta_{d}$ la 2-catégorie
dont les objets sont les catégories possédant les limites projectives de
type $d'$ et les limites inductives de type $d''$, dont les 1-flèches sont
les foncteurs commutant à ces limites (les \emph{$\Delta$-foncteurs}), et
dont les 2-flèches sont les transformations naturelles.\footnote{Le
manuscrit écrit « la catégorie des catégories ayant des $\lim$ de type $d$,
avec foncteurs y commutant ». Que $\Delta$ soit une 2-catégorie et non une
catégorie est constant dans tout le dossier — toutes les propriétés
universelles y sont énoncées « à équivalence près » — et nous le disons
d'emblée.}

Une catégorie $C$ étant donnée, un \emph{paquet de cônes} $\delta =
(\delta', \delta'')$ sur $C$ est un petit ensemble de cônes projectifs
$(x \to x_{\alpha})_{\alpha}$ de type $\in d'$ et de cônes inductifs de
type $\in d''$. Un foncteur $u : C \to E$ vers un objet $E$ de $\Delta$
est \emph{$\delta$-exact} s'il envoie chaque cône de $\delta$ sur un cône
limite de $E$. On note $\underline{\mathrm{Hom}}_{\delta}(C, E)$ la
sous-catégorie pleine de $\underline{\mathrm{Hom}}(C, E)$ formée des
foncteurs $\delta$-exacts, et $\underline{\mathrm{Hom}}_{\Delta,\delta}(C,E)$
la même chose lorsque $C$ est déjà un objet de $\Delta$ et que l'on
demande de plus que $u$ soit un $\Delta$-foncteur.

En langage d'aujourd'hui, le couple $(C, \delta)$ est une \emph{esquisse}
(une esquisse à limites projectives et inductives, au sens d'Ehresmann),
et le théorème qui suit est l'existence de la catégorie
\emph{$\Delta$-libre engendrée par l'esquisse}.\footnote{Ehresmann
introduit les esquisses en 1966--1968 ; la notion est dans l'air à la date
que l'inventaire propose pour ces feuillets, et rien ici ne la cite. Le mot
« enveloppe » est celui du manuscrit, et nous le gardons partout.}

\pagerange{3}{3}
\subsection*{L'énoncé}

\textbf{Théorème.} \emph{Soient $d$, $\Delta = \Delta_{d}$ comme ci-dessus,
$C$ une catégorie essentiellement petite et $\delta$ un paquet de cônes sur
$C$. Il existe un objet $\overline{C}_{\delta}$ de $\Delta$ et un foncteur
$i : C \to \overline{C}_{\delta}$ tels que}
\begin{enumerate}
\item[a)] \emph{$i$ est $\delta$-exact ;}
\item[b)] \emph{pour tout $E \in \mathrm{Ob}\,\Delta$, le foncteur
$\varphi \mapsto \varphi \circ i$ induit une équivalence de catégories}
\[
\underline{\mathrm{Hom}}_{\Delta,\, i(\delta)}(\overline{C}_{\delta}, E)
\;\xrightarrow{\ \approx\ }\;
\underline{\mathrm{Hom}}_{\delta}(C, E).
\]
\end{enumerate}
\emph{Autrement dit, $(\overline{C}_{\delta}, i)$ est un objet 2-initial de
la 2-catégorie des objets $E$ de $\Delta$ munis d'un foncteur $\delta$-exact
$C \to E$. De plus $\overline{C}_{\delta}$ est essentiellement
petite.}\footnote{La petitesse est essentielle et elle tient à ce que $d$ est
un \emph{ensemble} de types : la construction s'arrête à un ordinal
$\omega$ dont le cardinal majore les cardinaux des catégories d'indices en
jeu. Si l'on admettait tous les petits diagrammes, la conclusion serait
fausse : la complétion libre d'une petite catégorie par toutes les petites
limites projectives n'est pas petite (celle de la catégorie ponctuelle par
les produits quelconques est $(\mathrm{Ens})^{\circ}$). Le manuscrit y
revient page 76 : « supposons $d'$ petit, alors une $R$ admettant une petite
sous-catégorie $d'$-cogénératrice est petite ». Dans l'exemple des pages 70
à 73, où toutes les limites sont admises, l'enveloppe est
$\widehat{I^{\circ}}{}^{\circ}$ et n'est pas petite ; voir ci-dessous.}

Le \emph{N.B.} de la page 3 dit ce que l'énoncé unifie : le « problème 1 »
(pas de cônes du tout, $\delta' = \delta'' = \emptyset$ : la
$\Delta$-enveloppe d'une catégorie nue) et le « problème 2 » ($C$ déjà dans
$\Delta$, et $\delta$ formé des cônes limites de $C$ auxquels on adjoint des
flèches $(x \to y)$ d'un ensemble $\Sigma$ vues comme cônes à un seul
sommet : c'est la localisation $C\Sigma^{-1}$ prise dans $\Delta$). Ces deux
problèmes sont nommés dès ici et jamais énoncés dans le dossier ; nous les
prenons dans ce sens, qui est le seul compatible avec toutes leurs
occurrences.\footnote{« Pb 1 » et « Pb 2 » reviennent pages 17, 19, 20, 69
à 75. Le problème 1 est l'adjonction $\mathrm{Cat} \rightleftarrows \Delta$
de la page 82 ; le problème 2 est la représentabilité du 2-foncteur
$C \mapsto \underline{\mathrm{Hom}}_{\Delta}(J, C; \Sigma)$ de la même
page.}

\pagerange{3}{6}
\subsection*{La construction transfinie}

La démonstration est une récurrence transfinie qui est devenue la méthode
standard pour ce genre d'énoncé.\footnote{La forme générale en est donnée
par Kelly, \emph{A unified treatment of transfinite constructions} (1980) ;
la construction ici est celle-là, pour un cas particulier, écrite une
douzaine d'années plus tôt.} On construit une suite transfinie de petites
catégories sous $C$, chacune munie d'un paquet de cônes :
\[
\begin{cases}
(C_{0}, \delta_{0}) = (C, \delta), \\
(C_{i+1}, \delta_{i+1}) = K(C_{i}, \delta_{i}), \qquad
\tau_{i} : C_{i} \to C_{i+1}, \\
C_{i} = \varinjlim_{j < i} C_{j} \ \text{(limite dans } \mathrm{Cat}),
\quad \delta_{i} = \bigcup_{j<i} \mathrm{Im}\,\delta_{j}
\ \text{ si $i$ est un ordinal limite,}
\end{cases}
\]
où $K$ est une « opération » à expliciter, et l'on note $\alpha_{i} : C \to
C_{i}$. L'opération $K$ doit avoir quatre propriétés, énoncées page 4 pour
le pas $C_{i} \to C_{i+1}$ :

\begin{enumerate}
\item[a)] $\tau_{i}$ est \emph{$\delta_{i}$-préexact}, ce qui signifie deux
choses pour chaque cône projectif $(x \to x_{\alpha}) \in \delta'_{i}$ (et
dualement pour les cônes inductifs) : deux flèches $u, v : z \rightrightarrows
x$ égalisées par toutes les projections deviennent égales dans $C_{i+1}$ ;
et toute famille compatible $(z \to x_{\alpha})_{\alpha}$ se factorise, dans
$C_{i+1}$, par $\tau_{i}(x)$ ;
\item[b)] tout diagramme de type $d'$ (resp. $d''$) de $C_{i}$ acquiert dans
$C_{i+1}$ un cône qui appartient à $\delta'_{i+1}$ (resp. $\delta''_{i+1}$) ;
\item[c)] $\delta_{i+1}$ est formé de l'image de $\delta_{i}$ et des cônes
ajoutés en b) ;
\item[d)] pour tout $E \in \mathrm{Ob}\,\Delta$, $\varphi \mapsto \varphi
\circ \tau_{i}$ induit une équivalence
$\underline{\mathrm{Hom}}_{\delta_{i+1}}(C_{i+1}, E) \to
\underline{\mathrm{Hom}}_{\delta_{i}}(C_{i}, E)$.
\end{enumerate}

De d) on tire par récurrence l'équivalence $(d^{*})$ :
$\underline{\mathrm{Hom}}_{\delta_{i}}(C_{i}, E) \simeq
\underline{\mathrm{Hom}}_{\delta}(C, E)$ pour tout $i$. Il reste à trouver
un ordinal $\omega$ tel que $C_{\omega} \in \mathrm{Ob}\,\Delta$, que
$\alpha_{\omega}$ soit $\delta$-exact, et que
$\underline{\mathrm{Hom}}_{\delta_{\omega}}(C_{\omega}, E)$ soit exactement
$\underline{\mathrm{Hom}}_{\Delta,\,\alpha_{\omega}(\delta)}(C_{\omega}, E)$.
On prend un cardinal $\pi$ strictement plus grand que le nombre de flèches
de toute catégorie d'indices en jeu (celles des types $d'$, $d''$ et celles
des cônes de $\delta$), et $\omega$ le plus petit ordinal de cardinal
$\geqslant \pi$ ; alors tout $j < \omega$ est « grand devant » ces
cardinaux, ce qui donne :
\[
(*) \qquad \forall j < \omega,\quad C_{j} \to C_{\omega}
\ \text{est } \delta_{j}\text{-exact},
\]
d'où $C_{\omega} \in \mathrm{Ob}\,\Delta$ (par $(*)$ et b)) et
\[
(**) \qquad \delta_{\omega} = (\text{cônes limites de types $d'$, $d''$
de } C_{\omega}) \cup \alpha_{\omega}(\delta),
\]
qui est le point 3\textsuperscript{o}.\footnote{La page 5 dit de $(*)$ :
« cela résulte aisément de a) et c) et de la grandeur de $\omega$ », avec
en marge « détailler la démonstration ». L'argument attendu : un diagramme
de type $\Theta$ dans $C_{\omega}$ ne fait intervenir que
$\mathrm{card}\,\mathrm{Fl}\,\Theta < \pi$ flèches, donc provient d'un
$C_{j}$, $j < \omega$, et son cône limite y a été ajouté au pas suivant ;
la même remarque de « grandeur » est refaite, complètement cette fois, page
29. Nous ne la redonnons pas ici.}

Tout revient donc à construire $K$. Elle est composée de six opérations
élémentaires, trois pour les cônes projectifs ($\rho = \rho_{3}\rho_{2}
\rho_{1}$) et trois pour les cônes inductifs ($\sigma = \sigma_{3}\sigma_{2}
\sigma_{1}$), chacune « faisant environ la moitié du boulot » :
\begin{itemize}
\item $\rho_{1}$ : prendre l'enveloppe de $C$ dans $\Delta'' =
\Delta_{(\emptyset, d'')}$ — on ajoute librement les limites inductives de
type $d''$ — et ajouter à $\delta''$ tous les cônes de type $d''$ ainsi
créés ; cela assure b'') ;
\item $\rho_{2}$ : égaliser, pour chaque cône projectif $(x \to
x_{\alpha}) \in \delta''_{1}$, les couples de flèches $z \rightrightarrows
x$ que toutes les projections égalisent ; cela assure $a'_{1}$) ;
\item $\rho_{3}$ : ajouter, pour chaque cône projectif de $\delta''_{2}$ et
chaque famille compatible $(z \to x_{\alpha})$, une flèche $z \to x$ la
factorisant ; cela assure $a'_{2}$) ;\footnote{Les indices sont ceux du
manuscrit : la page 5 fait porter $\rho_{2}$ et $\rho_{3}$ sur $\delta''$
alors que $a'_{1}$, $a'_{2}$ concernent $d'$ (la page 4 écrit d'ailleurs
$d'_{i}$). Il y a une hésitation de notation entre « prime » et « seconde »
d'un feuillet à l'autre, sans conséquence : les trois opérations $\rho$
traitent les cônes projectifs et les trois $\sigma$, « définies de façon
analogue », les cônes inductifs.}
\item $\sigma_{1}, \sigma_{2}, \sigma_{3}$ : les mêmes, avec les rôles de
$d'$ et $d''$ échangés.
\end{itemize}
Chaque opération conserve c) et d) — c'est la vérification à faire, notée
en marge — et $K = \sigma\rho$ a les propriétés a) à d).

Le \emph{N.B.} de la page 6 mérite d'être retenu : l'ordre des six
opérations est indifférent, « on pourrait prendre $K' = \sigma_{1}\rho_{3}
\sigma_{2}\rho_{1}\rho_{2}\sigma_{3}$ ! », et la seule chose qui ait un sens
intrinsèque est la limite des itérés transfinis. C'est exactement ce que la
théorie ultérieure des constructions transfinies formalise : le point fixe
ne dépend pas de la présentation choisie de l'opération.

\pagerange{6}{9}
\subsection*{Conditions d'exactitude : la 2-catégorie $\Delta_{f}$}

Les pages 6 à 9 étendent l'énoncé à des conditions plus fines que
l'existence de limites. On se donne, pour $i = 0, 1, 2$, des familles
$f^{(i)} = (f^{(i)}_{\lambda})$ de $\Delta$-foncteurs entre petits objets de
$\Delta$, $f^{(i)}_{\lambda} : R^{(i)}_{\lambda} \to R'^{(i)}_{\lambda}$, et
l'on définit $\Delta_{f}$ comme la sous-2-catégorie pleine de $\Delta$
formée des $E$ tels que, pour tout $\varphi$ de la famille, le foncteur
$u \mapsto u \circ \varphi$, de $\underline{\mathrm{Hom}}_{\Delta}(R', E)$
vers $\underline{\mathrm{Hom}}_{\Delta}(R, E)$, soit
\begin{itemize}
\item[$a_{0}$)] surjectif sur les Hom si $\varphi$ est un $f^{0}_{\lambda}$ ;
\item[$a_{1}$)] injectif sur les Hom (fidèle) si $\varphi$ est un
$f^{1}_{\lambda}$ ;
\item[$a_{2}$)] essentiellement surjectif si $\varphi$ est un
$f^{2}_{\lambda}$.\footnote{La page 6 laisse la condition $a_{2}$) en
blanc ; nous la remplissons par la troisième composante d'une équivalence,
qui est ce que les propriétés $a_{2}$) de la page 7 (existence d'un $u'$
avec $\tau_{i} \circ u \simeq u' \circ \varphi$) construisent. Pour
$a_{1}$) la page écrit « bijectifs », biffé en « injectifs », et nous
suivons la correction. Ainsi $f^{1}$ impose la fidélité, $f^{0}$ et $f^{1}$
ensemble la pleine fidélité, les trois l'équivalence — c'est le
découpage que la Remarque de la page 13 nomme a), b), c).}
\end{itemize}

\textbf{Théorème.} \emph{Sous les hypothèses du théorème précédent, il
existe un $\Delta$-foncteur $\delta$-exact $C \to \overline{C}$ avec
$\overline{C} \in \mathrm{Ob}\,\Delta_{f}$, 2-universel pour ces
propriétés, et $\overline{C}$ est essentiellement petite.}

La démonstration suit le même schéma. Grâce au premier théorème on peut
supposer $C \in \mathrm{Ob}\,\Delta$ et $\delta$ formé de tous les cônes
limites de types $d'$, $d''$, de sorte que
$\underline{\mathrm{Hom}}_{\delta}(C, E) =
\underline{\mathrm{Hom}}_{\Delta,\delta}(C, E)$. On itère alors une
opération $K$ à valeurs dans $\Delta$, dont le pas $\tau_{i} : C_{i} \to
C_{i+1}$ doit :
\begin{enumerate}
\item[$a_{0}$)] pour $\varphi : R \to R'$ dans la famille et deux
transformations $\lambda, \mu : u' \to v'$ entre foncteurs $R' \to C_{i}$,
identifier $\tau_{i} * \lambda$ et $\tau_{i} * \mu$ dès que $\lambda *
\varphi = \mu * \varphi$ ;
\item[$a_{1}$)] pour $\varphi$ parmi les $f^{1}$, $f^{2}$, relever toute
transformation $\lambda' : u'\varphi \to v'\varphi$ en une transformation
$\lambda : \tau_{i} u' \to \tau_{i} v'$ avec $\lambda * \varphi = \tau_{i} *
\lambda'$ ;
\item[$a_{2}$)] pour $\varphi$ parmi les $f^{2}$, prolonger tout $u : R \to
C_{i}$ en un $u' : R' \to C_{i+1}$ avec $\tau_{i} \circ u \simeq u' \circ
\varphi$ ;
\item[b)] être $\delta_{i}$-exact ;
\item[c)] induire, pour $E \in \mathrm{Ob}\,\Delta_{f}$, une équivalence
$\underline{\mathrm{Hom}}_{\Delta,\delta_{i+1}}(C_{i+1}, E) \simeq
\underline{\mathrm{Hom}}_{\Delta,\delta_{i}}(C_{i}, E)$.
\end{enumerate}
Le cardinal $\pi$ majore maintenant aussi les $\mathrm{card}\,\mathrm{Fl}\,
R$, $\mathrm{card}\,\mathrm{Fl}\,R'$ des sources et buts de la famille $f$,
et pour $\omega$ le plus petit ordinal de cardinal $\pi$, $C_{\omega}$ est
dans $\Delta_{f}$ : un foncteur $R \to C_{\omega}$ se factorise par un
$C_{j}$, $j < \omega$, et les propriétés $a_{0}$)--$a_{2}$) au pas suivant
donnent la surjectivité, la fidélité, la surjectivité essentielle
voulues.\footnote{Ce second théorème donne, à titre d'exemple, les enveloppes
dans les 2-catégories de catégories \emph{avec conditions d'exactitude} —
celles où l'on demande que certaines limites en commutent à d'autres. Le
cas emblématique est celui des topos : les conditions de Giraud sont des
conditions de ce type, et c'est ainsi que la page 17 lira « $\Delta$ = bons
topos ». La page 9 s'arrête sur « il reste donc à faire la construction de
$K(C, \delta)$ » : cette construction n'est pas écrite, ni ici ni dans la
reprise des pages 26 à 31, qui ne fait que l'énumérer (propriétés
$e_{0}$), $e_{1}$), $e_{2}$)).}

\pagerange{10}{14}
\subsection*{Une typologie des conditions d'exactitude}

Les pages 10 à 14 recommencent, sous le même titre, pour classer les
conditions que l'on peut vouloir imposer et voir lesquelles suffisent. Sur
une catégorie $C$, on distingue :
\begin{enumerate}
\item[a)] l'existence inconditionnelle de certaines limites (les types
$d = (d', d'')$) ;
\item[b)] l'existence \emph{conditionnelle} : on se donne un diagramme
$\Theta$ muni de sous-diagrammes $\Theta'_{i} \to \Theta$, $\Theta''_{j} \to
\Theta$, et l'on demande que pour tout $\Theta \to C$ les limites des
restrictions aux $\Theta'_{i}$ et $\Theta''_{j}$ existent (par exemple :
les produits fibrés existent dès que la somme amalgamée existe) ;
\item[c)] l'exactitude de cônes : un paquet $(\delta, \lambda)$ étant donné,
tout foncteur $\delta$-exact doit être $\lambda$-exact ;
\item[d)] la forme la plus générale : un foncteur $\varphi : \Theta \to
\overline{\Theta}$ entre catégories munies de paquets $\delta$,
$\overline{\delta}$, et le foncteur de restriction
\[
\varphi^{*}_{C} : \underline{\mathrm{Hom}}(\overline{\Theta}, C)
\longrightarrow \underline{\mathrm{Hom}}(\Theta, C),
\]
dont on peut demander qu'il envoie $\underline{\mathrm{Hom}}_{\overline
{\delta}}$ dans $\underline{\mathrm{Hom}}_{\delta}$ (i), qu'il y soit fidèle
(i'), surjectif sur les Hom (ii), ou essentiellement surjectif (ii').
\end{enumerate}
La page 11 observe que (i) est une condition d'\emph{existence de flèches}
et (ii') une condition d'\emph{existence d'objets et de diagrammes}, et
retient finalement deux types de conditions qui permettent d'exprimer les
autres :
\[
1^{\circ})\ \ \mathrm{Im}\bigl(\varphi^{*}_{C} \mid
\underline{\mathrm{Hom}}_{\overline{\delta}}\bigr) \subset
\underline{\mathrm{Hom}}_{\delta}, \qquad
2^{\circ})\ \ \mathrm{Im}\bigl(\varphi^{*}_{C} \mid
\underline{\mathrm{Hom}}_{\overline{\delta}}\bigr) \supset
\underline{\mathrm{Hom}}_{\delta},
\]
la première impliquant les conditions d'exactitude habituelles, la seconde
les conditions d'existence (y compris conditionnelle) de limites.

Pour les \emph{foncteurs} $F : C \to C'$, on se donne des triplets
$(\Theta, \delta, \lambda)$ et l'on dit que $F$ est $\rho$-exact si pour tout
$u : \Theta \to C$ qui est $\lambda$-exact, $F \circ u$ l'est. La page 12
vérifie ce qui fait de cette classe $\Sigma_{\rho} \subset
\mathrm{Fl}(\mathrm{Cat})$ une bonne classe de flèches : stable par
isomorphisme, contenant les équivalences, stable par composition, et stable
par intersection. Symétriquement, un couple $\sigma = (\sigma_{1},
\sigma_{2})$ de familles de quadruplets $(R, \delta, \overline{R},
\overline{\delta}, \varphi)$ définit une classe d'objets $\Delta_{\sigma}
\subset \mathrm{Ob}(\mathrm{Cat})$ par les conditions $1^{\circ}$) pour
$\sigma_{1}$ et $2^{\circ}$) pour $\sigma_{2}$ ; elle est stable par
équivalence et par intersection. La 2-catégorie $\Delta_{\sigma,\rho}$ a
pour objets $\Delta_{\sigma}$ et pour flèches $\Sigma_{\rho}$.\footnote{Le
manuscrit demande, page 12, que les foncteurs identiques soient dans
$\Sigma_{\rho}$, ce qui impose de se restreindre aux $C$ où tout $\Theta
\to C$ $\delta$-exact est $\lambda$-exact — c'est pourquoi les objets sont
contraints par $\sigma$ et les flèches par $\rho$ séparément. La
\emph{Remarque} de la page 13 demande si la sous-2-catégorie $\Delta_{0}
\subset \Delta$ définie par la fidélité, la pleine fidélité ou
l'équivalence de $\varphi^{*}$ est « distinguée » ; la question est laissée
ouverte, et la page 14 ne fait que redresser les données (a), (b), (c) avec
lesquelles la reprise de la page 26 travaillera.}

C'est cette typologie qui est la vraie nouveauté des feuillets : la
définition d'une classe de catégories par des propriétés de
\emph{restriction le long de foncteurs entre esquisses} recouvre à la fois
les conditions d'existence et les conditions d'exactitude, et c'est sous
cette forme que le théorème sera réénoncé.

\pagerange{26}{31}
\subsection*{La reprise générale : $\Delta_{\sigma,\rho}$ (pages 26 à 31)}

La page 26 reprend le théorème depuis le début sous la forme la plus
générale du dossier. On se donne
\begin{enumerate}
\item[1)] un ensemble $\sigma$ de couples $(S, \lambda)$, $S$ une petite
catégorie et $\lambda = (\lambda', \lambda'')$ un petit paquet de cônes de
$S$ ;
\item[2)] un ensemble $\rho$ de systèmes $(R_{1}, \delta_{1}, R_{2},
\delta_{2}, \varphi)$, avec $\varphi : R_{2} \to R_{1}$ un foncteur entre
petites catégories munies de paquets, tel que $\varphi(\delta_{2}) \subset
\delta_{1}$ ;
\end{enumerate}
et $\Delta = \Delta_{\sigma,\rho}$ est la sous-2-catégorie de $\mathrm{Cat}$
dont les objets $E$ vérifient $a_{1}$) pour tout $(S, \lambda) \in \sigma$
et tout $u : S \to E$, les limites des cônes de $u(\lambda)$ existent, et
$a_{2}$) pour tout système de $\rho$, $u \mapsto u \circ \varphi$ induit une
équivalence $\underline{\mathrm{Hom}}_{\delta_{1}}(R_{1}, E) \simeq
\underline{\mathrm{Hom}}_{\delta_{2}}(R_{2}, E)$ ; et dont les flèches $E \to
E'$ sont les foncteurs qui envoient tout $u : S \to E$ $\lambda$-exact sur
un foncteur $\lambda$-exact. On suppose (A) que tout est petit, et (B) que
les $(R_{i}, \delta_{i})$ des systèmes de $\rho$ figurent dans
$\sigma$.\footnote{Les esquisses « mixtes » — existence de limites par
$\sigma$, exactitude et existence conditionnelle par $\rho$ — sont
exactement les \emph{esquisses} de la littérature ultérieure, et la
2-catégorie $\Delta_{\sigma,\rho}$ est celle des catégories qui les
« réalisent ». La condition (B) dit que les catégories d'où partent les
conditions d'exactitude sont elles-mêmes soumises aux conditions
d'existence ; sans elle les propriétés $e_{0}$)--$e_{2}$) ne se
propageraient pas.}

\textbf{Théorème.} \emph{Pour $C$ essentiellement petite munie d'un petit
paquet $\delta$, il existe $C' \in \mathrm{Ob}\,\Delta$ et $i : C \to C'$
$\delta$-exact, 2-universels : pour tout $E \in \mathrm{Ob}\,\Delta$,}
\[
\underline{\mathrm{Hom}}_{\Delta,\, i(\delta)}(C', E)
\;\xrightarrow{\ \approx\ }\;
\underline{\mathrm{Hom}}_{\mathrm{Cat},\,\delta}(C, E),
\]
\emph{et $C'$ est essentiellement petite.}

La démonstration (pages 27 à 30) est celle des pages 3 à 5, mise au net.
La construction $K$, indépendante de $i$, doit vérifier :
\begin{enumerate}
\item[a)] l'équivalence $\underline{\mathrm{Hom}}_{\delta_{i+1}}(C_{i+1}, E)
\to \underline{\mathrm{Hom}}_{\delta_{i}}(C_{i}, E)$ ;
\item[b)] $\tau_{i}(\delta_{i}) \subset \delta_{i+1} \subset
\tau_{i}(\delta_{i}) \cup \bigcup_{(S,\lambda) \in \sigma} \bigcup_{u : S
\to C_{i+1}} u(\lambda)$ ;
\item[c)] $\tau_{i}$ est préexact pour $\delta_{i}$ ;
\item[d)] pour $(S, \lambda) \in \sigma$ et $u : S \to C_{i}$, il existe un
foncteur $\overline{v} : \overline{S} \to C_{i+1}$, avec $\overline{v}
\beta = \tau_{i} u$ et $\overline{v}(\overline{\lambda}) \subset
\delta_{i+1}$, où $\beta : (S, \lambda) \to (\overline{S},
\overline{\lambda})$ est obtenu en \emph{dédoublant les sommets des cônes
critiques} (voir ci-dessous) ;
\item[$e_{0}$)--$e_{2}$)] les trois propriétés de la page 7, relatives aux
systèmes de $\rho$ : identification des 2-flèches, relèvement des
2-flèches, prolongement des foncteurs.
\end{enumerate}
De a) et b) suivent par récurrence (1) $\underline{\mathrm{Hom}}_{\delta_{i}}
(C_{i}, E) \simeq \underline{\mathrm{Hom}}_{\delta}(C, E)$ et (2)
$\alpha_{i}(\delta) \subset \delta_{i} \subset \alpha_{i}(\delta) \cup
\bigcup \tau_{ij} u(\lambda)$. Pour $\omega$ un ordinal limite grand devant
le cardinal $\Pi$ qui majore toutes les catégories d'indices, on prouve
successivement (3) que $\tau_{\omega,i} : C_{i} \to C_{\omega}$ est
$\delta_{i}$-exact — c'est c) et la grandeur de $\omega$, en toutes lettres
cette fois —, puis $A_{1}$) et $A_{2}$) (par d) et par e) respectivement),
et enfin 4) : tout $u : S \to C_{\omega}$ $\lambda$-exact se factorise à
isomorphisme près par un $u_{i} : S \to C_{i}$ avec $u_{i}(\lambda) \subset
\delta_{i}$. Le point 4) est le seul délicat, et la page 30 le démontre par
le dédoublement.

\textbf{Le dédoublement des sommets} (page 31). À $(S, \lambda)$ on associe
$(\overline{S}, \overline{\lambda})$ en ajoutant, pour chaque cône $c : (x
\to x_{\alpha})_{\alpha}$ de $\lambda'$, un nouvel objet $\overline{x}_{c}$,
une flèche $i_{c} : x \to \overline{x}_{c}$ et des flèches $q_{\alpha} :
\overline{x}_{c} \to x_{\alpha}$ avec $q_{\alpha} i_{c} = p_{\alpha}$ ; et
dualement pour $\lambda''$ :

\begin{tikzcd}[column sep=small, row sep=small]
& x_{1} & \\
x \arrow[r, "i_{c}"] \arrow[ur, bend left=20, "p_{1}"] \arrow[dr, bend right=20, "p_{3}"'] & \overline{x}_{c} \arrow[u, "q_{1}"] \arrow[r, "q_{2}"] \arrow[d, "q_{3}"'] & x_{2} \\
& x_{3} &
\end{tikzcd}

Le paquet $\overline{\lambda}$ est formé des cônes de sommet $\overline
{x}_{c}$. Alors, pour tout $E \in \mathrm{Ob}\,\Delta$, tout foncteur $u :
S \to E$ se prolonge de façon unique en un $\overline{u} : \overline{S} \to
E$ qui soit $\overline{\lambda}$-exact (on envoie $\overline{x}_{c}$ sur la
limite du diagramme $u(x_{\alpha})$), et $u$ est $\lambda$-exact si et
seulement si $\overline{u}(i_{c})$ est un isomorphisme pour tout $c$. C'est
la réduction, aujourd'hui classique, de l'exactitude d'un cône à
l'\emph{inversibilité d'une flèche de comparaison} ; elle ramène le point 4)
à la propriété d) et à un argument de cardinalité sur les $\bar{u}_{j}
(i_{c})$.\footnote{La page 30 se termine par « cqfd », mais deux renvois
notés (*) et (**) restaient à écrire ; la page 31 écrit (**), le
dédoublement, et commence (*), la définition du foncteur $\delta$-préexact,
qui s'arrête après trois lignes. La reprise est donc complète pour tout ce
qui concerne l'ordinal $\omega$ et laisse à l'état d'énoncé, comme la
première version, l'opération $K$ elle-même. L'existence de $K$ n'est pas
en doute — les six opérations des pages 5--6 et leurs analogues pour
$e_{0}$)--$e_{2}$) la fournissent — mais elle n'est écrite nulle part dans
le dossier.}

\pagerange{17}{24}
\section*{Théories algébriques et topos classifiants (pages 17 à 24)}

\pagerange{17}{17}
\subsection*{Le topos enveloppe d'une catégorie à limites finies}

Le feuillet 17 applique ce qui précède au cas qui intéresse le dossier
entier. On prend $\Delta_{0}$ = les catégories à limites projectives
finies, avec les foncteurs exacts à gauche ; et $\Delta$ = les « bons
topos », c'est-à-dire les catégories satisfaisant aux conditions
d'exactitude de Giraud (limites projectives finies, sommes disjointes
universelles, relations d'équivalence effectives universelles), avec pour
flèches les foncteurs exacts à gauche commutant aux limites inductives —
les foncteurs image inverse.\footnote{Le manuscrit écrit « $\Delta$ = bons
topos (conditions $\supset$ b) c) de Giraud), foncteurs commutant aux
$\varprojlim$ et exacts à g. ». Le mot « bon » signale que la condition de
\emph{taille} de Giraud — l'existence d'une petite famille génératrice —
n'est pas demandée : $\Delta$ contient des catégories qui ne sont pas des
topos de Grothendieck, et c'est précisément ce qui donne son contenu au
Théorème 4) ci-dessous. Nous appellerons « topos » un topos de Grothendieck
et « objet de $\Delta$ » un bon topos.}

Une \emph{théorie algébrique} $T_{0}$ relative à $\Delta_{0}$ est
représentée par une petite catégorie $C \in \mathrm{Ob}\,\Delta_{0}$ :
$T_{0}(E) = \underline{\mathrm{Hom}}_{\Delta_{0}}(C, E)$ pour $E \in
\Delta_{0}$. Pour $E \in \mathrm{Ob}\,\Delta$ on a alors
\[
T_{0}(E) = \underline{\mathrm{Hom}}_{\Delta_{0}}(C, E)
\;\xrightarrow{\ \approx\ }\;
\underline{\mathrm{Hom}}_{\Delta}(\widehat{C}, E)
\;\simeq\; \underline{\mathrm{Hom}}_{\mathrm{top}}(E, \widehat{C})^{\circ}
\quad (\text{si $E$ est un topos}),
\]
où $\widehat{C}$ est la catégorie des préfaisceaux sur $C$. C'est le
théorème de Diaconescu sous sa forme la plus simple : un foncteur exact à
gauche $C \to E$ se prolonge de façon unique en un foncteur image inverse
$\widehat{C} \to E$, et les morphismes de topos $E \to \widehat{C}$
correspondent aux foncteurs plats $C \to E$, qui sont ici les foncteurs
exacts à gauche puisque $C$ a les limites finies.\footnote{Diaconescu,
\emph{Change of base for toposes with generators} (1975), postérieur à la
fourchette de l'inventaire ; le cas des préfaisceaux sur une catégorie à
limites finies est dans SGA 4 IV 4.9.4. L'énoncé « $\widehat{C}$ est la
$\Delta_{0} \to \Delta$-enveloppe de $C$ », qui fait de Diaconescu un cas
particulier du théorème d'enveloppe, est la façon propre au dossier de le
voir. Le passage à l'opposé dans $\underline{\mathrm{Hom}}_{\mathrm{top}}
(E, \widehat{C})^{\circ}$ est une convention : un morphisme de topos $f \to
g$ est une transformation $f^{*} \to g^{*}$, ou $g^{*} \to f^{*}$ selon
l'auteur, et le manuscrit prend la seconde.} Ainsi $\widehat{C}$ est le
\emph{topos classifiant} de la théorie $T_{0}$, et le passage $\Delta_{0}
\to \Delta$ est le passage d'une théorie « à limites finies » à sa
réalisation dans les topos.

Les deux problèmes suivent. \emph{Problème 1} : la $\Delta_{0}$-enveloppe
d'une petite catégorie $I$ est la sous-catégorie pleine $\rho_{\Delta_{0}}(I)
\subset \underline{\mathrm{Hom}}(I, \mathrm{Ens})^{\circ}$ engendrée par $I$
sous les limites finies — les foncteurs « de présentation finie » —, et sa
$\Delta$-enveloppe est $\widehat{\rho_{\Delta_{0}}(I)}$. \emph{Problème 2} :
une sous-théorie d'une $\Delta$-théorie représentable par un topos $R$
correspond à un sous-topos. Précisément, $\Sigma$ étant un ensemble de
flèches de $R$, on considère $E \mapsto \underline{\mathrm{Hom}}_{\Delta,
\Sigma}(R, E)$ (les images inverses rendant inversibles les flèches de
$\Sigma$) ; on sature $\Sigma$ par les trois conditions ST1 (deux des trois
flèches d'un triangle commutatif), ST2 (stabilité par changement de base),
ST3 (nature locale) de SGA 4 I 9.1.1, d'où $\Sigma'$, et $R' = R\Sigma'^{-1}$
est un sous-topos de $R$ avec
\[
\underline{\mathrm{Hom}}_{\Delta,\Sigma}(R, E) \;\xrightarrow{\ \approx\ }\;
\underline{\mathrm{Hom}}_{\Delta}(R', E) ;
\]
tout sous-topos s'obtient ainsi, et si $R \simeq \widehat{C}$ les
sous-topos correspondent aux topologies $T'$ sur $C$.\footnote{Ce sont les
« sous-topos $R'$ comme catégories de fractions » de SGA 4 IV 9 ; le
renvoi à I 9.1.1 est celui du manuscrit. La page 18 ajoute qu'un sous-topos
donné se décrit toujours par un \emph{petit} ensemble $\Sigma$ : on prend
$C \subset R$ petite génératrice et les familles couvrantes de $C$ pour
$T'$.} Un sous-topos donné se décrit par un petit $\Sigma$ ; et la page 18
« corrige » : si $T_{R}$ est algébrique, la sous-théorie $T_{R'}$ l'est
aussi, déduite par un petit ensemble d'axiomes. Un plongement de théories
$T_{R'} \to T_{R}$ (i.e. $T_{R'}(E) \to T_{R}(E)$ pleinement fidèle pour tout
$E$) correspond à un morphisme $R \to R'$ de $\Delta$ tel que
$\underline{\mathrm{Hom}}_{\Delta}(R', E) \to \underline{\mathrm{Hom}}_{
\Delta}(R, E)$ soit pleinement fidèle pour tout $E$ : c'est la
caractérisation des plongements de topos $R' \to R$.

\pagerange{18}{22}
\subsection*{Le théorème : théories algébriques et topos}

\textbf{Théorème 4.} \emph{a) Un objet $R$ de $\Delta$ représente une théorie
algébrique si et seulement si $R$ est un topos, c'est-à-dire admet une
petite sous-catégorie génératrice. b) Soient $I$ un ensemble, $R \in
\mathrm{Ob}\,\Delta$ et $b : I \to R$ une famille d'objets. Alors $b$ est
basique pour $T_{R}$ si et seulement si la sous-catégorie pleine de $R$
engendrée par les produits finis d'objets de $b(I)$ est strictement
génératrice.}\footnote{« Algébrique » n'est pas défini sur ces feuillets ;
le sens, constant dans le dossier (pages 62--63, 71, 84--87), est
\emph{constructible} : une théorie est algébrique si elle s'obtient à
partir de la structure vide sur des objets de base par une suite,
éventuellement transfinie, de trois opérations — ajouter des flèches,
imposer des égalités, imposer des conditions d'exactitude —, ce qui est
exactement la présentation d'une \emph{théorie géométrique} par des
axiomes. « Basique » signifie que la donnée d'une $T$-structure sur $E$
détermine fidèlement la famille de ses objets de base, i.e. que $T_{R}(E)
\to E^{I}$ est fidèle (page 86, Remarque 1). Le théorème 4 a) est ainsi
l'énoncé « tout topos de Grothendieck est le topos classifiant d'une
théorie géométrique, et réciproquement », qui sera publié par Joyal et
Wraith puis Makkai et Reyes (1977) ; le manuscrit le sait par construction,
puisque « topos » y signifie « objet de $\Delta$ à petit générateur ». La
condition b) est dans le même esprit : un générateur strict donne les
sortes et les opérations.}

La démonstration (pages 18 à 22), dans le sens « suffisance », construit la
théorie représentée par $(R, b)$ en trois pas, qui sont ceux de toute
présentation :
\begin{enumerate}
\item[$1^{\circ}$)] la structure vide sur des objets de base $(X_{i})_{i \in
I}$, représentée par le topos $R_{0} = \widehat{J_{0}}$, $J_{0} =
\rho_{\Delta_{0}}(I)$ (problème 1 ; $R_{0}$ est le topos des ensembles
$I$-multisimpliciaux, le classifiant de $I$ objets) ;\footnote{Pour $I$
réduit à un point, $\widehat{J_{0}}$ est le topos des foncteurs sur les
ensembles finis, le \emph{classifiant d'un objet} (\emph{object
classifier}). « Ensembles $I$-multisimpliciaux » est le nom du manuscrit,
avec un $\omega$ ajouté au-dessus : les objets de $J_{0}$ sont les produits
finis d'objets de base, et un préfaisceau sur $J_{0}$ est une famille
d'ensembles multi-indexée.}
\item[$2^{\circ}$)] on ajoute des flèches : $J_{1}$ est la sous-catégorie
pleine de $R$ engendrée par $b(I)$ et ses produits finis, et $R_{1} =
\rho_{\Delta}(J_{1})$ ;
\item[$3^{\circ}$)] on met les axiomes d'égalité entre composés, en vertu de
$R$, d'où $R_{2}$ et $T_{R_{2}}(E) \simeq \underline{\mathrm{Hom}}(J_{1},
E)$, et $T_{R} \to T_{R_{2}}$ est pleinement fidèle ;
\item[$4^{\circ}$)] on ajoute les axiomes explicites qui expriment la
continuité des foncteurs $J_{1} \to E$ : transformer les cônes inductifs
exacts de $R$ et les familles couvrantes en cônes exacts et familles
couvrantes.
\end{enumerate}
On obtient $I \to R_{0} \to R_{1} \to R_{2} \to R_{3} \simeq R$, les
morphismes $R_{1} \to R_{2} \to R_{3}$ étant des localisations (le
schéma de la page 22 les note « loc »), et $\widehat{J_{1}} \to R$ est un
foncteur de localisation. Le \emph{N.B.} de la page 19 condense $3^{\circ}$
et $4^{\circ}$ en un seul pas : trois pas minimum, objets, flèches (aux
produits finis), axiomes sur les flèches. Le \emph{Re N.B.} de la page 20
donne le raccourci si l'on veut seulement qu'un topos soit « algébrique » :
choisir $C$ petite à limites finies avec $R \simeq \widehat{C}$, prendre la
théorie $T^{\Delta_{0}}_{C}$ représentée par $\widehat{C}$, et ajouter les
axiomes de continuité. C'est le mode d'emploi, aujourd'hui standard, pour
présenter un topos $\widetilde{C}$ comme classifiant : la théorie « à
limites finies » de $C$, plus les axiomes de recouvrement.

La \emph{nécessité} (page 20) demande trois lemmes de construction, tous
ramenés au théorème d'enveloppe : a) la structure vide sur $I$ est
représentée par $\widehat{\widetilde{I}}$, $\widetilde{I} = \rho_{\Delta_{0}}
(I)$, et $\widetilde{I}$ $\Delta$-engendre $\widehat{\widetilde{I}}$ « en
deux pas » (produits, puis limites inductives quelconques) ; b) mettre des
axiomes sur $T_{R}$ est le problème 2 ; c) ajouter un petit ensemble $H$ de
flèches revient à un \emph{produit 2-fibré de théories}
\[
T' \simeq T_{R} \times^{(2)}_{T_{I}} T_{C},
\]
où $C$ est la catégorie libre engendrée par les objets de base et les
flèches de $H$, $T_{C}(E) = \underline{\mathrm{Hom}}_{\mathrm{Cat}}(C, E)$ et
$T_{I}$ de même. Comme $T_{C}$ et $T_{I}$ sont représentables (problème
1), il faut savoir que \emph{les produits 2-fibrés de topos existent} —
ce que la page 21 démontre — et d) que les \emph{2-limites inductives
transfinies de topos} existent (page 22).

\pagerange{21}{22}
\subsection*{Produits 2-fibrés de topos et passage à la limite}

Pour $f^{*} : T \to R$ et $g^{*} : T \to S$ deux images inverses, on
choisit de petites sous-catégories génératrices $T_{0}$, $R_{0}$, $S_{0}$
avec $f^{*}(T_{0}) \subset R_{0}$, $g^{*}(T_{0}) \subset S_{0}$, et $U_{0}$ la
2-somme amalgamée de $R_{0}$ et $S_{0}$ sous $T_{0}$ dans $\mathrm{Cat}$, de
sorte que $\underline{\mathrm{Hom}}(U_{0}, E) \simeq \underline{\mathrm{Hom}}
(R_{0}, E) \times^{2}_{\underline{\mathrm{Hom}}(T_{0}, E)}
\underline{\mathrm{Hom}}(S_{0}, E)$. Si $R_{0}$, $S_{0}$ sont stables par
limites finies, la condition qu'un $h : U_{0} \to E$ provienne de
$\underline{\mathrm{Hom}}_{\Delta}(R, E) \times^{2}_{\underline{\mathrm{Hom}}
_{\Delta}(T,E)} \underline{\mathrm{Hom}}_{\Delta}(S, E)$ s'exprime par des
axiomes d'exactitude et de continuité sur $h$ ; la théorie $T_{R}
\times^{(2)}_{T_{T}} T_{S}$ est donc une sous-théorie pleine de
$T_{\rho_{\Delta}(U_{0})}$, et par le problème 2 elle est représentée par
un \emph{sous-topos} de $\rho_{\Delta}(U_{0})$ : le produit 2-fibré
$R \times^{2}_{T} S$ existe dans $(\mathrm{Top})$, et c'est une 2-limite
amalgamée $R \amalg_{T} S$ dans $\Delta$.\footnote{C'est l'existence des
produits fibrés de topos, SGA 4 IV 8.2 ; le manuscrit dit « cette
construction est $\pm$ connue, redonnons-en une autre ». La sienne, par la
somme amalgamée des sous-catégories génératrices puis passage au
sous-topos, est celle que l'on donne aujourd'hui en logique catégorique :
le produit fibré de topos classifiants classifie la « somme amalgamée » des
théories. Il note en passant que si $T_{0} \to S_{0}$ est surjectif sur les
objets alors $R \to R \amalg_{T} S$ a une image génératrice.}

Le passage à la limite transfinie (page 22) est le point où la méthode
tire sur sa taille : une 2-limite inductive transfinie de topos dans
$\Delta$ est un topos, pourvu que l'on choisisse des petites catégories
génératrices stables par les morphismes de transition — « par exemple à
l'aide des filtrations cardinales, mais c'est un mauvais moyen » — et que
l'on prenne $C = \varinjlim_{\alpha} C_{\alpha}$ dans $\mathrm{Cat}$, la
limite étant un sous-topos de $\widehat{C}$ ; l'énoncé utile est que si
$J \to R_{\alpha}$ engendre $R_{\alpha}$ pour tout $\alpha$, alors $J \to R$
engendre $R$. La page reporte le schéma des trois pas, avec en marge l'idée
qui gouverne les feuillets 81 à 87 : « $R$ apparaît comme une catégorie de
fractions de $\widehat{J_{1}}$, i.e. $J_{1}$ est l'\emph{analyseur} libre
sur $I$ engendré par les flèches induites par $R$ ».\footnote{Le mot
« analyseur » est celui de Lazard (\emph{Lois de groupes et analyseurs},
1955) ; il revient dans le titre de la page 89, « Relation avec analyseurs
de Lazard », qui reste sans texte. Un analyseur est, en substance, une
théorie algébrique présentée par ses opérations $n$-aires ; la catégorie
$J_{1}$ des produits finis d'objets de base avec les flèches de $R$ est
exactement la \emph{théorie de Lawvere} de la structure, et c'est le
rapprochement que le dossier fait sans le nom de Lawvere, dont la thèse
date de 1963.}

\pagerange{23}{24}
\subsection*{Deux feuillets d'exemples : le topos d'un objet}

La page 23 note d'abord ce que ces théories ont en commun : $T(E) \simeq
\underline{\mathrm{Hom}}_{\mathrm{top}}(E, R)^{\circ} \simeq
\underline{\mathrm{Hom}}_{\Delta}(R, E)$ est stable par limites inductives
filtrantes — les conditions sur les structures « convergent » — et
$T_{R}(\mathrm{Ens}) \simeq \underline{\mathrm{Pt}}(R)$ est la catégorie des
points du topos $R$.\footnote{Le manuscrit écrit $\underline{\mathrm{Pt}}
(R)^{\circ} \simeq \underline{\mathrm{Fib}}(R)$ ; l'opposé est la même
convention de sens des 2-flèches que ci-dessus, et $\underline{\mathrm{Fib}}$,
les foncteurs fibres, est un autre nom des points.} La \emph{Question} qui
suit est la bonne : si $T(E)$ a des limites projectives finies pour tout
$E$, et si les $T(E) \to T(E')$ y commutent, $T$ est-elle la restriction à
$\Delta$ d'une théorie algébrique sur $\Delta_{0}$, i.e. $R \simeq
\widehat{C}$ avec $C$ stable par limites finies ? C'est la question de
savoir quand un topos classifiant est un topos de préfaisceaux, i.e. quand
une théorie géométrique est équivalente à une théorie \emph{cartésienne}
(à limites finies) ; les pages 77 et 78 y répondront pour les théories
définies par des limites, et la réponse générale est la même : lorsque les
modèles ont assez d'objets de présentation finie.

Le second feuillet prend $R = \widehat{C}$ avec $C$ la catégorie opposée
aux ensembles finis, de sorte que $R = \underline{\mathrm{Hom}}(\mathrm{Ens}_{f},
\mathrm{Ens})$ est le \emph{topos classifiant d'un objet} : ses points sont
les ensembles, et un morphisme $E \to R$ est un objet de
$E$.\footnote{Le manuscrit écrit « $C$ = catégorie opposée aux ens. finis,
donc $R$ = ens. semi-simpliciaux », avec une note marginale « pts à les
$\Delta_{n}$ » ; les deux descriptions ne coïncident pas (les ensembles
simpliciaux sont les préfaisceaux sur les ordinaux finis, dont les points
ne sont pas les ensembles), et tout ce qui suit sur la page — les points
« de p. f. » $\mathbb{N}$, les sous-topos décrits par la structure $X$ —
est vrai du classifiant d'un objet et faux des ensembles simpliciaux. Nous
lisons donc $C = (\mathrm{Ens}_{f})^{\circ}$, comme la page l'écrit, et
nous rattachons « $\Delta_{0}$, $\Delta_{1}$ » aux ensembles à $0$ et $1$
éléments.} On cherche ses sous-topos, « classés par leurs points » : le
feuillet en trouve six, rangés en un treillis,

\begin{tikzcd}[column sep=small, row sep=small, nodes={font=\scriptsize}]
& R & & \\
R^{+} \arrow[ur, no head] & & V \arrow[ul, no head] & \\
& R^{+} \wedge V = \varphi_{1} \arrow[ul, no head] \arrow[ur, no head] & & \overline{(\varphi_{0})} \arrow[ul, no head] \\
& & \emptyset \arrow[ul, no head] \arrow[ur, no head] &
\end{tikzcd}

où $R^{+}$ classifie les objets $X$ tels que $X \to e$ soit épimorphique
(les objets « habités »), $V$ ceux tels que $X \to e$ soit un
monomorphisme (les sous-objets de l'objet final : $V$ représente $E
\mapsto \mathcal{P}(e_{E})$), $\varphi_{1} = R^{+} \wedge V$ le point
« $X \simeq e$ », $\overline{(\varphi_{0})}$ le point « $X = \emptyset$ »,
et $\emptyset$ le topos vide. $V$ est un sous-topos fermé qui n'est pas
ouvert ; $V$ et $\varphi_{1}$ ne sont pas des complémentaires mais ont
pour « complémentaires faibles » $\emptyset$ et $\varphi_{0}$.\footnote{Ces six
sous-topos sont corrects, mais ils ne sont pas tous les sous-topos de $R$,
et le « en résumé : 6 sous-topos » de la page 24 ne doit pas être lu comme
une classification. Les sous-topos du classifiant d'un objet correspondent
aux théories géométriques quotients de la théorie d'un objet, et la
condition « $X$ a au plus $n$ éléments », qui s'écrit $\top \vdash
\bigvee_{i<j} x_{i} = x_{j}$ pour $n+1$ variables, est géométrique : elle
définit pour chaque $n$ un sous-topos dont les points sont les ensembles de
cardinal $\leqslant n$, et qui n'est pas dans la liste. Le feuillet
restreint sa recherche, comme le dit son titre, aux sous-topos
\emph{ouverts} et à leurs fermés complémentaires, et aux intersections de
ceux-ci ; pour ceux-là la liste est plausible. La mention « $V$ n'est ni
ouvert ni \ldots » est la sienne, illisible ensuite.}

Le feuillet se ferme sur une liste de questions — les topos classifiants
et sous-topos pour les structures « $f : X \to Y$ sans plus », « $f : X \to
X$ » (c'est $B_{\mathbb{N}}$, le classifiant du monoïde libre à un
générateur), « $e \to X$ », « $X \times Y \to Z$ », « $X \times X \to X$
(!) », « $X \times X \rightrightarrows X$ (!!) » — qui sont, chacune, le
classifiant d'une signature sans axiomes. La marque d'insistance sur les
deux dernières est justifiée : le classifiant d'une loi de composition
binaire sans axiome, et celui d'un couple de lois, sont les premiers cas où
la question des sous-topos rejoint l'algèbre universelle proprement dite.

\pagerange{33}{34}
\section*{Deux feuilles de travail : théories de Lawvere, monades, descente (pages 33 et 34)}

Deux feuillets sans ordre linéaire, en anglais et en français, notent ce
que la littérature du moment offre au projet. Le premier pose une
\emph{théorie de Lawvere} $\mathcal{L}$ — une petite catégorie dont les
objets sont les puissances $n = \mathbb{1}^{n}$ d'un objet $\mathbb{1}$ —
et l'adjonction entre les modèles $\underline{\mathrm{Hom}}'(\mathcal{L}^{\circ},
\mathrm{ENS})$ (foncteurs conservant les produits) et les algèbres
$\mathrm{Alg}(\mathfrak{T})$ d'une monade ; il note les propriétés de
l'objet $P$ qui rend $\mathrm{Hom}(P, -)$ un bon foncteur fibre (puissances
$P_{I} = \coprod_{I} P$, relations d'équivalence effectives, $f : X \to Y$
effectif si et seulement si $\mathrm{Hom}(P, X) \to \mathrm{Hom}(P, Y)$
est surjectif, produits fibrés), et enfin le \emph{théorème de
monadicité} : un foncteur $U : A \to B$ admettant un adjoint à gauche $F$
est monadique ($A \simeq \mathrm{Alg}(\mathbb{T})$, $\mathbb{T} = UF$) si
$U$ est conservatif et si $A$ a et $U$ préserve les coégalisateurs des
paires réflexives — le critère de Beck, dont les noms de Beck, Barr et
Linton sont notés au bas de la feuille.\footnote{Le mot du manuscrit est
\emph{tripleable}, terme de l'époque (une « triple » est une monade), et
les trois noms sont peu nets. Le critère de Beck date de 1967 ; Linton,
\emph{Some aspects of equational categories} (1966), Barr, \emph{Coequalizers
and free triples} (1970). Le second point du critère, laissé illisible,
est nécessairement celui-là. C'est le seul endroit du dossier où le nom de
Lawvere est écrit, et il désigne la notion que les pages 18--22 et 84--87
construisent sous le nom d'analyseur.}

Le second feuillet est le \emph{critère de descente} pour un morphisme de
topos $f : B \to B'$, d'image inverse $f^{*} : B' \to B$. On forme la
comonade $\pi = f^{*}f_{*}$ sur $B$, exacte à gauche et accessible, la
catégorie $B_{\pi}$ de ses coalgèbres, et le foncteur de comparaison $B'
\to B_{\pi}$, $X \mapsto f^{*}X$. Le théorème est que ce foncteur est une
équivalence dès que $f^{*}$ est conservatif : c'est le théorème de Beck
appliqué à l'adjonction $f^{*} \dashv f_{*}$, la seconde condition du
critère — que $f^{*}$ préserve les noyaux de couples coréflexifs — étant
acquise parce que $f^{*}$ est exact à gauche. Un morphisme de topos dont
l'image inverse est conservative est donc un morphisme de descente
effective.\footnote{Le manuscrit hésite sur la définition : « $f$
conservatif $\Leftarrow$ (déf) $f^{*}$ fidèle ($=$ conservatif) — les
deux ». Pour un foncteur exact à gauche entre topos, fidèle et conservatif
sont équivalents, ce qui justifie le « $=$ ». Il écrit ensuite $\mathrm{CoAlg}
(B, \pi) = B_{\pi}$, « $B \to B_{\pi}$ conservatif » (le mot « fidèle »
biffé), et « $B_{\pi} \xrightarrow{\simeq} B$ (th) » : la source de
l'équivalence doit être $B'$ et non $B$, comme le croquis, qui place
$B_{\pi}$ entre $B$ et $B'$, le montre ; nous avons corrigé. C'est le
théorème de descente pour les morphismes surjectifs de topos, qui sera
énoncé sous cette forme par Joyal et Tierney (1984) ; l'adjonction
$\underline{\mathrm{Hom}}_{\mathrm{Ens}}(f^{*}X, Y) \simeq
\underline{\mathrm{Hom}}_{G}(X, f_{*}Y)$ avec $f_{*}(Y) =
\underline{\mathrm{Hom}}(G, Y)$, notée au bas du feuillet, est l'exemple
du point canonique $(\mathrm{Ens}) \to BG$ du topos des $G$-ensembles,
dont l'image inverse est le foncteur d'oubli : les coalgèbres de la
comonade $Y \mapsto \underline{\mathrm{Hom}}(G, Y)$ sur les ensembles sont
les $G$-ensembles, et c'est le cas le plus simple du théorème.} Le dernier alinéa dit l'usage prévu :
un topos $B$ étant donné, on cherche un topos $B'$ ayant assez de points
et un morphisme $B \to B'$ auquel appliquer le critère, avec le croquis
$C \subset \widehat{C}$. C'est l'amorce d'une question constante dans le
cahier des anneaux : un sous-topos « a-t-il assez de points », auquel cas
ses propriétés se vérifient sur les modèles ensemblistes.
\pagerange{36}{50}
\section*{Le cahier des axiomes d'anneaux et de leurs topos classifiants (pages 36 à 50)}

\pagerange{36}{38}
\subsection*{Le dispositif}

Un cahier à spirale, quadrillé, fait fonctionner la machine des pages 17 à
22 sur l'exemple qui la justifie : les anneaux commutatifs et leurs
propriétés. Les conventions sont fixées une fois pour toutes :
\begin{itemize}
\item $R$ est la catégorie des schémas affines de type fini sur
$\mathbb{Z}$, opposée à celle des algèbres de type fini sur $\mathbb{Z}$ ;
elle a les limites projectives finies ;
\item $X_{\mathrm{ann}} = \widehat{R}$ est le \emph{topos classifiant des
anneaux} : un anneau $A$ d'un topos $\varepsilon$ est un foncteur exact à
gauche $R \to \varepsilon$, $X \mapsto X(A)$, i.e. un morphisme de topos
$\varepsilon \to \widehat{R}$ ; les anneaux ensemblistes sont les points de
$\widehat{R}$, et $T(\mathrm{Ens}) = \mathrm{Ind}(R^{\circ})$ est la catégorie
des anneaux, limite inductive filtrante de ses sous-anneaux de type
fini ;\footnote{C'est le cas $C = R$ de l'équivalence $T_{0}(E) \simeq
\underline{\mathrm{Hom}}_{\mathrm{top}}(E, \widehat{C})^{\circ}$ de la page 17.
La formule $T(\mathrm{Ens}) = \mathrm{Ind}(C^{\circ})$ est écrite, à
l'encre bleue, page 108 ; celle du cahier est $\tau = \mathrm{Ind}(S)$
(page 90). L'objet « anneau universel » est $A = \mathrm{Spec}\,
\mathbb{Z}[t]$, que la page 93 note $O = E'_{\mathbb{Z}} = \mathrm{Sp}
(\mathbb{Z}[t])$.}
\item pour une propriété $P$ d'anneaux, $X_{P} \subset X_{\mathrm{ann}}$ est
le sous-topos classifiant les anneaux ayant la propriété $P$, et l'on
cherche à l'écrire comme $\widetilde{R}_{\mathrm{top}\,P}$, catégorie des
faisceaux sur $R$ pour une topologie « $\mathrm{top}\,P$ », ou, mieux,
comme $\widehat{R_{P}}$, préfaisceaux sur une catégorie $R_{P}$ déduite de
$R$ par localisation.
\end{itemize}

Le principe qui donne les topologies est énoncé une fois, page 38, au cas
« anneau local », et appliqué partout : un axiome de la forme « pour toute
section $f$ de $A$ sur $U$, telle famille de $U$ est couvrante » se vérifie
sur le \emph{cas universel} $U = A$, $f$ = la section identique, et la
famille couvrante universelle engendre la topologie. Le sous-topos $X_{P}$
est alors $\widetilde{R}_{\mathrm{top}\,P}$, et $P$ est vérifiée par $A$
dans $\varepsilon$ si et seulement si le morphisme $\varepsilon \to
\widehat{R}$ se factorise par $\widetilde{R}_{\mathrm{top}\,P}$, i.e. si le
foncteur $X \mapsto X(A)$ transforme familles couvrantes en familles
épimorphiques.\footnote{C'est le mécanisme général : le topos des
faisceaux sur $(R, J)$ classifie les foncteurs exacts à gauche
\emph{$J$-continus} $R \to \varepsilon$ (Diaconescu). Lorsque $R$ est
l'opposé des anneaux de présentation finie, un tel foncteur est un anneau
$A$ et la continuité dit exactement que $X'(A) \to X(A)$ est surjectif pour
$X' \to X$ couvrant. Chaque ligne du cahier est donc une \emph{définition}
de la topologie par la propriété, et l'énoncé $X_{P} \simeq \widetilde{R}
_{\mathrm{top}\,P}$ y est tautologique ; le contenu est dans
l'identification de la topologie à une topologie connue, ou du sous-topos à
un topos de préfaisceaux.}

Une carte des implications ouvre le cahier (page 36) : local $\Rightarrow$
anneau $\Leftarrow$ réduit $\Leftarrow$ intègre $\Leftarrow$ corps
$\Rightarrow$ pseudo-corps $\Rightarrow$ réduit et à spectre compact ;
strictement local $\Rightarrow$ local.\footnote{Le feuillet dessine aussi
une flèche de « spectre compact » vers « anneau réduit », que la
transcription rend dans ce sens ; l'implication est fausse (un anneau
local artinien non réduit a un spectre réduit à un point), et nous ne la
reportons pas. Le second réseau du feuillet, illisible, redistribue les
mêmes noms avec quelques autres : spectre connexe, spectre irréductible,
intègre normal, régulier, ultra local, complet.}

\pagerange{37}{37}
\subsection*{Sous-topos ouverts du classifiant}

La page 37 rappelle comment les sous-topos ouverts de $B = \widehat{R}$ se
lisent en termes de structures. Un ouvert de $B$ est un sous-objet $S$ de
l'objet final $e_{B}$, i.e. un crible $C_{R}$ de $R$ ; le sous-topos ouvert
correspondant est $B' = B_{/S} \xrightarrow{i} B$, avec $i_{!}(\mathbb{1}
_{B'}) = S$, et il classifie la sous-structure $T'$ de $T = T_{B}$ définie
par
\[
T'(\varepsilon) = \{\, f : \varepsilon \to B \mid f^{*}(S) \simeq
e_{\varepsilon} \,\},
\]
c'est-à-dire les anneaux $A$ tels que la « condition ouverte » $S$ soit
vérifiée partout. Les exemples 16) de la page 43 ($n$ inversible, $n =
0$) sont de ce type.

\pagerange{38}{43}
\subsection*{La liste}

Nous donnons les vingt-deux articles dans l'ordre du cahier, avec pour
chacun l'axiome, le cas universel, et le topos classifiant proposé.

\begin{enumerate}
\item[1)] \emph{Anneau} (commutatif unitaire) : $X_{\mathrm{ann}} =
\widehat{R}$.

\item[2)] \emph{Anneau local} : pour toute section $f$ de $A$ sur $U$, $U$
est réunion des ouverts $U_{f}$ et $U_{1-f}$ ; cas universel $U = A$, $f$
la diagonale. $X_{\mathrm{anloc}} = \widetilde{R}_{\mathrm{zar}}$, le
\emph{topos de Zariski}.\footnote{C'est le théorème de Hakim, \emph{Topos
annelés et schémas relatifs} (1972) : le gros topos de Zariski classifie
les anneaux locaux. La date de l'inventaire ne permet pas de dire si le
cahier lui est antérieur ; il ne cite personne. La formulation par le
recouvrement $\{U_{f}, U_{1-f}\}$ du cas universel est celle qui rend
l'énoncé évident : l'axiome des anneaux locaux dit précisément que $f$ ou
$1-f$ est inversible localement.}

\item[4)] \emph{Anneau strictement local} : local, et tout revêtement étale
surjectif $U' \to U$ a localement une section. Cas universel $U =
\mathbb{A}^{n}$, $U' = \mathrm{Spec}(\mathcal{O}_{\mathbb{A}^{n}}[T]/(F))_{F'}$
avec $F(T) = T^{n} + a_{1}T^{n-1} + \cdots + a_{n}$ le polynôme unitaire
universel. $X_{\mathrm{anstrloc}} = \widetilde{R}_{\mathrm{et}}$, le topos
étale.\footnote{Anneaux strictement henséliens ; c'est aussi dans Hakim, et
c'est le théorème que Wraith reprendra (\emph{Generic Galois theory of local
rings}, 1979). Le cas universel donné — les revêtements étales standards
définis par un polynôme unitaire et sa dérivée — est exactement le
« étale local » de SGA 1.}

\item[3)] \emph{Anneau local hensélien} : $X_{\mathrm{anlochen}} =
\widetilde{R}_{\mathrm{hen}}$, pour la topologie où $(X_{i} \to X)$ est
couvrante si et seulement si $X_{i}(A) \to X(A)$ est surjectif pour tout
anneau local hensélien $A$. La définition est ici tautologique ; le cas
universel est celui de 4) sans la condition de séparabilité.

\item[5)] \emph{Anneau ultralocal} : tout recouvrement fppf $X' \to X$ de
$R$ induit une surjection $X'(A) \to X(A)$. $X_{\mathrm{anultloc}} =
\widetilde{R}_{\mathrm{fppf}}$, cas universel $U = \mathbb{A}^{n}$, $U' =
\mathrm{Spec}\,\mathcal{O}_{\mathbb{A}^{n}}[T]/(F)$. Le mot est du cahier
et n'a pas de nom reçu ; la \emph{Question} posée (tout anneau de type
fini a-t-il une $\mathcal{O}$-algèbre ultralocale qui soit limite
inductive d'algèbres fidèlement plates de type fini ?) demande si le topos
fppf a assez de points, et reste ouverte sur la page.\footnote{Elle a une
réponse positive, comme tout topos cohérent (Deligne, SGA 4 VI 9) — le
topos fppf sur $R$ est cohérent, ses recouvrements étant finis. Mais la
description des points comme « anneaux ultralocaux » ne les identifie à
aucune classe classique : un corps algébriquement clos est ultralocal
(article 15), un anneau strictement hensélien ne l'est pas en général.}

\item[6)] \emph{Anneau réduit} : $\mathcal{O}_{A} \simeq \varinjlim_{n}
\mathrm{Ker}(x \mapsto x^{n})$ ; cas universel : les $V(f^{n}) \to V(f)$.
$X_{\mathrm{red}} \simeq \widetilde{R}_{\mathrm{topred}} \simeq
\widehat{R_{\mathrm{red}}}$, où topred est la topologie où $(X_{i} \to Y)$
est couvrante si l'une des $X_{i}$ a une section sur $Y_{\mathrm{red}}$, et
$R_{\mathrm{red}}$ la catégorie des schémas réduits de $R$, qui est aussi
la catégorie de fractions de $R$ pour les immersions $X_{\mathrm{red}} \to
X$.\footnote{L'égalité $X_{\mathrm{red}} \simeq \widehat{R_{\mathrm{red}}}$
est le fait que « réduit » est une théorie \emph{cartésienne} (l'axiome
$x^{n} = 0 \Rightarrow x = 0$ est une implication entre équations, sans
disjonction ni existentiel), dont le classifiant est donc un topos de
préfaisceaux, sur l'opposé des anneaux réduits de présentation finie. Que
$R_{\mathrm{red}}$ ait des limites finies — non induites par $R$ : le
produit fibré de deux schémas réduits n'est pas réduit — vient de ce que
$X \mapsto X_{\mathrm{red}}$ est adjoint à droite de l'inclusion. Le
\emph{N.B.} de la page 38 s'en inquiète (« $X_{\mathrm{red}}$ peut être
\ldots, ce qui se voit sur $R_{\mathrm{red}}$ \ldots\ des lim. finies ») ;
c'est bien la remarque à faire.}

\item[7)] \emph{Anneau intègre} : a) le morphisme $u : A \amalg A \to D$
vers le produit fibré $D = \{(x,y) : xy = 0\}$ est un épimorphisme
($xy = 0$ implique $x = 0$ ou $y = 0$, localement) ; b) $\mathrm{Ker}(0_{A},
1_{A} : e \rightrightarrows A) = \emptyset$ ($0 \neq 1$). $X_{\mathrm{int}}
\simeq \widetilde{R}_{\mathrm{topint}}$, topint étant la topologie où
$(X_{i} \to Y)$ est couvrante si, pour toute composante irréductible
$Y_{\alpha}$ de $Y$ munie de sa structure réduite, l'une des $X_{i}$ a une
section au-dessus de $Y_{\alpha}$. Les faisceaux sont les $F \in
\widehat{R}$ tels que $F(X) \to \prod_{i} F(X_{i}) \rightrightarrows
\prod_{i,j} F(X_{i} \times_{X} X_{j})$ soit exacte pour la famille des
composantes irréductibles.\footnote{Ici une disjonction apparaît, et le
classifiant n'est plus un topos de préfaisceaux : « intègre » est une
théorie \emph{cohérente}. La topologie décrite est la traduction de la
disjonction : recouvrir $\mathrm{Spec}\,A$ par les $V(x)$ et $V(y)$
lorsque $xy = 0$.}

\item[8)] \emph{Anneau irréductible} : $A_{\mathrm{red}}$ est intègre ; la
condition s'écrit avec le nilradical $N$ comme en 7), et donne
$X_{\mathrm{irr}} \simeq \widetilde{R}_{\mathrm{topirr}}$, la topologie où
$(X_{i} \to X)$ est couvrante si pour toute composante irréductible $Y$ de
$X$ l'une des $X_{i}$ a une section sur $Y_{\mathrm{red}}$. Le \emph{N.B.}
qui suit est le premier signal de la limite de la méthode : \emph{cette
topologie n'est pas quasi-compacte} — un schéma à une infinité de
composantes irréductibles a des recouvrements sans sous-recouvrement
fini —, de sorte que $X_{\mathrm{irr}}$ n'est pas un topos
cohérent.\footnote{Un schéma affine de type fini sur $\mathbb{Z}$ est
noethérien et n'a qu'un nombre fini de composantes irréductibles ; la
non-quasi-compacité vient des recouvrements par des familles infinies
$(X_{i})$ dont chaque composante ne rencontre qu'un nombre fini. La
remarque est juste et elle est importante : l'énoncé « irréductible »,
$\forall f, g\ (fg \text{ nilpotent} \Rightarrow f \text{ nilpotent} \vee
g \text{ nilpotent})$, est géométrique mais non cohérent, à cause de la
disjonction infinie cachée dans « nilpotent ».}

\item[10)] \emph{Anneau compact}, i.e. à spectre compact (Hausdorff) :
pour toute section $f$, il existe un unique idempotent $e$ tel que $ef$
soit inversible dans $eA$ et $(1-e)f$ nilpotent. Cas universel : pour
chaque $n$, soit $D_{n} \subset A \times A \times A$ le lieu des $(g, f, e)$
avec $e^{2} = e$, $f^{n}(1-e) = 0$, $g(1-e) = 0$, $f^{n}g = e$ ; la
projection $D_{n} \to A$ sur $f$ est un monomorphisme ($e$ et $g$ sont
déterminés par $f$), et la condition est que la famille $(D_{n} \to
A)_{n}$ soit couvrante.\footnote{La page 42 écrit les équations de $D_{n}$
sans l'exposant $n$, qui est pourtant ce qui distingue $D_{n}$ de $D_{1}$ ;
la ligne suivante, qui fait couvrir $A$ par les $D_{n}$, exige la
puissance, et nous l'avons rétablie.} $X_{\mathrm{cpct}} \simeq
\widetilde{R}_{\mathrm{top\,cpct}}$ : $(X_{i} \to X)$ couvrante si
$X_{i}(\mathcal{O}) \to X(\mathcal{O})$ est surjectif pour tout anneau
artinien local $\mathcal{O}$. Marge : « topologie non cohérente ! A-t-elle
assez de points ?? », et « encore un topos non cohérent (bien que sous-topos
d'un topos cohérent) ».\footnote{Les articles 9) et 10) portent les numéros
du cahier ; la numérotation y est reprise deux fois (le 12) de la page 41
devient le 13) de la page 42). « Spectre compact » signifie spectre de
dimension zéro, i.e. anneau dont tout premier est maximal ; la condition
donnée est la caractérisation par les idempotents, et la « quantification
sur $n$ » (le $f^{n}$ nilpotent) est ce qui rend la théorie non cohérente,
comme en 8).}

\item[11)] \emph{Corps} : $A^{*} \amalg e \to A$ (l'inclusion des
inversibles et la section nulle) est épimorphique — tout élément est
inversible ou nul. $X_{\mathrm{corps}} \simeq \widetilde{R}_{\mathrm{top\,
corps}}$, la topologie où les familles couvrantes sont celles qui ont des
sections sur chaque fibre.\footnote{Le cahier écrit « $A^{*} \amalg A$ »,
avec pour seconde flèche $0_{A}$ ; la page 94 écrit $e \sqcup O^{*} \to O$,
qui est ce que le sens exige, et nous suivons la page 94. L'axiome
$\forall x\ (x = 0 \vee \exists y\ xy = 1)$ est cohérent ; avec « $0 \neq
1$ » c'est la théorie géométrique des corps, dont le classifiant est le
topos des faisceaux sur les anneaux de présentation finie pour la
topologie engendrée par les recouvrements $\{\mathrm{Spec}\,A_{f},
\mathrm{Spec}\,A/f\} \to \mathrm{Spec}\,A$ — c'est ce que « sections sur
chaque fibre » veut dire. Les pages 92 à 95 y reviennent longuement.}

\item[12)--13)] \emph{Corps} $=$ compact et intègre $=$ pseudo-corps intègre
$=$ pseudo-corps local, et \emph{pseudo-corps} $=$ anneau compact réduit :
pour toute section $f$, $A \xrightarrow{\sim} A_{f} \times A/fA$, i.e. il
existe un unique idempotent $e$ avec $ef$ inversible dans $eA$ et $(1-e)f =
0$ ; le cas universel dit que $i_{1} : D_{1} \hookrightarrow A$ est un
isomorphisme. $X_{\mathrm{ps.corps}} \simeq \widetilde{R}_{\mathrm{top\,ps\,
corps}} \simeq \widehat{R_{\mathrm{cons}}}$, où $R_{\mathrm{cons}}$ est la
catégorie des « schémas compacts réduits associés aux schémas de type fini
sur $\mathbb{Z}$ », localisée de $R$ pour les flèches à extensions
résiduelles triviales.\footnote{Un pseudo-corps est un anneau
\emph{absolument plat}, ou \emph{régulier au sens de von Neumann} : tout
$x$ s'écrit $x = x^{2}y$. La condition $A \simeq A_{f} \times A/fA$ pour
tout $f$ en est une caractérisation, et l'unicité de l'idempotent en fait
une théorie cartésienne, d'où le classifiant en préfaisceaux
$\widehat{R_{\mathrm{cons}}}$. La catégorie $R_{\mathrm{cons}}$ est
vraisemblablement l'opposée de celle des anneaux absolument plats de
présentation finie ; l'« anneau absolument plat universel » $T(A)$
associé à un anneau $A$, dont le spectre est $\mathrm{Spec}\,A$ muni de la
topologie constructible, est l'objet de la thèse de J.-P. Olivier
(Montpellier, 1970--71), et le nom « cons » y renvoie. Le cahier ne cite
pas Olivier.}

\item[14)] \emph{Corps séparablement clos} $=$ corps strictement hensélien :
$X \simeq \widetilde{R}_{\mathrm{top\,corps\,sépclos}} \simeq
\widetilde{(R_{\mathrm{cons}})}_{\mathrm{et}}$, faisceaux sur $R_{\mathrm{cons}}$
pour la topologie étale.

\item[15)] \emph{Corps algébriquement clos} $=$ corps ultralocal : $X \simeq
\widetilde{(R_{\mathrm{cons}})}_{\mathrm{fppf}}$, la topologie fppf sur
$R_{\mathrm{cons}}$ étant celle des familles finies surjectives.\footnote{Le
classifiant des corps algébriquement clos est un exemple classique de la
logique catégorique (Makkai--Reyes 1977, Johnstone) ; sa description
comme faisceaux sur les anneaux absolument plats de présentation finie
pour les familles finies surjectives est, à la lettre près, celle qui y
est donnée.}

\item[16)] \emph{Anneaux annulés par $n$}, \emph{où $n$ est inversible},
\emph{$\Lambda$-algèbres} : trois sous-topos \emph{ouverts}. Pour $n1_{A} =
0$ : $X_{\mathbb{Z}/n} \simeq \widetilde{R}_{\mathrm{top}\,n\text{-tors}}
\simeq \widehat{R_{\mathbb{Z}/n}}$, $R_{\mathbb{Z}/n}$ la sous-catégorie
pleine des $\mathbb{Z}/n$-schémas ; pour $n$ inversible, de même avec
$\mathbb{Z}[1/n]$. Plus généralement (pages 43 à 45), soit $S \to e_{R}$
un monomorphisme de $\mathrm{Pro}\,R$, i.e. un épimorphisme d'anneaux
$\mathbb{Z} \to \Lambda$ ; les $\Lambda$-algèbres de $\varepsilon$ sont les
anneaux $A$ dont la structure de $\Lambda$-algèbre existe, et elle est
alors unique ; $R_{/S} \subset R$ est la sous-catégorie pleine des schémas
qui sont des $\Lambda$-schémas, avec $X \mapsto X \times S$ pour adjoint,
et
\[
X_{\Lambda\text{-alg}} \simeq \widetilde{R}_{\Lambda\text{-top}} \simeq
\widehat{R_{\Lambda}} = \widehat{R_{/S}},
\]
sous-topos ouvert de $\widehat{R}$ défini par $S$. Exemples : $X_{\mathbb{Z}
/n}$, $X_{\mathrm{car}\,p}$, $X_{\mathbb{Z}[1/n]}$, et $X_{\mathbb{Q}} =
X_{\mathrm{car}\,0} = \bigcap_{n \geqslant 2} X_{\mathbb{Z}[1/n]}$.\footnote{
$\mathbb{Z} \to \Lambda$ épimorphique signifie que $\Lambda$ est un
quotient d'un localisé de $\mathbb{Z}$, i.e. $\mathbb{Z}/n$, $\mathbb{Z}
[1/n]$, $\mathbb{Q}$ ou leurs produits finis ; c'est ce que le cahier
nomme « mono $S \to e_{R}$ dans $\mathrm{Pro}\,R$ ». L'unicité de la
structure de $\Lambda$-algèbre est un cas particulier du fait qu'un
épimorphisme d'anneaux $\mathbb{Z} \to \Lambda$ rend $\Lambda \otimes_{
\mathbb{Z}} \Lambda \to \Lambda$ bijectif.} La page 45 traite le cas d'une
famille filtrante décroissante $S_{\alpha}$ de sous-objets de $e_{R}$ : la
sous-structure $\bigcap_{\alpha} T_{S_{\alpha}}$ est classifiée par $\bigcap
_{\alpha} \widehat{R_{/S_{\alpha}}} = \widetilde{\varprojlim_{\alpha}
R_{/S_{\alpha}}}$, la catégorie $\varprojlim_{\alpha} R_{/S_{\alpha}}$
étant une catégorie de fractions de $R$ pour les flèches $i$ telles que $i
\times S$ soit inversible ; c'est le cas de $X_{\mathbb{Q}}$, qui est un
sous-topos non ouvert, intersection d'ouverts.

\item[17)--19)] \emph{Anneau parfait de caractéristique $p$} ($p1_{A} = 0$
et $x \mapsto x^{p}$ bijectif) : $X_{p\text{-parf}} \simeq \widehat{R_{p\text{-parf}}}$,
$R_{p\,\mathrm{parf}}$ la catégorie des clôtures parfaites des schémas de
type fini sur $\mathbb{F}_{p}$ ; \emph{compact parfait} et \emph{corps
parfait}, en croisant avec 13) et 11).\footnote{Le passage à la clôture
parfaite $X \mapsto X_{\mathrm{parf}}$ est adjoint à l'inclusion des schémas
parfaits, comme la réduction en 6) ; « parfait » est cartésien
($\forall x\, \exists! y\ y^{p} = x$), d'où le classifiant en préfaisceaux.
Le \emph{N.B.} « parfait $\Rightarrow$ réduit » est immédiat : Frobenius
étant injectif, $x^{p} = 0$ entraîne $x = 0$.}
\end{enumerate}

\pagerange{46}{49}
\subsection*{Caractéristique, et le recollement de deux sous-topos}

L'article 20), \emph{anneau à caractéristique} (localement, $A$ est de
caractéristique $p > 0$ ou de caractéristique $0$), est celui où la
méthode achoppe et où le cahier fait sa meilleure observation. La
condition s'écrit : la famille
\[
\bigl\{\, V(p \cdot 1_{A}) \hookrightarrow e_{\varepsilon}\ (p \text{
premier}),\quad \textstyle\bigcap_{n \geqslant 1} (e_{\varepsilon})_{n1_{A}}
\hookrightarrow e_{\varepsilon} \,\bigr\}
\]
est épimorphique, et le second membre est une \emph{intersection infinie}
de sous-objets : « donc ne donne pas l'existence d'un topos classifiant »,
c'est-à-dire pas par la voie ouverte (les conditions ouvertes de la page
37). Ce que l'on peut dire est que $\varepsilon \to \widehat{R}$ se factorise
localement sur $\varepsilon$ par la « réunion » $V = X_{\mathrm{car} > 0}$
des sous-topos ouverts $\widehat{R_{\mathbb{Z}/p}}$ et par le sous-topos $W =
X_{\mathbb{Q}}$, qui n'est pas ouvert. On pose $X_{\mathrm{car}} =
\widetilde{R}_{\mathrm{top.car}}$, où « top car » est la borne inférieure
des « top car $p$ » : $(X_{i} \to X)$ est couvrante si pour tout $p$ il
existe $i(p)$ tel que $(X_{i(p)})_{p} \to X_{p}$ ait une section. C'est un
\emph{topos non cohérent} : les familles $(X_{p} \to X)_{p}$ sont
couvrantes et, si $X_{\mathbb{Q}} \neq \emptyset$, n'ont pas de
sous-famille finie couvrante.\footnote{La théorie « $A$ a une
caractéristique » s'écrit $\top \vdash \bigvee_{p} (p = 0) \vee
\bigwedge_{n} (n \text{ inversible})$ ; la conjonction infinie n'est pas
géométrique, et c'est ce que le cahier constate en disant que
l'intersection $\bigcap_{n}$ ne donne pas de condition ouverte. Le
sous-topos $\sup(V, W)$ existe néanmoins, comme borne supérieure de deux
sous-topos, et $X_{\mathrm{car}}$ est ce sup ; la question est de le
décrire.}

Le problème général est alors posé (pages 47 et 48) : $B$ un topos, $V$ et
$W$ deux sous-topos, décrire le sous-topos $Z = \sup(V, W)$ classifiant les
structures qui sont \emph{localement} de type $V$ ou de type $W$. On
cherche $Z$ muni de deux ouverts $Z' = Z_{/e'}$, $Z'' = Z_{/e''}$ avec $Z' \subset
V$, $Z'' \subset W$ et $e_{Z} = \sup(e', e'')$ ; alors $Z' = Z \cap V$ est un
ouvert de $V$ et $Z''$ un ouvert de $W$, et la condition nécessaire et
suffisante trouvée page 48 est qu'il existe un ouvert $Z'$ de $V$ et un
ouvert $Z''$ de $W$ tels que $Z'' \wedge V$ soit un ouvert de $V$, $Z' \wedge
W$ un ouvert de $W$, et
\[
V = \sup(Z', Z'' \wedge V), \qquad W = \sup(Z'', Z' \wedge W) ;
\]
$Z$ s'obtient alors en \emph{recollant} $Z'$ et $Z''$ le long de leur
intersection. Dans le cas $B = \widehat{R}$, $V = X_{\mathrm{car} > 0}$, $W =
X_{\mathrm{car}\,0}$, la condition est satisfaite, $V \wedge W$ étant l'ouvert
correspondant à l'anneau nul : « donc OK ».\footnote{C'est le recollement
de topos le long d'un ouvert commun, i.e. la construction du topos
classifiant d'une théorie « locale » à partir de ceux de ses deux
branches ; les pages 47--48 sont de brouillon et l'énoncé ci-dessus est
notre mise en forme de ce que la page 48 conclut. Le point délicat, que
le cahier voit, est que $W = X_{\mathbb{Q}}$ n'est pas ouvert dans $B$ mais
que $V \wedge W$ (le lieu où $p = 0$ et $p$ inversible, i.e. l'anneau nul)
l'est.} La page 49 conclut par une description en préfaisceaux :
\[
X_{\mathrm{car}} \simeq \widehat{R_{\mathrm{car}}},
\]
$R_{\mathrm{car}}$ étant la « somme amalgamée » sous le schéma vide des
catégories $R_{p}$ de schémas de type fini sur les corps premiers,
$\mathbb{Q}$ compris.\footnote{Cette dernière identification est énoncée
sans démonstration, et elle demande une précaution que le cahier ne prend
pas : les schémas de type fini sur $\mathbb{Q}$ ne sont pas dans $R$ (ils
sont dans $\mathrm{Pro}\,R$, comme $X_{\mathbb{Q}} = \bigcap_{n}
X_{\mathbb{Z}[1/n]}$ le dit page 44). La catégorie $R_{\mathrm{car}}$ n'a
pas d'objet final — $\mathrm{Spec}\,\mathbb{F}_{p}$ et $\mathrm{Spec}\,
\mathbb{Q}$ n'ont pas de produit non vide — et ses foncteurs plats vers
les ensembles sont les anneaux de caractéristique $p$ ou $0$, l'anneau nul
excepté, ce qui est bien la description voulue des points ; que
$\widehat{R_{\mathrm{car}}}$ soit le sous-topos $\sup(V, W)$ de
$\widehat{R}$ est plausible et non établi ici.}

Les articles 21) et 22) élargissent. 21) : reprendre 1) à 19) au-dessus
d'un sous-topos $Z(p)$ de $\widehat{R}$ pour chaque caractéristique, et
recoller. 22) \emph{Anneaux quasi-connexes de rang $\leqslant n$} : le
spectre a au plus $n$ composantes connexes, i.e. localement $A \simeq A_{1}
\times \cdots \times A_{n}$ avec les $A_{i}$ connexes, i.e. il existe
localement des idempotents orthogonaux $e_{1}, \ldots, e_{n}$ de somme $1$
avec $A/\sum_{j \neq i} e_{j}A$ connexe. On forme le sous-objet $\Pi_{n}
\subset A^{n}$ des tels $n$-uples, et pour chaque $j$ le plus grand
sous-objet $\Pi_{n,j}$ sur lequel $A_{j}$ est connexe ; la condition est que
$\Pi'_{n} = \bigcap_{j} \Pi_{n,j} \to e_{\varepsilon}$ soit épimorphique. La
page 50 note que la formation de $\Pi_{n,j}$ ne commute pas visiblement aux
changements de topos de base, « donc il n'est pas clair qu'on a un topos
modérateur », et propose la structure plus riche — un $\mathfrak{S}_{n}$-torseur
$P$ et une décomposition de $A_{P}$ en $n$ facteurs connexes — pour
laquelle la classification est claire.\footnote{« Connexe » ($e^{2} = e
\Rightarrow e = 0 \vee e = 1$, avec $0 \neq 1$) est cohérent, et « au plus
$n$ composantes » aussi ($\bigvee_{i} e_{i} = 0$ pour $n+1$ idempotents
orthogonaux) ; ces théories ont donc des topos classifiants, sous-topos de
$\widehat{R}$ (Makkai--Reyes). Ce que la page 50 voit juste est la
différence entre la théorie \emph{avec} la décomposition pour structure —
dont le classifiant est le produit fibré de $n$ copies du classifiant des
anneaux connexes, muni de son action de $\mathfrak{S}_{n}$ — et la théorie
qui l'affirme seulement par un existentiel, dont le classifiant est le
quotient de la première par $\mathfrak{S}_{n}$. Le « pseudo-torseur des
choix » qu'elle décrit est le lien entre les deux.}

\pagerange{41}{41}
\subsection*{La question des points}

Le \emph{N.B.} de la page 41 formule la contrainte qui pèse sur tout le
cahier : les sous-topos 1) à 11) ont des points — ils sont même définis par
leurs points, puisque chacun a été défini par une condition sur les
modèles ensemblistes —, mais il n'est pas clair qu'une \emph{intersection}
de plusieurs d'entre eux en ait assez ; pour vérifier que deux
intersections coïncident, ou s'incluent, il ne suffit donc pas de le
vérifier sur les structures ensemblistes, « sauf si on vérifie à chaque
fois que les intersections précédentes ont assez de points ! ». C'est la
raison des marges « topologie non cohérente, a-t-elle assez de points ?? »
des articles 8) et 10).\footnote{Le théorème de Deligne (SGA 4 VI 9, 1972)
dit qu'un topos cohérent a assez de points ; il règle tous les articles
cohérents, et les intersections finies de sous-topos cohérents d'un topos
cohérent, qui sont cohérentes. Il ne règle pas 8), 10) et 20), qui ne le
sont pas, et pour lesquels la question reste, sur ces pages, ouverte.
Elle l'est toujours en général : un topos non cohérent peut n'avoir aucun
point.}

\pagerange{51}{52}
\section*{Propriétés typiques (pages 51 et 52)}

Deux feuillets de questions demandent, pour des propriétés de modules et
d'algèbres, si elles sont \emph{typiques} — le mot du manuscrit pour :
définissables par un sous-topos du classifiant, ce qui exige d'être stable
par limites inductives filtrantes et par changement de topos.\footnote{En
termes d'aujourd'hui, une propriété typique est une propriété
\emph{géométrique} : exprimable par une formule de logique géométrique
(conjonctions finies, disjonctions quelconques, existentiels), qui sont
exactement celles que préservent les images inverses des morphismes de
topos. La stabilité par limites filtrantes en est la trace sur les points.}
Les réponses :
\begin{itemize}
\item $M$ \emph{plat} est typique : $N \mapsto M \otimes N$ exact se ramène à
l'injectivité de $(\sum f_{i}A) \otimes M \to M$ pour toute famille finie de
sections, et il suffit de l'exiger pour les sections universelles $t_{i}$
sur $U_{n} = A^{n}$.\footnote{C'est le critère de platitude par les
relations linéaires (Bourbaki, Lazard), et il est bien géométrique.}
\item $M$ \emph{fidèlement plat} : « est-ce typique ? », encerclé ; la
condition $M \otimes N \neq 0$ pour tout $N \neq 0$ n'est pas de la forme
voulue.
\item Algèbres de type fini, de présentation finie, lisses, étales, nettes :
\emph{pas typiques} ; radicielles (monomorphiques), plates, entières,
formellement nettes : oui ; surjectives : « ?! ».
\item Modules localement libres, libres, projectifs, de type fini, de
présentation finie, monogènes : les trois derniers « ne passent pas à la
limite inductive filtrante, donc pas typiques ».
\end{itemize}
Un encadré démontre que $M$ \emph{inversible} ($M \otimes N \simeq A$ pour un
$N$) entraîne $M$ de type fini : $M = \varinjlim M_{i}$ sur ses sous-modules
de type fini, $M \otimes N = \varinjlim M_{i} \otimes N$ et $M_{i} \otimes N
\subset A$ par platitude, donc l'un des $M_{i} \otimes N$ est $A$ et $M_{i} =
M$. Mais « $M$ inversible » n'est pas exprimable par limites inductives et ne
commute pas au changement de base, à cause du foncteur $\underline{\mathrm{Hom}}$ ;
la propriété typique voisine est « localement libre de rang un », dont le
classifiant est le topos $X_{\mathrm{mod}} = \widehat{\mathrm{Mod}_{\mathrm{pf}}
^{\circ}}$ des modules de présentation finie sur les algèbres de
présentation finie, restreint par une condition ouverte.\footnote{Le
manuscrit écrit $X_{\mathrm{mod}} = \widehat{\mathrm{Mod}^{\circ}_{\mathrm{pf}}}$
pour le classifiant des couples (anneau, module), qui est une théorie
cartésienne ; la fin de la phrase est nôtre. Les feuillets sont une liste
de questions et n'affirment rien de plus que ce qui est reporté ici.}

\pagerange{54}{61}
\section*{Catégories accessibles (pages 54 à 61)}

\pagerange{54}{55}
\subsection*{Le relevé}

Des feuillets jaunis, d'une autre plume, recopient et complètent les
énoncés de SGA 4, exposé I, § 8 et § 9 — les renvois numérotés en marge
(8.7.5, 8.9.5, 8.13.2, 9.1--9.3, 9.6, 9.12, 9.16, 9.22) sont ceux de
l'exposé — sur les \emph{catégories accessibles}. C'est l'outillage des
sections « Exemples » et « Compléments » : la catégorie $S = T(\mathrm{Ens})$
des modèles d'une théorie y sera reconnue comme une catégorie de ce type,
et c'est ce qui permettra de reconstruire la théorie à partir de ses
modèles.\footnote{Ces notions sont celles que Gabriel et Ulmer publient en
1971 (\emph{Lokal präsentierbare Kategorien}) sous le nom de catégories
\emph{localement présentables}, et que Makkai et Paré généraliseront en
1989 (catégories \emph{accessibles}). Le relevé n'utilise que SGA 4 ; la
terminologie de l'exposé I — « accessible » pour ce que Gabriel--Ulmer
appellent présentable — est conservée ici, et le dernier corollaire des
pages 60--61 est, mot pour mot, la caractérisation de Gabriel--Ulmer.}

Les définitions : un ensemble ordonné filtrant est \emph{grand devant}
$\pi$ si toute partie de cardinal $\leqslant \pi$ est majorée ; une
catégorie satisfait $L_{\pi}$ si les limites inductives indexées par de tels
ensembles y existent ; un objet $X$ est \emph{$\pi$-accessible} si
$\mathrm{Hom}(X, -)$ commute à ces limites ; une catégorie $E$ est
\emph{$\pi$-accessible} si elle satisfait $L_{\pi}$ et admet une petite
sous-catégorie strictement génératrice formée d'objets $\pi$-accessibles,
et \emph{accessible} s'il existe un tel $\pi$. Pour $\mathcal{C}$ petite,
$\mathrm{Ind}_{\pi}(\mathcal{C})$ est la sous-catégorie pleine de
$\mathrm{Ind}(\mathcal{C})$ des ind-objets indexés par des ensembles
ordonnés grands devant $\pi$, et $E_{\pi}$ la sous-catégorie pleine des
objets $\pi$-accessibles de $E$.

\textbf{Théorème.} \emph{Pour une catégorie $E$, sont équivalentes : a) $E$
est $\pi$-accessible ; b) $E$ est équivalente à $\mathrm{Ind}_{\pi}
(\mathcal{C})$, $\mathcal{C}$ petite ; c) $E_{\pi}$ est équivalente à une
petite catégorie et strictement génératrice. Sous ces conditions,
$E_{\pi'} \simeq \mathrm{Ind}_{\pi}(E_{\pi})_{\pi'}$ pour $\pi' \geqslant \pi$,
et les $E_{\pi'}$ forment une filtration cardinale de $E$.}

Les corollaires : un foncteur $f : E \to F$ entre catégories accessibles est
accessible si et seulement si $f \simeq \mathrm{Ind}_{\pi}(f_{\pi})$ pour un
$\pi$, si et seulement s'il est dans l'image essentielle de
$\underline{\mathrm{Hom}}(\mathcal{C}, F) \to \underline{\mathrm{Hom}}(E, F)$
pour une petite $\mathcal{C} \subset E$ ; pour $E$ $\pi$-accessible et $F$
satisfaisant $L_{\pi}$, $\underline{\mathrm{Hom}}(E, F)_{\pi} \simeq
\underline{\mathrm{Hom}}(E_{\pi}, F)$ ; pour $\mathcal{C}$ accessible, les
foncteurs $\mathcal{C} \to (\mathrm{Ens})$ exacts à gauche et accessibles
sont pro-représentables, et un foncteur $f : \mathcal{C} \to \mathcal{C}'$
admet un adjoint à gauche si et seulement s'il commute aux petites limites
projectives et est accessible ; enfin limites projectives et sections d'une
catégorie fibrée en catégories accessibles sur une petite base sont
accessibles (9.22).\footnote{Ce dernier énoncé — stabilité de
l'accessibilité par limites projectives — sera le théorème central de
Makkai--Paré ; le b) de la page 54, sur les limites inductives, est marqué
« à vérifier ». Le \emph{théorème d'adjonction} (« $f$ admet un adjoint à
gauche ssi il commute aux limites et est accessible ») est la forme que
prend, dans ce cadre, le théorème du foncteur adjoint, et c'est celle que
les pages 58 et 72 utiliseront.}

\pagerange{55}{55}
\subsection*{Accessibilité et générateurs}

La page 55 ajoute une \emph{Proposition} de stabilité — pour $F$
satisfaisant $L_{\pi}$, $\underline{\mathrm{Hom}}(E, F)_{\pi}$ est stable
par les petites limites inductives représentables dans $F$, et, si $F$ est
$\pi$-accessible, par les limites de type $I$ avec $\mathrm{card}\,I
\leqslant \pi$ — et une \emph{réciproque partielle} : une catégorie $E$
satisfaisant $L$, stable par noyaux de doubles flèches à foncteur $\mathrm{Ker}$
accessible et par produits fibrés, et admettant une petite sous-catégorie
génératrice par épimorphismes stricts universels, est accessible. D'où le
\textbf{Corollaire} : \emph{une catégorie $E$ stable par petites limites
inductives est accessible si et seulement si elle admet une petite
sous-catégorie génératrice par épimorphismes stricts universels et si le
foncteur noyau sur ses doubles flèches est accessible.} Exemples : les
topos ; toute catégorie à limites filtrantes exactes et petit générateur.
Contre-exemples : $(\mathrm{Ens})^{\circ}$ et $(\mathrm{Ab})^{\circ}$ ne
sont pas accessibles, parce que leurs foncteurs noyau ne commutent pas
aux limites inductives grandes devant $\pi$ (une limite inductive
d'épimorphismes n'est pas nécessairement un épimorphisme dans ces
catégories).\footnote{L'argument est le bon, et le fait est vrai : une
catégorie accessible complète est localement présentable, et $(\mathrm{Ens})$
ne peut être à la fois localement présentable et co-localement
présentable sans être un ensemble préordonné (Gabriel--Ulmer). Le
manuscrit l'obtient directement.}

\pagerange{56}{56}
\subsection*{Foncteurs strictement générateurs}

La page 56 est un énoncé complet, sous forme de liste, sur un foncteur
$\varphi : \mathcal{C} \to E$ d'une petite catégorie. Il définit
$\hat{\varphi} : E \to \widehat{\mathcal{C}}$, $X \mapsto \mathrm{Hom}_{E}
(\varphi(-), X)$, et, si $E$ a des petites limites inductives, $\bar{\varphi}
: \widehat{\mathcal{C}} \to E$, $F \mapsto \varinjlim_{\mathcal{C}_{/F}}
\varphi$, adjoint à gauche de $\hat\varphi$. Sont équivalentes :
\begin{enumerate}
\item[a)] pour tout $X$, $\varinjlim_{\mathcal{C}_{/X}} \varphi(Y)
\xrightarrow{\sim} X$ ;
\item[b)] $\hat\varphi$ est pleinement fidèle ;
\end{enumerate}
et l'on dit alors que $\varphi$ est \emph{strictement générateur}.
Conditions équivalentes sous hypothèses : b') si $E$ est cocomplète,
$\bar\varphi$ est un foncteur de localisation ; c') si $E$ a les limites
filtrantes et les $\mathcal{C}_{/X}$ sont filtrantes (par exemple
$\mathcal{C}$ à limites finies et $\varphi$ exact à gauche), $\mathrm{Ind}
(\mathcal{C}) \to E$ est une localisation ; c) si $\hat\varphi$ prend ses
valeurs dans $\mathrm{Ind}(\mathcal{C})$, le foncteur ${}^{i}\varphi : E \to
\mathrm{Ind}(\mathcal{C})$ est pleinement fidèle ; d) si $\varphi$ a un
adjoint $\psi$, $\psi$ est pleinement fidèle ; e) si $\varphi$ est
pleinement fidèle, $\varphi(\mathcal{C})$ engendre par épimorphismes
stricts ; f) pour toute $F$ cocomplète, la restriction $\underline{\mathrm{Hom}}'
(E, F) \to \underline{\mathrm{Hom}}(\mathcal{C}, F)$ des foncteurs
cocontinus est pleinement fidèle ; g) si $\mathcal{C}$ a les limites
inductives finies et $\varphi$ y commute et est pleinement fidèle, cette
restriction identifie $\underline{\mathrm{Hom}}'(E, F)$ aux foncteurs
$\mathcal{C} \to F$ commutant aux limites inductives finies ; h) si $E$ est
un topos, $\mathcal{C}$ stable par limites finies et $\varphi$ exact à
gauche, $\varphi$ définit $E$ comme \emph{sous-topos} de
$\widehat{\mathcal{C}}$.\footnote{Un foncteur strictement générateur est ce
qu'on appelle aujourd'hui un foncteur \emph{dense} ; $\hat\varphi$ est son
nerf, $\bar\varphi$ sa réalisation. Le point h) est la version « site » :
un topos muni d'une petite sous-catégorie dense stable par limites finies
est un sous-topos du topos des préfaisceaux sur elle, avec pour inclusion
le nerf et pour réflecteur exact à gauche la réalisation. Toutes ces
équivalences sont correctes ; e) demande la stricte épimorphicité
\emph{universelle} pour la réciproque, ce que la page 55 précise.}

\pagerange{58}{58}
\subsection*{Cocomplète à générateur dense : la question de l'accessibilité}

La page 58 caractérise les catégories $E$ vérifiant l'une des conditions
équivalentes : a) $E$ est cocomplète et admet une petite sous-catégorie
pleine strictement génératrice ; a') $E$ est cocomplète et il existe
$\varphi : \mathcal{C} \to E$ strictement générateur, $\mathcal{C}$
petite ; a'') il existe une petite $\mathcal{C}$ et un foncteur pleinement
fidèle $E \to \widehat{\mathcal{C}}$ admettant un adjoint à gauche. Dans une
telle catégorie tout foncteur contravariant vers $(\mathrm{Ens})$ qui
transforme les petites limites inductives en limites projectives est
représentable, donc $E$ a les petites limites projectives, et dans a'') le
foncteur $E \to \widehat{\mathcal{C}}$ y commute ; si de plus $E$ est
accessible, un foncteur $E \to \widehat{\mathcal{C}}$ a un adjoint à gauche
si et seulement s'il est accessible et commute aux petites limites
projectives.\footnote{L'assertion de représentabilité est vraie, mais pour
une raison qui n'est pas sur la page : une catégorie cocomplète à petite
sous-catégorie dense est \emph{totale} (Street--Walters 1978, Kelly 1986),
et une catégorie totale est « compacte » au sens d'Isbell, i.e. y
représente tout foncteur $E^{\circ} \to (\mathrm{Ens})$ transformant les
colimites en limites. La page l'énonce entre crochets, comme un fait
connu.} La \emph{Question} : une telle catégorie est-elle accessible ? Il
faut et suffit que $E \to \widehat{\mathcal{C}}$ soit accessible, i.e. que
$E$ soit stable dans $\widehat{\mathcal{C}}$ par les limites filtrantes
grandes devant un cardinal $\pi$. La marge propose $(\mathrm{Ens})^{\circ}$
comme contre-exemple.\footnote{Que $(\mathrm{Ens})^{\circ}$ ne soit pas
accessible est acquis (page 55). Qu'elle satisfasse a), i.e. que $(\mathrm{Ens})$
ait une petite sous-catégorie \emph{codense}, dépend de considérations de
théorie des ensembles (existence de cardinaux mesurables) que la page ne
soulève pas ; le contre-exemple est donc conditionnel, et la question,
sous la forme « a) implique-t-il accessible ? », a une réponse négative
que ce contre-exemple ne suffit pas à établir.} Suivent les \emph{exemples}
qui motivent le tout — a) les topos et leurs opposés, b) $T(\mathrm{Ens})$
pour $T$ une structure algébrique définissable par petites limites
projectives, c) $T(\mathcal{C})$ pour $\mathcal{C}$ accessible — et deux
\emph{conséquences} : un morphisme de structures $T \to T'$ donne un
foncteur accessible $T(\mathcal{C}) \to T'(\mathcal{C})$ commutant aux
petites limites projectives, qui a donc un adjoint à gauche (\emph{les
structures $T'$-$T$-libres}) ; et un foncteur accessible $\mathcal{C} \to
\mathcal{C}'$ commutant aux limites donne de même $T(\mathcal{C}) \to
T(\mathcal{C}')$ avec adjoint à gauche. À prouver, enfin : pour
$\mathcal{C}$, $\mathcal{C}'$ accessibles, la catégorie
$\underline{\mathrm{Hom}}^{*}(\mathcal{C}^{\circ} \times \mathcal{C}'^{\circ},
\mathrm{Ens})$ des bifoncteurs commutant aux limites en chaque variable
est accessible, et c'est $T_{\mathcal{C}^{\circ}}T_{\mathcal{C}'^{\circ}}
(\mathcal{D}) \simeq T_{\mathcal{C}'^{\circ}}T_{\mathcal{C}^{\circ}}(\mathcal{D})$,
théorie « définie par $(\mathcal{C} \times \mathcal{C}')^{\circ}$ »,
universelle pour les couplages commutant aux limites en chaque
argument.\footnote{C'est le produit tensoriel de théories, et sa
commutativité ; la première conséquence est le théorème d'existence des
foncteurs « libres » entre catégories d'algèbres (Gabriel--Ulmer, et déjà
Freyd 1964 par le théorème du foncteur adjoint). Le point b) des
exemples — $T(\mathrm{Ens})$ est accessible pour $T$ définie par petites
limites — est l'énoncé qui fait de toute catégorie de modèles d'une
esquisse à limites une catégorie localement présentable.}

\pagerange{59}{61}
\subsection*{Le lemme sur $\mathrm{Ind}_{\pi}(\mathcal{C})$, et la caractérisation}

\textbf{Lemme.} \emph{Soient $\mathcal{C}$ petite, $\pi$ un cardinal, $E =
\mathrm{Ind}_{\pi}(\mathcal{C})$. Alors (i) $E$ satisfait $L_{\pi}$, et
l'inclusion $E \to \mathrm{Ind}(\mathcal{C})$ est $\pi$-accessible ; (ii) les
objets $\pi$-accessibles de $E$ sont les facteurs directs d'objets de
$\mathcal{C}$, de sorte que $E_{\pi}$ est l'enveloppe karoubienne de
$\mathcal{C}$, et $E_{\pi} = \mathcal{C}$ si $\mathcal{C}$ est karoubienne ;
(iii) $\mathcal{C}$ est strictement génératrice dans $E$ ; (iv) pour tout
$X$, $\mathcal{C}_{/X}$ est filtrante et grande devant $\pi$.}
\emph{Réciproquement, une $\mathcal{U}$-catégorie $E$ est équivalente à un
$\mathrm{Ind}_{\pi}(\mathcal{C})$ dès que : (i) $E$ satisfait $L_{\pi}$ ;
(ii) $E_{\pi}$ est petite et strictement génératrice ; (iii) pour tout $X$,
$E_{\pi/X}$ est filtrante et grande devant $\pi$ — et (iii) résulte de
(i), (ii) si $E$ est cocomplète.}

Puis les \emph{filtrations cardinales} : pour $\pi' \geqslant \pi$,
$\mathrm{Fil}^{\pi'}(E) = \mathrm{Ind}_{\pi}(\mathcal{C})_{\pi'}$ est formé
des ind-objets indexés par des $I$ grands devant $\pi$ et admettant une
partie cofinale de cardinal $\leqslant \pi'$ ; a) $\mathrm{Fil}^{\pi'}(E)
\subset E_{\pi'}$ ; b) si $\pi' = 2^{c}$ avec $c \geqslant \pi$, il y a
égalité ; c) si $E$ est cocomplète, les $\mathrm{Fil}^{\pi'}(E)$ forment
une filtration cardinale ; d) $E \simeq \mathrm{Ind}(E_{\pi'})$. La page 60
démontre l'axiome c) des filtrations cardinales pour les
$\mathrm{Fil}^{\pi'}$ (tout $X = \varinjlim_{I} X_{i}$ est la limite des
$X_{I'} = \varinjlim_{I'} X_{i}$ sur les parties filtrantes $I' \subset I$
de cardinal $\leqslant \pi'$, l'ensemble $J$ de ces parties étant filtrant
et grand devant $\pi'$), et la page 61, la démonstration du
\textbf{Corollaire} : \emph{une $\mathcal{U}$-catégorie cocomplète $E$ est
équivalente à un $\mathrm{Ind}_{\pi}(\mathcal{C})$, $\mathcal{C}$ petite, si
et seulement si elle admet une petite sous-catégorie strictement
génératrice formée d'objets accessibles. « Ce sont ces $E$ qu'il faudrait
appeler accessibles. »}\footnote{Ce sont, exactement, les catégories
\emph{localement présentables} de Gabriel et Ulmer, et l'énoncé est leur
théorème de caractérisation ; la démonstration de la page 61 (tout objet
$\pi$-accessible est facteur direct d'une $\varinjlim_{I} Y_{i}$ de
cardinal $\leqslant \pi$, puis $F^{\pi} = E_{\pi}$ par SGA 4 I 9.9) est
celle qui deviendra standard. La dernière phrase, sur le nom, dit que
l'exposé I de SGA 4 avait pris « accessible » pour une notion plus faible ;
la terminologie moderne a suivi le manuscrit sans le savoir, en réservant
« accessible » au cas non cocomplet (Makkai--Paré) et « localement
présentable » à celui-ci.}
\pagerange{62}{68}
\section*{Compléments : théories sur $\Delta$, et le passage aux préfaisceaux (pages 62 à 68)}

Les feuillets « Compléments » (titre de sa main, page 62) reprennent le
projet des pages 17 à 22 au niveau de généralité des pages 26 à 31 : $\Delta$
est une 2-catégorie de catégories à limites de type $d = (d', d'')$, une
\emph{théorie} sur $\Delta$ est un 2-foncteur $T : \Delta \to \mathrm{Cat}$,
elle est \emph{2-représentable} s'il existe $R \in \mathrm{Ob}\,\Delta$ avec
$T(C) \simeq \underline{\mathrm{Hom}}_{\Delta}(R, C)$, et $T_{R}$ désigne la
théorie représentée par $R$.\footnote{La définition précise — une théorie
est une catégorie 2-cofibrée sur $\Delta$ — est donnée page 83 ; nous
l'utilisons dès ici. Dans toute la suite « $R$ représente $T$ » signifie
cela, et « la catégorie classifiante » de $T$ désigne $R$, ce que la page
78 nomme ainsi et la page 111 « catégorie modulaire ».}

\pagerange{62}{63}
\subsection*{Caractériser les familles d'objets de base}

\emph{Problème (1).} Pour $R \in \mathrm{Ob}\,\Delta$ et $I \to R$ une famille
d'objets, caractériser la situation où $I$ est une famille d'objets de base
pour une théorie \emph{algébrique} $T_{R}$. La réponse constructive : $R$
est l'aboutissement d'une suite transfinie $(R_{\alpha})$ de catégories sur
$R$, pseudo-système inductif dans $\Delta$, avec $R_{0} = I$, $R_{\alpha+1}$
déduit de $R_{\alpha}$ soit en rendant inversible un petit ensemble de
flèches, soit en égalisant des couples de flèches, soit en ajoutant un
petit ensemble de flèches, et $R_{\alpha} = \varinjlim_{\beta < \alpha}
R_{\beta}$ aux ordinaux limites. Deux \emph{N.B.} précisent : (1) la
question ne se pose vraiment que si $(d', d'')$ n'est pas petit, mais même
alors la construction transfinie a un sens, pourvu qu'on admette dans la
définition récurrente des limites indexées par des ensembles bien
ordonnés ; (2) si $\widetilde{I}$ est la petite sous-catégorie pleine de $R$
engendrée par $I$ et par les sources et buts des flèches ajoutées, il est
plausible que $\underline{\mathrm{Hom}}_{\Delta}(R, C) \to
\underline{\mathrm{Hom}}_{\mathrm{Cat}}(\widetilde{I}, C) = T_{\widetilde{I}}
(C)$ soit pleinement fidèle, de sorte que $T_{R}$ se déduise de la théorie
algébrique $T_{\widetilde{I}}$ « en rajoutant seulement des axiomes » —
peut-être pas en une fois, mais par une construction transfinie. D'où deux
énoncés candidats pour « $T_{R}$ est algébrique » : 1) $R$ se déduit de
$\widetilde{J}$, $J$ petite, en rendant inversibles les flèches d'un
ensemble $\Sigma$ et en égalisant les couples d'un ensemble $\Sigma'$ ;
2) il existe un pseudo-système inductif transfini $(R_{\alpha})$ dans
$\Delta$ avec $R_{0} = \widetilde{J}$, chaque pas rendant inversibles et
égalisant de petits ensembles, et $R_{\alpha} = R$ pour $\alpha$ assez
grand dans l'univers.\footnote{Les trois opérations sont celles de toute
présentation d'une théorie : sortes (les objets de $I$), opérations (les
flèches ajoutées), axiomes (égalités et inversibilités, ces dernières
codant les conditions d'exactitude par le dédoublement de la page 31).
L'énoncé 1) « en une fois » sera établi pour $\Delta$ à limites de type
$d'$ petit aux pages 74 à 76 (le « Cor Main ») ; l'énoncé 2) est la
définition que les pages 81 à 87 adopteront.}

\pagerange{63}{64}
\subsection*{Conditions intrinsèques, conditions à la Schanuel, changement de $\Delta$}

\emph{Problème (2).} Conditionner les théories algébriques par des
propriétés intrinsèques plutôt que par la constructibilité. La seule
condition nécessaire visible : $T$ commute aux 2-limites projectives de
$\Delta$, en particulier transforme les foncteurs pleinement fidèles
(resp. fidèles) en foncteurs pleinement fidèles (resp. fidèles), et
$T(\underline{\mathrm{Hom}}_{\mathrm{Cat}}(J, C)) \simeq \underline{\mathrm{Hom}}_{
\mathrm{Cat}}(J, T(C))$ ; « probablement pas suffisante ».\footnote{La page
écrit le membre de gauche sans $T$ ; la formule complète est celle de la
page 65, encadrée. La question de caractériser intrinsèquement les
2-foncteurs représentables est celle que la théorie des esquisses résout
par la notion d'\emph{accessibilité} du foncteur ; le manuscrit ne dispose
que de la commutation aux 2-limites, qui ne suffit pas en effet.}

\emph{Problème (3).} Conditions d'exactitude \emph{généralisées}, « comme
suggérées par Schanuel » : on se donne $\rho \to \sigma$ dans $R$ — un
morphisme entre petites catégories de $\Delta$ — et l'on demande sur $C$
que $\underline{\mathrm{Hom}}_{\Delta}(\sigma, C) \to \underline{\mathrm{Hom}}
_{\Delta}(\rho, C)$ soit une équivalence, ou pleinement fidèle, ou fidèle.
Cela inclut l'égalisation de flèches, et les remarques sur l'existence
des limites restent valables.\footnote{C'est la forme d) de la typologie
des pages 10 à 14, désormais attribuée : Schanuel, dont l'idée d'imposer
la bijectivité de $\mathrm{Hom}(\sigma, -) \to \mathrm{Hom}(\rho, -)$ est
la notion d'\emph{orthogonalité} de la théorie ultérieure (Freyd--Kelly
1972). C'est la seule attribution d'idée du dossier, avec Lazard et
Lawvere.}

\emph{Problème (4).} Changement de $\Delta \subset \Delta'$. Le diagramme de
la page 64 place les théories $\mathrm{Th}_{\Delta}$, les théories
algébriques $\mathrm{Thal}_{\Delta}$ et les catégories $\Delta_{\mathrm{al}}$
qui les représentent, avec la restriction $\mathrm{Th}_{\Delta'} \to
\mathrm{Th}_{\Delta}$ en haut et la $\Delta'$-enveloppe $\rho_{\Delta',
\Delta} : \Delta \to \Delta'$ en bas, les verticales étant les
2-équivalences $\mathrm{Thal} \simeq \Delta_{\mathrm{al}}$ :

\begin{tikzcd}[column sep=small, row sep=small]
\mathrm{Th}_{\Delta'} \arrow[rr, "\text{restriction}"] \arrow[d] & & \mathrm{Th}_{\Delta} \arrow[d] \\
\mathrm{Thal}_{\Delta'} \arrow[rr, "\text{restr.}"] \arrow[d, no head, "\wr" description] & & \mathrm{Thal}_{\Delta} \arrow[d, no head, "\wr" description] \\
\Delta'_{\mathrm{al}} \arrow[rr] & & \Delta_{\mathrm{al}}
\end{tikzcd}

La flèche du bas est l'enveloppe $\rho_{\Delta',\Delta} : \Delta' \to
\Delta$, et elle va \emph{dans le même sens} que la restriction : si $T'$
est représentée par $R' \in \Delta'$, sa restriction à $\Delta$ est
représentée par la $\Delta$-enveloppe de $R'$, puisque
$\underline{\mathrm{Hom}}_{\Delta'}(R', C) \simeq \underline{\mathrm{Hom}}_{\Delta}
(\rho(R'), C)$ pour $C \in \mathrm{Ob}\,\Delta$. Le cube dessiné à droite
est le même diagramme avec les inclusions $\mathrm{Thal} \subset \mathrm{Th}$
et $\Delta_{\mathrm{al}} \subset \Delta$ en profondeur.\footnote{La page 64
ne comporte que le diagramme, son cube, et la ligne « 2-foncteur d'où
$\Delta'$-enveloppe, défini ainsi : $\Delta' \to \Delta$,
$\rho_{\Delta',\Delta}$ ». Ici $\Delta \subset \Delta'$ signifie que
$\Delta$ exige plus de limites que $\Delta'$, comme $\Delta_{1} \subset
\Delta_{0}$ page 74. Le contenu de la question — quand la restriction
est-elle pleinement fidèle, quelle est son image — est donné aux pages 77
et 78 pour $\Delta_{1} \subset \Delta_{0}$ et repris page 102
($\lambda$-variable).}

\pagerange{65}{65}
\subsection*{Limites dans $T(C)$, et le carré 2-cartésien}

\textbf{Proposition.} \emph{Soit $T = T_{R}$ 2-représentable, et $\delta$ un
type de diagrammes tel que a) les limites projectives de type $\delta$
existent dans $C$, b) elles y commutent aux limites inductives de type
$d''$. Alors les limites projectives de type $\delta$ existent dans $T(C)
\simeq \underline{\mathrm{Hom}}_{\Delta}(R, C)$ et se calculent « argument
par argument » : les foncteurs construction $r_{c} : T(C) \to C$, $\xi
\mapsto \xi(c)$, y commutent. Un morphisme de théories $T \to T'$ induit
$T(C) \to T'(C)$ commutant à ces limites, et un $\Delta$-foncteur $C \to C'$
vers une $C'$ satisfaisant les mêmes hypothèses induit $T(C) \to T(C')$ de
même.}

La preuve est celle qu'on attend : une limite de type $\delta$ de
foncteurs $R \to C$ se calcule dans $\underline{\mathrm{Hom}}(R, C)$
argument par argument, et le foncteur limite commute aux limites de type
$d'$ (toujours) et aux limites inductives de type $d''$ (par b)), donc reste
un $\Delta$-foncteur. Conséquences : pour $J$ petite,
$\underline{\mathrm{Hom}}_{\mathrm{Cat}}(J, C)$ est dans $\Delta$ et
\[
\boxed{\ T_{R}\bigl(\underline{\mathrm{Hom}}_{\mathrm{Cat}}(J, C)\bigr)
\simeq \underline{\mathrm{Hom}}_{\mathrm{Cat}}\bigl(J, T_{R}(C)\bigr)\ } ;
\]
en particulier, si $(\mathrm{Ens}) \in \mathrm{Ob}\,\Delta$, avec $J^{\circ}$ :
\[
T_{R}(\widehat{J}) \simeq \underline{\mathrm{Hom}}(J^{\circ}, T_{R}(\mathrm{Ens})).
\]
Si $C \to C'$ est un $\Delta$-foncteur pleinement fidèle (resp. fidèle),
$T(C) \to T(C')$ l'est aussi ; et si $T \to T'$ est associé à $R' \to R$
« $\Delta$-générateur » — la sous-catégorie pleine de $R$ dans $\Delta$
engendrée par l'image de $R'$ est $R$ —, alors le carré

\begin{tikzcd}
T(C) \arrow[r, hook] \arrow[d] & T(C') \arrow[d] \\
T'(C) \arrow[r, hook] & T'(C')
\end{tikzcd}

est 2-cartésien : une $T$-structure sur $C'$ dont la $T'$-structure
sous-jacente provient de $C$ provient de $C$.\footnote{Le cas
particulier annoncé page 66 — $T' = \tau^{I}$, la structure vide sur $I$
objets — n'est pas écrit ; c'est le carré (8.2.3) de la page 104, et la
raison en est donnée là : un $\Delta$-foncteur $R \to C'$ dont les valeurs
sur une famille $\Delta$-génératrice sont dans $C$ est à valeurs dans $C$,
puisque $C \subset C'$ est stable par les limites en jeu. La proposition
sur les limites dans $T(C)$ est la remarque, aujourd'hui banale, que la
catégorie des modèles d'une théorie à limites dans $C$ a les limites que
$C$ a, calculées sur les objets sous-jacents ; la condition b) est ce qui
la rend vraie pour une théorie à limites \emph{inductives} aussi.}

\pagerange{67}{68}
\subsection*{Théories définies par limites projectives : $R^{\circ} \subset S \subset \widehat{R^{\circ}}$}

Le point (4), $d'' = \emptyset$ — les structures définies seulement par des
limites projectives —, réunit les « faits particuliers » qui gouverneront
la suite du dossier. Dans ce cas :

$1^{\circ}$) les $T(C)$ ont toutes les limites projectives que $C$ a, et
tous les foncteurs $T(C) \to T(C')$, $T(C) \to T'(C)$, $T(C) \to C$ y
commutent ;

$2^{\circ}$) $C \to \widehat{C}$ est un $\Delta$-foncteur, donc, si
$(\mathrm{Ens}) \in \mathrm{Ob}\,\Delta$ et $T = T_{R}$,
\[
T_{R}(C) \longrightarrow T_{R}(\widehat{C}) \simeq
\underline{\mathrm{Hom}}_{\mathrm{Cat}}(C^{\circ}, S), \qquad S = T_{R}(\mathrm{Ens}),
\]
est pleinement fidèle et identifie $T_{R}(C)$ à la sous-catégorie pleine
des $\varphi : C^{\circ} \to S$ tels que, pour toute construction $r \in R$
et le foncteur $\tilde{r} : S \to (\mathrm{Ens})$ qu'elle définit, $\tilde{r}
\circ \varphi : C^{\circ} \to (\mathrm{Ens})$ soit représentable ; et il
suffit de l'exiger pour une famille $(r_{i})$ qui $\Delta$-engendre $R$,
i.e. le carré

\begin{tikzcd}[column sep=small]
T(C) \arrow[r] \arrow[d] & \underline{\mathrm{Hom}}_{\mathrm{Cat}}(C^{\circ}, S) \arrow[d, "\varphi \mapsto (\tilde{r}_{i} \circ \varphi)"] \\
C^{I} \arrow[r] & \underline{\mathrm{Hom}}_{\mathrm{Cat}}(C^{\circ}, (\mathrm{Ens}))^{I}
\end{tikzcd}

est 2-cartésien. \emph{Donc $T$ est connu, à équivalence près, quand on
connaît $S$ et les $\tilde{r}_{i} : S \to (\mathrm{Ens})$}, ou, de façon plus
intrinsèque, l'image de $R$ dans $\underline{\mathrm{Hom}}(S, \mathrm{Ens}) =
\widehat{S^{\circ}}$.\footnote{C'est le théorème par lequel une théorie à
limites projectives est déterminée par sa catégorie de modèles
ensemblistes $S$ : un modèle dans $C$ est un foncteur $C^{\circ} \to S$
« représentable en chaque construction ». La version pour les limites
finies est la dualité de Gabriel--Ulmer, et c'est la forme que la page 72
lui donnera ; celle-ci, pour $d'$ quelconque, ne demande que le
plongement de Yoneda et l'échange des limites.}

$3^{\circ}$) $S \simeq \underline{\mathrm{Hom}}_{\Delta}(R, \mathrm{Ens})
\subset \widehat{R^{\circ}}$, et cette sous-catégorie \emph{contient}
$R^{\circ}$, parce que les foncteurs représentables covariants $\mathrm{Hom}
(r, -)$ commutent aux limites projectives :
\[
\boxed{\ R^{\circ} \xrightarrow{\ \partial\ } S\ }, \qquad
\tilde{r} = \mathrm{Hom}_{S}(\partial r, -) : S \to (\mathrm{Ens}),
\]
le foncteur associé à la construction $r$ étant celui que $\partial r$
représente \emph{dans} $S$. Ainsi $R$ est connu quand on connaît $S$ et
l'image de $R$ dans $\widehat{S^{\circ}}$, car $R \to \widehat{S^{\circ}}$
est le composé de deux foncteurs pleinement fidèles $R \xrightarrow{
\partial^{\circ}} S^{\circ} \to \widehat{S^{\circ}}$. La sous-catégorie
pleine $\Sigma \approx R^{\circ}$ de $S$ se décrit intrinsèquement : elle
est formée des $\varphi \in \mathrm{Ob}\,S$ tels que, pour tout $C \in
\mathrm{Ob}\,\Delta$ et tout $\xi \in T(C)$ définissant $\tilde\xi :
C^{\circ} \to S$, le foncteur $x \mapsto \mathrm{Hom}_{S}(\varphi,
\tilde\xi(x))$ soit représentable ; $\Sigma$ est stable dans $S$ par
limites \emph{inductives} de type $d'$, et il y a une équivalence
explicite $\Sigma^{\circ} \simeq R$.\footnote{Que $R^{\circ} \subset \Sigma$
est clair : pour $\varphi = \partial r$, $\mathrm{Hom}_{S}(\partial r,
\tilde\xi(x)) = \tilde\xi(x)(r) = \mathrm{Hom}_{C}(x, \xi(r))$ est
représenté par $\xi(r)$. L'inclusion inverse est affirmée sans
démonstration, et le point $4^{\circ}$), qui l'abordait, s'arrête page 68
après trois lignes barrées. Pour $d'$ = les diagrammes finis, la page 77
la démontrera par une autre voie : $\Sigma$ est alors la sous-catégorie des
objets de présentation finie de $S$. La page 106 reprendra la description
mot pour mot. La stabilité de $\Sigma$ par limites inductives de type $d'$
vient de ce que $R$ a les limites projectives de type $d'$ et que $\partial$
les transforme en limites inductives de $S$ — ce qui se vérifie sur les
Hom : $\mathrm{Hom}_{S}(\partial(\varprojlim r_{i}), F) = F(\varprojlim
r_{i}) = \varprojlim F(r_{i})$ puisque $F$ commute aux limites.}

\pagerange{69}{80}
\section*{Exemples : les théories définies par des limites (pages 69 à 80)}

\pagerange{69}{69}
\subsection*{Trois cas triviaux}

(1) $d' = d'' = \emptyset$, $\Delta = \mathrm{Cat}$ : les problèmes 1 et 2
ont des solutions évidentes, les théories algébriques sont exactement les
théories représentées par des \emph{petites} catégories, et les $I \to R$
basiques sont les applications surjectives sur les objets. (2) Avec une
condition d'exactitude, par exemple $\Delta$ = les groupoïdes : le
problème 1 est le groupoïde fondamental d'une catégorie (rendre toutes les
flèches inversibles), le problème 2 est trivial, et le reste comme en (1) ;
la condition s'écrit avec $\rho = [a \xrightarrow{u} b]$, « on rend $u$
inversible ». (3) $\Delta$ = les catégories (pré)ordonnées : $\rho = [a
\rightrightarrows b]$, « on rend $u$ égal à $v$ ».\footnote{Les deux
exemples montrent que « condition d'exactitude » couvre, dans le
vocabulaire du dossier, toute condition de la forme « $\mathrm{Hom}(\sigma,
C) \to \mathrm{Hom}(\rho, C)$ est une équivalence » ; les groupoïdes et les
préordres sont les deux cas extrêmes.}

\pagerange{70}{73}
\subsection*{Toutes les limites projectives : $\Delta_{1}$}

Le cas $d'$ = tous les petits diagrammes, $d'' = \emptyset$, sans condition
d'exactitude ; notons $\Delta_{1}$ cette 2-catégorie.\footnote{La page 70
écrit « $d''$ = tous les petits diagrammes, $d'' = \emptyset$ », lapsus
pour $d'$ ; la page 74 nomme $\Delta_{1}$ cette 2-catégorie, et nous
adoptons ce nom dès ici. Les objets de $\Delta_{1}$ sont les catégories
complètes, ses flèches les foncteurs continus.}

\emph{Problème 1.} Pour $I$ petite et $C \in \mathrm{Ob}\,\Delta_{1}$,
\[
\underline{\mathrm{Hom}}(I, C) \simeq \underline{\mathrm{Hom}}(I^{\circ},
C^{\circ})^{\circ} \simeq \underline{\mathrm{Hom}}_{\Delta_{1}^{\circ}}
(\widehat{I^{\circ}}, C^{\circ})^{\circ} \simeq \underline{\mathrm{Hom}}_{
\Delta_{1}}(\widehat{I^{\circ}}{}^{\circ}, C),
\]
donc la $\Delta_{1}$-enveloppe de $I$ est $\widetilde{I} =
\widehat{I^{\circ}}{}^{\circ}$, l'opposé d'un topos de préfaisceaux, et $I
\to \widetilde{I}$ fait des objets de $I$ une famille
\emph{cogénératrice}.\footnote{C'est la propriété universelle du topos de
préfaisceaux $\widehat{I^{\circ}}$ comme complétion libre de $I^{\circ}$ par
les petites limites inductives, dualisée. L'enveloppe n'est pas petite, et
c'est la raison de la restriction « $d$ petit » du théorème des pages 3
et 27.}

\emph{Problème 2.} $C \in \mathrm{Ob}\,\Delta_{1}$, $\Sigma \subset \mathrm{Fl}
\,C$ petit : rendre inversibles les flèches de $\Sigma$ dans $\Delta_{1}$.
On peut supposer que $\Sigma$ contient les isomorphismes, est stable par
composition, et vérifie « $vu \in \Sigma$, $v \in \Sigma \Rightarrow u \in
\Sigma$ » ; alors $\Sigma$ admet un calcul de fractions, les limites
projectives de $C$ existent dans $C\Sigma^{-1}$ et $C \to C\Sigma^{-1}$ y
commute, pourvu qu'on suppose aussi $\Sigma$ stable par produits infinis de
morphismes, ce qui règle les produits ; reste à vérifier que $C\Sigma^{-1}$
est une $\mathcal{U}$-catégorie. Si $I$ est cogénératrice par
monomorphismes stricts dans $C$, elle l'est encore dans $C\Sigma^{-1}$ ; et
la limite inductive $R$ dans $\Delta_{1}$ d'un pseudo-système filtrant
$(R_{\alpha})$ à cogénérateurs $I \to R_{\alpha}$ admet encore $I$ pour
cogénérateur, parce que $R$ s'obtient de $\varinjlim^{\mathrm{Cat}}
R_{\alpha}$ en rendant inversibles des flèches — « mais attention : le fait
d'être strictement cogénérateur n'est pas transitif ».

\textbf{Théorème.} \emph{Les théories algébriques pour $\Delta_{1}$ sont
celles qui sont 2-représentables par des $R \in \mathrm{Ob}\,\Delta_{1}$
admettant une petite famille strictement cogénératrice. Pour que $\rho : I
\to R$ soit basique, il faut et suffit que ses images soient strictement
$\Delta_{1}$-cogénératrices.}

La nécessité est claire par ce qui précède ; pour la suffisance, soit
$I'$ la sous-catégorie pleine de $R$ d'objets $\rho(I)$ : $T_{\widetilde{I}'}$
est algébrique de base $I$, et $\widetilde{I}' \to R$ est surjectif sur
les objets (essentiellement) et sur les Hom, de sorte que $R$ se déduit
de $\widetilde{I}'$ en rendant inversibles certaines flèches. Les
\emph{Remarques} : (1) la construction se fait en quatre étapes, objets de
base, flèches, égalités, inversibilités, et 3\textsuperscript{o}--4\textsuperscript{o}
se condensent en une catégorie de fractions ; (2) les quotients $R'$ de $R$
dans $\Delta_{1}$ par inversion de flèches correspondent bijectivement aux
ensembles $\Sigma$ de flèches stables par composition, par changement de
base, contenant les isomorphismes, et stables par produits quelconques ;
(3) si $R$ a une petite famille cogénératrice, alors
\[
\underline{\mathrm{Hom}}_{\Delta_{1}}(R, \mathrm{Ens}) \simeq R^{\circ} :
\]
les foncteurs $R \to (\mathrm{Ens})$ commutant aux limites sont les
représentables ; donc, pour $T = T_{R}$ et $S = T(\mathrm{Ens})$, on a $R =
S^{\circ}$.\footnote{C'est le théorème du foncteur adjoint spécial, qui
demande, outre la complétude et la petite famille cogénératrice, que $R$
soit \emph{bien puissante} (que les sous-objets d'un objet forment un
ensemble) — hypothèse que la page ne mentionne pas et sans laquelle
l'énoncé est faux. Elle est satisfaite dans tous les cas que le dossier
considère : $R = S^{\circ}$ avec $S$ localement présentable ou $S$ un
topos. La remarque (2) est correcte telle quelle, et c'est la description
des localisations d'une catégorie complète par ses « classes de flèches
saturées ».}

\textbf{Corollaire.} \emph{Le 2-foncteur $T \mapsto T(\mathrm{Ens})$ est une
2-équivalence entre les théories définissables par limites projectives
quelconques, sans axiome d'exactitude, et les $\mathcal{U}$-catégories $S$
possédant les petites limites projectives et une petite famille d'objets
génératrice par épimorphismes stricts, i.e. telle que $S \to
\widehat{I^{\circ}}$ soit conservatif. La donnée d'une base de $T$ équivaut
à celle d'une famille cogénératrice de $S$. Pour tout $C \in \mathrm{Ob}\,
\Delta_{1}$,}
\[
T(C) \simeq \underline{\mathrm{Hom}}_{\Delta_{1}}(S^{\circ}, C),
\]
\emph{les foncteurs $S^{\circ} \to C$ continus s'interprétant comme les
foncteurs $S^{\circ} \times C^{\circ} \to (\mathrm{Ens})$ représentables en
chaque argument, ou comme les foncteurs $C^{\circ} \to S$ tels que pour tout
$Y \in S$ le composé avec $\mathrm{Hom}_{S}(Y, -)$ soit
représentable.}\footnote{La page écrit « $S$ avec $\varprojlim$ quelc. » ;
mais $S = R^{\circ}$ est l'opposée d'une catégorie \emph{complète}, donc
$S$ est \emph{cocomplète}, et c'est la remarque (4) de la page 72 qui le
dit aussitôt (« $R$ a aussi des $\varinjlim$ quelconques »). Les deux se
concilient : $S = \underline{\mathrm{Hom}}_{\Delta_{1}}(R, \mathrm{Ens})$ a
les limites projectives calculées argument par argument, et les limites
inductives de $R$. « Génératrice par épimorphismes stricts » est la
condition duale de « cogénératrice par monomorphismes stricts » du
théorème ; avec la bonne puissance implicite, ce corollaire est le
dictionnaire de la page 90 pour $\lambda$ = toutes les limites. Un
exemple non trivial de telle $S$ : les espaces compacts, opposée des
$C^{*}$-algèbres commutatives, qui n'est pas localement présentable —
signe que ce cas $\Delta_{1}$ déborde le cadre des esquisses petites.}

(5) Un foncteur $\varphi : R \to R'$ entre catégories complètes dont la
source a une petite famille cogénératrice est un $\Delta_{1}$-foncteur
si et seulement s'il admet un adjoint à gauche $\psi$ ; il est un
« plongement », i.e. correspond à une sous-théorie pleine $T' = T_{R'}
\subset T = T_{R}$, si et seulement si $\psi$ est pleinement fidèle. En
termes de $S = R^{\circ}$ et $S' = R'^{\circ}$ : les sous-théories pleines
de $T$ correspondent aux sous-catégories pleines $S' \subset S$, complètes,
dont le foncteur d'inclusion $i = \varphi^{\circ}$ admet un adjoint à
gauche $j = \psi^{\circ}$,
\[
\mathrm{Hom}_{S}(x, i(y')) \simeq \mathrm{Hom}_{S'}(jx, y') ;
\]
ce sont les sous-catégories \emph{réflexives}, et elles sont
nécessairement stables par limites projectives dans $S$.\footnote{La page
73 écrit « admettant des $\varprojlim$ quelconques (mais pas
nécessairement stables par lesdites) » ; c'est inexact : le foncteur
d'inclusion $i$, adjoint à droite, commute à toutes les limites
projectives, de sorte que la limite dans $S'$ d'un diagramme de $S'$ est
aussi sa limite dans $S$. Une sous-catégorie pleine réflexive est toujours
stable par limites. Nous avons corrigé. La contravariance est celle du
foncteur $R \mapsto T_{R}$ : à $\varphi : R \to R'$ correspond $T_{R'} \to
T_{R}$, et $\varphi^{\circ} : S \to S'$ est le foncteur « structure libre »
tandis que $\psi^{\circ} : S' \to S$ est le foncteur d'oubli.} Le passage
$S' \to S$ « structure libre définie par une $S$-structure » n'a pas
d'analogue pour $C$ quelconque : en général on ne peut définir de foncteur
$T(C) \to T'(C)$ redonnant $j$ pour $C = \mathrm{Ens}$, sauf si $C$ est un
topos et si $T$, $T'$ sont définies par limites finies. Enfin
(\emph{N.B.}, page 73) : \emph{la 2-catégorie des théories algébriques pour
$\Delta_{1}$ est 2-équivalente à celle des catégories $S$ possédant les
petites limites projectives et une petite sous-catégorie strictement
génératrice, avec pour morphismes les foncteurs commutant aux limites
projectives}, un morphisme $T' \to T$ correspondant au foncteur d'oubli
$T'(\mathrm{Ens}) \to T(\mathrm{Ens})$.\footnote{« Mais, en renversant les
flèches pour les morphismes ? », demande la page. La réponse est non : à $T' \to T$
correspond $S' \to S$ dans le même sens, le foncteur d'oubli. La
2-équivalence est covariante sur ces morphismes et c'est
$R \mapsto T_{R}$ qui est contravariant.}

\pagerange{74}{76}
\subsection*{Limites de type $d'$ : enveloppes et théorème}

Le cas $d'$ quelconque, $d'' = \emptyset$, sans condition d'exactitude ;
$d'_{0}$ = tous les diagrammes finis, $d'_{1}$ = tous les petits diagrammes,
et $\Delta \supset \Delta_{0} \supset \Delta_{1}$ les 2-catégories
correspondantes.\footnote{Ordre des inclusions : plus il y a de limites
exigées, moins il y a de catégories. $\Delta_{0}$ est ainsi la 2-catégorie
des catégories à limites finies de la page 17.}

\emph{Problème 1.} La $\Delta$-enveloppe $\rho_{\Delta}(I)$ d'une petite
catégorie $I$ : soit $\overline{I} = \widehat{I^{\circ}}{}^{\circ}$ sa
$\Delta_{1}$-enveloppe, dans laquelle $I$ se plonge pleinement fidèlement,
et $\widetilde{I}$ la plus petite sous-catégorie strictement pleine de
$\overline{I}$ contenant $I$ et stable par limites projectives de type
$d'$. Alors $\widetilde{I}$ est la $\Delta$-enveloppe. Il faut pour cela
le \emph{problème 1'}, ou problème auxiliaire : pour $C \in \mathrm{Ob}\,
\Delta$, le foncteur canonique de $C$ dans sa $\Delta \to \Delta_{1}$
enveloppe $\rho(C)$ est pleinement fidèle. Or
\[
\rho(C) = \underline{\mathrm{Hom}}_{\Delta}(C, \mathrm{Ens})^{\circ} =
T_{C}(\mathrm{Ens})^{\circ},
\]
le foncteur $C \to \rho(C)$ associant à tout objet le foncteur covariant
qu'il représente — d'où la pleine fidélité, par Yoneda — et l'on vérifie
la propriété universelle par $\underline{\mathrm{Hom}}_{\Delta_{1}}(\rho(C),
D) \simeq \underline{\mathrm{Hom}}_{\Delta}(C, D)$. Alors les foncteurs
$I \to C$ correspondent aux $\Delta_{1}$-foncteurs $\overline{\psi} :
\overline{I} \to \rho(C)$ envoyant $I$ dans l'image essentielle de $C$, qui
envoient $\widetilde{I}$ dans $C$, d'où la restriction
$\underline{\mathrm{Hom}}_{\Delta}(\widetilde{I}, C) \to \underline{\mathrm{Hom}}
(I, C)$ et son inverse.\footnote{La démonstration est celle des
« complétions libres par une classe de limites » : plonger dans la
complétion par toutes les limites et prendre la clôture. Elle est correcte
sous la réserve, que la page fait, que $\widetilde{I}$ soit stable par les
limites de type $d'$ ; la page 74 signale que si $d'$ n'est pas stable par
sous-diagrammes il faut prendre la clôture « transfiniment au besoin ».
L'identification $\rho(C) = \underline{\mathrm{Hom}}_{\Delta}(C,
\mathrm{Ens})^{\circ}$ demande, comme en (3) ci-dessus, que $\rho(C)$ soit
une $\mathcal{U}$-catégorie, ce qui est le cas pour $C$ petite.}

\emph{Problème 2.} Pour $C \in \Delta$ et $\Sigma$ un ensemble de flèches
(et de couples de flèches à égaliser), soit $\rho(C, \Sigma)$ la
$\Delta_{1}$-catégorie de fractions de $\rho(C)$ définie par $\Sigma$, et
$C'$ la plus petite $\Delta$-sous-catégorie strictement pleine de $\rho(C,
\Sigma)$ contenant l'image de $C$, i.e. contenant $C\Sigma^{-1}$ ; alors
$\underline{\mathrm{Hom}}_{\Delta,\Sigma}(C, D) \simeq \underline{\mathrm{Hom}}
_{\Delta}(C', D)$, la pleine fidélité de la restriction résultant de ce
que $C'$ se déduit de $C\Sigma^{-1}$ en ajoutant transfiniment des limites
de type $d'$. \emph{N.B.} : dans les deux cas la catégorie résultante se
déduit de $I$ (resp. $C\Sigma^{-1}$) en rajoutant des limites de type
$d'$, d'où :

\textbf{Théorème.} \emph{Les théories algébriques de type $\Delta$ sont
celles qui sont 2-représentées par des $R \in \mathrm{Ob}\,\Delta$ ayant une
petite sous-catégorie $d'$-génératrice, i.e. telle que $R$ soit la réunion
d'une suite transfinie de sous-catégories strictement pleines déduites
l'une de l'autre en rajoutant des limites de type $d'$. Un foncteur $I \to
R$, $I$ discret, est basique si et seulement s'il est $d'$-cogénérateur.}

La page 76 ramène le second énoncé au premier comme pour $\Delta_{1}$ :
$I' \subset R$ la sous-catégorie pleine d'objets $I$, $T_{R} \to
T_{\widetilde{I}'}$ pleinement fidèle, et si $d'$ satisfait aux conditions
de stabilité on trouve un $\Delta$-foncteur pleinement fidèle et
essentiellement surjectif $\widetilde{I}'\Sigma^{-1} \to R$, d'où le
\textbf{Corollaire principal} : \emph{on peut construire les théories
algébriques pour $\Delta$ en trois pas, comme pour $\Delta_{1}$ : une
catégorie $I'$ ayant $I$ pour ensemble d'objets (pas $1^{\circ}$ et
$2^{\circ}$), et une catégorie de fractions de $\widetilde{I}'$ (pas
$3^{\circ}$)}, la conclusion valant par exemple si $d'$ est formé de
diagrammes discrets stable par sous-diagrammes, ou de tous les diagrammes
finis ; pour un $d'$ général « il semble qu'il faille une construction
récurrente transfinie ». \emph{Question} : si $d'$ est petit, une $R$ à
petite sous-catégorie $d'$-cogénératrice est petite.\footnote{C'est
l'énoncé « toute théorie à limites de type $d'$ se présente en un pas par
une esquisse » pour $d'$ petit, i.e. l'équivalence entre théories
constructibles et esquisses petites, et la réponse à la question de la
page 62. La \emph{Question} finale répond à la note de la page 3 sur la
petitesse : si $d'$ est petit, la clôture transfinie d'une petite
catégorie sous les limites de type $d'$ est petite, par le même argument
d'ordinal $\omega$.}

\pagerange{77}{78}
\subsection*{Limites finies : $R$, $T$, $S$, $\Sigma$ se déterminent mutuellement}

Le cas $d'$ = les diagrammes finis, $d'' = \emptyset$, sans condition
d'exactitude, est celui de la page 17 et il est ici résolu complètement.
Comme $d'$ est petit, les $R \in \mathrm{Ob}\,\Delta_{0}$ qui représentent
des théories algébriques sont les catégories essentiellement petites à
limites finies. Étant donnés
\begin{quote}
$T$ une théorie algébrique pour $\Delta_{0}$ ; \\
$R \in \mathrm{Ob}\,\Delta_{0}$ qui la représente ; \\
$S = T(\mathrm{Ens}) = \underline{\mathrm{Hom}}_{\mathrm{lex}}(R, \mathrm{Ens})$ ; \\
$\Sigma \subset S$ la sous-catégorie pleine des structures représentables,
\end{quote}
tout se détermine mutuellement :
\begin{align*}
R \ \text{donne}\quad & T = T_{R} : C \mapsto \underline{\mathrm{Hom}}_{\mathrm{lex}}(R, C),
\quad S = \underline{\mathrm{Hom}}_{\mathrm{lex}}(R, \mathrm{Ens}),
\quad \Sigma \simeq R^{\circ} \subset S ; \\
\Sigma \ \text{donne}\quad & R = \Sigma^{\circ},
\quad T : C \mapsto \underline{\mathrm{Hom}}_{\mathrm{lex}}(\Sigma^{\circ}, C),
\quad S = \underline{\mathrm{Hom}}_{\mathrm{lex}}(\Sigma^{\circ}, \mathrm{Ens}) \supset \Sigma ;
\end{align*}
et $S$ se déduit de $\Sigma$ par le

\textbf{Lemme.} \emph{Le foncteur d'inclusion $\Sigma \to S$ induit une
équivalence $\mathrm{Ind}(\Sigma) \xrightarrow{\ \approx\ } S$.}

En effet $\Sigma \to S$ est universel pour les foncteurs exacts à gauche
$\Sigma \to D$ vers une $D$ à limites inductives filtrantes ; or $\mathrm{Ind}
(\Sigma)$ a les limites filtrantes et les limites projectives finies, $\Sigma
\to \mathrm{Ind}(\Sigma)$ y commute (SGA 4 I 8.9.5), et les foncteurs
exacts $\Sigma \to D$ correspondent aux foncteurs $\mathrm{Ind}(\Sigma) \to
D$ commutant aux limites filtrantes et aux limites finies — la nécessité
est triviale, et pour la suffisance $\mathrm{Ind}(\Sigma) \to \mathrm{Ind}(D)
\to D$ est exact (SGA 4 I 8.9.8).

\textbf{Corollaire.} \emph{Pour $X \in S$, on a $X \in \Sigma$ si et seulement
si $\mathrm{Hom}_{S}(X, -)$ commute aux limites inductives filtrantes.}

Nécessaire par construction des Hom dans $\mathrm{Ind}$ ; suffisant parce
que $X = \varinjlim X_{i}$ avec $X_{i} \in \Sigma$, filtrante, et que
l'identité de $X$ se factorise alors par un $X_{i}$, de sorte que $X$ est
un facteur direct de $X_{i}$, donc dans $\Sigma$, qui est stable par
noyaux.\footnote{« Donc stable (c'est un noyau) » : un facteur direct est
le coégalisateur de l'identité et de l'idempotent, et $\Sigma = R^{\circ}$
a les limites inductives finies parce que $R$ a les limites finies. Une
catégorie à coégalisateurs est karoubienne ; il n'y a donc pas ici
d'écart entre $R^{\circ}$ et son enveloppe karoubienne, contrairement à ce
qui se produirait pour une esquisse à produits finis seulement (page 78,
\emph{Question}). L'ensemble lemme + corollaire est la \emph{dualité de
Gabriel--Ulmer} (1971) : les petites catégories à limites finies
correspondent, par $R \mapsto \underline{\mathrm{Hom}}_{\mathrm{lex}}(R,
\mathrm{Ens})$ et $S \mapsto (S_{\mathrm{pf}})^{\circ}$, aux catégories
localement finiment présentables. Les renvois marginaux — SGA 4 I 8.7.5,
8.9.5, 8.7.3, 8.9.8, 8.7.8 — disent d'où le manuscrit tire chaque pas.}

\emph{Donc $\Sigma$ se reconstitue à partir de $S$ comme la sous-catégorie
pleine des objets « de présentation finie »} :
\begin{align*}
S \ \text{donne}\quad & \Sigma = S_{\mathrm{pf}} \subset S, \quad R = \Sigma^{\circ},
\quad T : C \mapsto \underline{\mathrm{Hom}}_{\mathrm{lex}}(S_{\mathrm{pf}}{}^{\circ}, C) ; \\
T \ \text{donne}\quad & S = T(\mathrm{Ens}), \quad \Sigma = T(\mathrm{Ens})_{\mathrm{pf}},
\quad R = \Sigma^{\circ}.
\end{align*}

\textbf{Corollaire.} \emph{Pour qu'une $\Delta_{1}$-théorie $T$ (par
limites projectives quelconques) provienne d'une théorie définissable par
limites finies, il faut et suffit que $S = T(\mathrm{Ens})$ soit strictement
engendrée par la sous-catégorie pleine de ses objets de présentation
finie, ou encore que $S \to \widehat{\Sigma}$ soit pleinement fidèle ; alors
$\Sigma^{\circ}$ est la catégorie classifiante de la théorie sur $\Delta_{0}$
qui donne naissance à $T$.}

\textbf{Corollaire.} \emph{Le 2-foncteur de restriction des théories
algébriques sur $\Delta_{0}$ aux théories sur $\Delta_{1}$ est pleinement
fidèle, et son image essentielle est formée des morphismes $T_{1} \to
T'_{1}$ tels que $S = T(\mathrm{Ens}) \to S' = T'(\mathrm{Ens})$ transforme
objets de présentation finie en objets de présentation finie — « on
pourrait les appeler les morphismes de présentation finie, ou cohérents,
de théories ».}\footnote{Ces deux corollaires sont la réponse à la
\emph{Question} de la page 23 : une théorie est « à limites finies » si et
seulement si ses modèles ensemblistes forment une catégorie localement
finiment présentable ; et un morphisme entre théories à limites finies est
un foncteur d'oubli entre catégories LFP qui préserve la présentation
finie, i.e. dont l'adjoint à gauche envoie $\Sigma$ dans $\Sigma'$ — ce
que le dictionnaire de la page 90 écrira « foncteurs $\tau \to \tau'$
commutant aux $\varinjlim$ quelconques et envoyant $S$ dans $S'$ ».
« Cohérent » est le bon mot : dans le vocabulaire des topos, un morphisme
de topos classifiants est cohérent lorsqu'il préserve les objets de
présentation finie, et c'est ici le même énoncé. Le manuscrit écrit
« $\mathbf{1}$-fidèle », dont le sens dans le dossier (page 97 : $\mathrm{Cat}
_{\lambda} \hookrightarrow \mathrm{Cat}$ est 1-fidèle) est « pleinement
fidèle sur les 1-flèches, et sur les 2-flèches » ; nous disons pleinement
fidèle.}

\emph{Question} (page 78) : pour $d'$ = les diagrammes finis
\emph{discrets} (produits finis), $R$ petite à produits finis, $\Sigma =
R^{\circ} \subset S = \underline{\mathrm{Hom}}_{\mathrm{prod}}(R, \mathrm{Ens})$ ;
on a $\Sigma \subset S_{\mathrm{pf}} = \Sigma_{0}$ et $S_{\mathrm{pf}} \simeq
\rho_{\Delta,\Delta_{0}}(\Sigma)$, et la question de retrouver $S$ à partir
de $\Sigma$ se ramène à celle d'une catégorie $\Sigma_{0} \in \Delta_{0}$
munie d'une sous-catégorie pleine $\Sigma$ stable par sommes
finies.\footnote{C'est le cas des théories de Lawvere, où $S$ est la
catégorie des algèbres, $\Sigma$ celle des algèbres \emph{libres} de type
fini, et $S_{\mathrm{pf}}$ celle des algèbres de présentation finie, qui
est la clôture de $\Sigma$ par coégalisateurs. La question est celle de
reconnaître, dans $S_{\mathrm{pf}}$, les objets libres ; elle n'a pas de
réponse intrinsèque en général (une catégorie LFP peut être la catégorie
des algèbres de plusieurs théories de Lawvere non équivalentes), et la
page la laisse ouverte.}

\pagerange{80}{80}
\subsection*{La liste des cas à traiter}

Le feuillet 80 dresse la liste des 2-catégories $\Delta$ pour lesquelles
le programme devrait être mené : (1) limites finies, foncteurs exacts ;
(1') limites de type $d'$ ; (1'') préordres à inf finis ; (2) limites
finies et colimites finies, foncteurs exacts ; (2') prétopos ; (2'')
catégories semi-additives ; (2''') additives ; (2'''') abéliennes ; (3)
$(\mathrm{Cat}_{d})$ avec $(d', d'')$ petits ; (4) limites finies et limites
filtrantes exactes ; (4'') préordres à inf finis, sup quelconques
distributifs ; (4') limites finies, colimites quelconques, sommes
disjointes universelles, limites filtrantes exactes ; (5) limites et
colimites quelconques ; (5') topos ; (6') topos essentiels ($\simeq
\widehat{C}$, ayant assez de points essentiels) avec les morphismes
essentiels. « Je m'attends à des ennuis (de sortie de l'Univers) avec (6)
et (6') ».\footnote{La liste est le programme de la logique catégorique
des années qui suivent : (1) théories cartésiennes, (2) régulières, (2')
cohérentes, (4') géométriques et leurs classifiants (5'), (2'')--(2'''')
les théories additives. Les cas (4) et (4') sont ceux où les
« catégories accessibles » des pages 54 à 61 interviennent ; (6') est le
cas des topos de préfaisceaux, i.e. des théories à limites finies, et
l'« ennui » prévu est celui de la page 84 : la catégorie des constructions
$B_{T}$ n'est pas visiblement une $\mathcal{U}$-catégorie.}

\pagerange{81}{87}
\section*{L'adjonction entre théories et catégories, et les constructions explicites (pages 81 à 87)}

\pagerange{81}{83}
\subsection*{$\sigma \dashv \rho$}

Soit $\text{Thé}_{\Delta}$ la 2-catégorie des théories sur $\Delta$ — les
catégories 2-cofibrées sur $\Delta$, i.e. les pseudo-foncteurs $T : \Delta
\to \mathrm{Cat}$ — et $\tau$ la théorie identique, $\tau(C) = C$. Deux
2-foncteurs :
\[
\sigma : \Delta \to \text{Thé}_{\Delta}, \quad \sigma(R) = \underline{\mathrm{Hom}}
_{\Delta}(R, -) ; \qquad
\rho : \text{Thé}_{\Delta} \to \Delta, \quad \rho(T) = B_{T} = \underline{\mathrm{Hom}}
(T, \tau),
\]
$B_{T}$ étant la \emph{catégorie des constructions} de $T$ : un objet en
est une façon de fabriquer, fonctoriellement en $C$ et en la structure, un
objet de $C$ à partir d'une $T$-structure sur $C$. On a
$\underline{\mathrm{Hom}}(\sigma(R), T) \simeq T(R)$ (Yoneda pour les
2-foncteurs), en particulier $\rho\sigma(R) = \underline{\mathrm{Hom}}(\sigma
(R), \tau) \simeq \tau(R) = R$ : $\sigma$ est 2-pleinement fidèle. La
question est de savoir quand
\[
\underline{\mathrm{Hom}}(\sigma(R), T) \longrightarrow \underline{\mathrm{Hom}}
_{\Delta}(\rho(T), R) = \sigma(\rho(T))(R)
\]
est une équivalence : \emph{pour tout $R$ si et seulement si $T$ est
2-représentable par $\rho(T)$}. Le morphisme d'unité $T \to \sigma\rho(T)$
est l'objet universel, et l'on demande si l'on peut décrire un objet
canonique de $T(\rho(T))$, i.e. la $T$-structure universelle sur la
catégorie des constructions.\footnote{C'est l'adjonction, aujourd'hui
classique, entre une théorie et sa « catégorie syntaxique » : pour une
théorie représentable, $B_{T} = R$ et la structure universelle est
l'identité de $R$, vue comme objet de $T(R) = \underline{\mathrm{Hom}}_{
\Delta}(R, R)$. Le manuscrit ajoute que $\rho(T) \in \mathrm{Ob}\,\Delta$
lorsque $\Delta$ est défini par des propriétés d'exactitude, puisque les
limites dans $\underline{\mathrm{Hom}}(T, \tau)$ se calculent argument par
argument, mais qu'il n'est pas clair que $B_{T}$ soit une
$\mathcal{U}$-catégorie — c'est le « grain de sel » de la page 84, et la
raison d'être des constructions explicites qui suivent.}

La page 82 énonce les deux propriétés de $\Delta$ dont tout dépend :
(1) pour $I$ petite, $C \mapsto \underline{\mathrm{Hom}}_{\mathrm{Cat}}(I, C)$
est 2-représentable par la $\Delta$-enveloppe $\widetilde{I}$ — l'inclusion
$\Delta \to \mathrm{Cat}$ a un 2-adjoint à gauche — ; (2) pour $J \in
\mathrm{Ob}\,\Delta$ et $\Sigma \subset \mathrm{Fl}(J)$ petit, $C \mapsto
\underline{\mathrm{Hom}}_{\Delta}(J, C ; \Sigma)$ est 2-représentable par
$J\Sigma^{-1}$. \emph{Conséquence} : les 2-limites inductives existent
dans $\Delta$. Pour un pseudo-foncteur $i \mapsto A_{i}$, soit $D =
\varinjlim A_{i}$ dans $\mathrm{Cat}$ ; alors $\underline{\mathrm{Hom}}_{
\mathrm{Cat}}(D, C) \simeq \varprojlim \underline{\mathrm{Hom}}_{\mathrm{Cat}}
(A_{i}, C)$, et la 2-limite dans $\Delta$ est le sous-2-foncteur des
$\tilde{u} : \widetilde{D} \to C$ tels que chaque $\tilde{u} \circ
\tilde\alpha_{i}$ soit un $\Delta$-foncteur, condition qui s'exprime en
rendant inversibles certaines flèches $\Sigma$ de $\widetilde{D}$ — du
moins si les $\Delta$-foncteurs sont définis par commutation à des limites
qui existent dans toutes les catégories de $\Delta$ ; et il faut que
$\Sigma$ soit petit, ce qui n'est garanti que si $d'$, $d''$ sont petits,
sinon il faut travailler avec des sous-catégories génératrices
petites.\footnote{« Construction des 2-limites dans $(\mathrm{Cat})$ : faite
dans SGA 4 VI en prenant la catégorie totale de la 2-catégorie cofibrée
et en rendant inversibles les morphismes cartésiens » — le renvoi est à
SGA 4 VI 6.}

\pagerange{83}{85}
\subsection*{Les constructions explicites $B'_{T}$}

Puisque $B_{T}$ n'est pas visiblement petite, on définit pas à pas une
catégorie $B'_{T}$ de \emph{constructions explicites} pour toute théorie
algébrique $T$ — la construction accompagnant la définition récurrente des
théories algébriques —, et l'on montre en même temps que $B'_{T} \in
\mathrm{Ob}\,\Delta$ et que $B'_{T}$ représente $T$ lorsque $T$ est
représentable.
\begin{enumerate}
\item[a)] \emph{Structure vide sur $I$ objets} : la théorie $\tau^{I}$,
$\tau^{I}(C) = C^{I}$, représentée par la $\Delta$-enveloppe $\widetilde{I}$
de la catégorie discrète $I$ ; $B'_{\tau^{I}} = \widetilde{I}$, que l'on
explicite en prenant les limites de type $d$ à partir des objets $i$ et les
nouvelles flèches que les axiomes d'exactitude imposent.
\item[b)] \emph{Ajouter des axiomes} : un ensemble $\Sigma$ de flèches de
$B'_{T}$ à rendre inversibles ; $B'_{T'} = B'_{T}\Sigma^{-1}$ dans $\Delta$.
Si $T$ est représentée par $R$ et $B'_{T} \simeq R$, alors $T'$ est
représentée par $B'_{T'}$ ; et si $T \to \tau^{I}$ est fidèle, $T' \to T
\to \tau^{I}$ aussi.\footnote{La marge demande de « préciser aussi les
axiomes du type flèches rendues égales », qui ne résultent pas des
inversions si $d' = d'' = \emptyset$ ; dans le cas général une égalité
$u = v$ s'impose en rendant inversible la flèche de l'égalisateur, comme
le dédoublement de la page 31 le fait pour les cônes.}
\item[c)] \emph{Ajouter des flèches} : un ensemble $H$ muni de $s, b : H
\to \mathrm{Ob}\,B'_{T}$ ; $T'(C)$ est la catégorie des $x \in T(C)$ munis
d'une famille $\varphi(h) : s(h)(x) \to b(h)(x)$, $h \in H$. Alors $B'_{T'}$
s'obtient de $B'_{T}$ en ajoutant librement les flèches $H$ (une catégorie
qui n'est pas dans $\Delta$), en prenant la $\Delta$-enveloppe $B'^{0}_{T'}$,
puis en rendant inversibles les flèches qui expriment que $B'_{T} \to
B'^{0}_{T'}$ est un $\Delta$-foncteur. Si $T$ est représentée par $R \simeq
B'_{T}$, $T'$ l'est par $B'_{T'}$ ; $T' \to T$ est fidèle.
\item[d)] \emph{Composer des structures} : $T = \prod T_{\alpha}$ est
algébrique et $B'_{T}$ est la 2-somme des $B'_{T_{\alpha}}$ dans $\Delta$ ;
d') à objets de base fixés, le produit 2-fibré $T'(C) = \prod^{(2)}_{\alpha,
C^{I}} T_{\alpha}(C)$ des $T_{\alpha}$ au-dessus de $\tau^{I}$, avec pour
$B'_{T'}$ la somme amalgamée des $B'_{T_{\alpha}}$ sous $\widetilde{I}$ ;
d'') limites projectives de théories, « il suffirait de le postuler pour
une catégorie d'indices bien ordonnée ».
\end{enumerate}

\emph{Remarque 1} : on ne rencontre ainsi que des théories représentées par
des $B_{T} = R$ munis d'une famille génératrice « fidèle » $I \to \mathrm{Ob}
\,R$, i.e. telle que $\underline{\mathrm{Hom}}_{\Delta}(R, C) \to C^{I}$ soit
fidèle pour tout $C$.

\pagerange{86}{87}
\subsection*{Réciproque : d'une catégorie à générateurs à une théorie}

\emph{Remarque 2}. Soit réciproquement $R \in \mathrm{Ob}\,\Delta$ avec des
$\Delta$-générateurs $I \to \mathrm{Ob}\,R$.

$1^{\circ}$) $R$ petite. On prend $I = \mathrm{Ob}\,R$, $T_{0} = \tau^{I}$,
$R_{0} = \widetilde{I}$ ; puis $H = \mathrm{Fl}(R)$ envoyé dans $\mathrm{Ob}\,
R_{0} \times \mathrm{Ob}\,R_{0}$, d'où $T_{1}$, $R_{1}$ ; puis on exprime que
la composition de $R$ devient celle de $R_{1}$, i.e. que l'application
« préfoncteur » $R \to R_{1}$ est un foncteur, d'où $T_{2}$, $R_{2} =
\widetilde{R}$ avec $T_{2}(C) \simeq \underline{\mathrm{Hom}}_{\mathrm{Cat}}
(R, C)$ ; enfin on exprime que $R \to R_{2}$ est un $\Delta$-foncteur en
rendant inversibles certaines flèches de $R_{2}$, d'où $T_{3}$ et $R_{3} =
R$. « On a appliqué une fois a), une fois c), deux fois b), et pas du tout
d). »

$2^{\circ}$) $R$ quelconque avec $I \to \mathrm{Ob}\,R$ générateurs. Soit
$\rho_{0}$ la sous-catégorie pleine engendrée par $I$, de sorte que
$\underline{\mathrm{Hom}}_{\Delta}(R, C) \to \underline{\mathrm{Hom}}_{\mathrm{Cat}}
(\rho_{0}, C) \simeq \underline{\mathrm{Hom}}_{\Delta}(\tilde\rho_{0}, C)$ est
fidèle ; on considère $R_{0} = \tilde\rho_{0} \to R$, les flèches $\Sigma$
qui deviennent inversibles, $R_{1} = R_{0}\Sigma^{-1}$, et on construit une
suite transfinie $R_{\alpha} \to R$ dans $\Delta$ dont toutes les théories
$\sigma(R_{\alpha})$ sont algébriques, le passage aux ordinaux limites se
faisant par d) ou d').

$4^{\circ}$) Précisément (page 87) : pour $R \in \mathrm{Ob}\,\Delta$ et $u : I
\to R$ un ensemble générateur, i.e. tel que $\underline{\mathrm{Hom}}_{\Delta}
(R, C) \to \underline{\mathrm{Hom}}_{\Delta}(\widetilde{I}, C)$ soit fidèle
pour tout $C$, on pose $R_{0} = \widetilde{I}$ avec $u_{0} : \widetilde{I}
\to R$ prolongeant $u$, $\Sigma_{\alpha}$ = les flèches de $R_{\alpha}$
qui deviennent inversibles dans $R$, $R_{\alpha+1} = R_{\alpha}
\Sigma_{\alpha}^{-1}$ au sens de $\Delta$, et $R_{\lim} = \varinjlim_{\Delta}
R_{\alpha}$ ; alors, \emph{pour $\alpha$ assez grand dans l'univers,
$R_{\alpha} \xrightarrow{\sim} R$}. Aux ordinaux limites on veut que $R_{\alpha}
\to R$ soit conservatif ; si les $\Delta$-foncteurs commutent aux noyaux
et conoyaux, les $R_{\alpha} \to R$ sont fidèles et les $R_{\alpha}$ sont des
sous-catégories ; pour tout cardinal $c$, si $\alpha$ est plus élevé que
$c$, $R_{\alpha}$ est stable sous les limites de types $\delta' \in d'$,
$\delta'' \in d''$ de cardinal $\leqslant c$ ; « donc si $d'$ et $d''$ sont
petits » — la phrase reste inachevée, et la marge dit ce qu'elle
concluait : la construction aboutit et $R$ est
algébrique.\footnote{C'est le théorème 4 a) de la page 18, généralisé de
« bons topos » à un $\Delta$ quelconque à types petits : \emph{un objet de
$\Delta$ à petite famille génératrice est la catégorie classifiante d'une
théorie algébrique}, et la démonstration est la construction transfinie
de la page 62. Les deux marges — « je ne sais s'il est 2-fidèle ??
pleinement fidèle », « on doit pouvoir prouver \ldots\ fidèle \ldots\ si
non grand \ldots » — sont sur le point qui manque : la conservativité aux
ordinaux limites. Dans le cas des théories à limites finies elle est
automatique et la conclusion est le théorème de la page 75.}

\pagerange{89}{91}
\section*{Les dictionnaires (pages 89 à 91)}

\pagerange{89}{89}
\subsection*{$\lambda$ fini : ce que $S$ retient de $R$}

Un demi-feuillet fixe, pour $\lambda$ un ensemble de diagrammes finis et
$R \in \mathrm{Cat}_{\lambda}$, les inclusions
\[
\Sigma = R^{\circ} \subset S = \underline{\mathrm{Hom}}_{\lambda}(R, \mathrm{Ens})
\subset \widehat{R^{\circ}} = \widehat{\Sigma}
\]
et les propriétés de $S$ : a) $S$ a les limites inductives quelconques ;
b) $S$ est « engendrée » par ses objets de type fini — ceux dont
$\mathrm{Hom}(X, -)$ commute aux limites filtrantes —, i.e. $S \to
\widehat{S_{0}}$ est conservatif, ou même pleinement fidèle, $S_{0}$ étant
la sous-catégorie pleine de ces objets ; c) $\Sigma \subset S_{0}$ et $\Sigma$
engendre $S_{0}$ par limites inductives finies, et $S$ par limites
inductives quelconques ; d) plus précisément $S_{0}$ est l'enveloppe de la
$\lambda$-catégorie $\Sigma$ sous les colimites finies, et $S$ l'enveloppe
de $S_{0}$ sous les colimites quelconques.\footnote{Pour $\lambda$ = tous
les diagrammes finis, c'est la page 77 : $S_{0} = \Sigma$, $S = \mathrm{Ind}
(\Sigma)$. Pour $\lambda$ plus petit (produits finis, par exemple),
$\Sigma \subsetneq S_{0}$ et d) est la description de $S$ comme
$\mathrm{Ind}$ de la clôture de $\Sigma$ par colimites finies, ce qui est
correct, avec la réserve que « objets de type fini » devrait ici se lire
« de présentation finie » — le manuscrit emploie l'un pour l'autre.} Le
titre qui suit, « 9. Relation avec analyseurs de Lazard », n'a pas de
texte.

\pagerange{90}{91}
\subsection*{$\lambda$ quelconque, puis $\lambda = (\lambda_{1}, \emptyset)$}

Le \emph{Dictionnaire pour $\lambda$ quelconque} met en regard les types
(théories) et les $\lambda$-catégories qui les représentent :

\begin{itemize}
\item un \emph{type} $T_{\lambda}$ $\longleftrightarrow$ une $\lambda$-catégorie
$R$ admettant une petite famille $\lambda$-génératrice ; un morphisme de
types $T_{\lambda} \to T'_{\lambda}$ $\longleftrightarrow$ un $\lambda$-foncteur
$R' \to R$ ;
\item un morphisme \emph{pleinement fidèle} (sous-théorie pleine $T' \subset
T$) $\longleftrightarrow$ un $\lambda$-foncteur de \emph{localisation} $R \to R'$
par rapport à un ensemble de flèches de $R$ stable par composition et par
$\lambda$-limites ;
\item un morphisme \emph{fidèle} $T' \to T$ (« structure plus riche »)
$\longleftrightarrow$ un $\lambda$-foncteur $f : R \to R'$ qui $\lambda$-engendre
$R'$ (« épimorphisme ? ») ;
\item un $\lambda$-type avec famille d'objets de base $(E_{i})_{i \in I}$
$\longleftrightarrow$ une $\lambda$-catégorie $R$ avec $I \to \mathrm{Ob}\,R$
qui $\lambda$-engendre $R$ ;
\item une limite projective de $\lambda$-théories $\longleftrightarrow$ une
limite inductive de $\lambda$-catégories.
\end{itemize}

Le \emph{Dictionnaire} pour $\lambda = (\lambda_{1}, \emptyset)$ — limites
projectives seulement — a cinq colonnes, et c'est le résumé le plus dense
du dossier ; nous le donnons ligne par ligne, avec pour colonnes le type
$T$, la $\lambda$-catégorie $R$, la $\lambda^{\circ}$-catégorie $S$, la
catégorie $\tau$ des modèles ensemblistes, et le topos $B$.

\emph{Les objets.} $T$ un $\lambda$-type, $T(C) = \underline{\mathrm{Hom}}_{
\lambda}(R, C) = \underline{\mathrm{Hom}}_{\lambda^{\circ}}(S^{\circ}, C)$, égal
à la sous-catégorie pleine de $\underline{\mathrm{Hom}}(C^{\circ}, \tau)$ des
foncteurs $u$ tels que les $\mathrm{Hom}(x, u(-))$, $x \in \mathrm{Ob}\,S$,
soient représentables, et à $\underline{\mathrm{Hom}}_{\mathrm{top}}(C, B)^{\circ}$
si $C$ est un topos. $R$ une $\lambda$-catégorie à petite famille
$\lambda$-génératrice ; $R = S^{\circ} = \mathrm{Pt\,ess}(B)$, la catégorie
des points essentiels du topos $B$. $S = R^{\circ}$, $\lambda^{\circ}$-catégorie
à petite famille $\lambda^{\circ}$-génératrice. $\tau$ une catégorie à
limites filtrantes (et limites projectives quelconques) munie d'une
sous-catégorie génératrice $S$ pour les foncteurs conservatifs, stable par
$\lambda^{\circ}$-limites, formée d'objets de présentation finie ; $\tau =
\underline{\mathrm{Hom}}_{\lambda}(R, \mathrm{Ens}) = T(\mathrm{Ens})$, $= \mathrm{Ind}
(S)$ si $\lambda_{1}$ est tous les diagrammes finis ; $\tau = \mathrm{Pt}
(B)$, $S = \mathrm{Pt\,ess}(B)$. $B$ un topos ayant suffisamment de points
essentiels et dont la catégorie des points essentiels est stable par
limites ; $B = \widehat{R} = \underline{\mathrm{Hom}}(S, \mathrm{Ens})$, le
topos classifiant de $T$.\footnote{Un point essentiel d'un topos $B$ est un
point $p$ dont l'image inverse $p^{*}$ a un adjoint à gauche $p_{!}$ ; pour
$B = \widehat{R}$, les points sont les foncteurs plats $R \to \mathrm{Ens}$
(ici les $\lambda$-foncteurs, i.e. $\tau$), et les points essentiels sont
les points représentables $\mathrm{Hom}(r, -)$, qui forment la
sous-catégorie $\Sigma = R^{\circ} \subset S = \tau$ de la page 89 ; c'est
elle que les colonnes $R$ et $S$ nomment, l'une $R$ et l'autre $R^{\circ}$,
à la convention près sur le sens des 2-flèches, qui porte les « $\circ$ »
de $\mathrm{Pt}(B)^{\circ}$ et $\mathrm{Pt\,ess}(B)^{\circ}$. Nous les
omettons, et nous lisons partout « la catégorie des points essentiels est
$\Sigma$, opposée de $R$ ». La colonne $\tau$ décrit une catégorie
localement présentable par sa sous-catégorie dense de présentables, et
« $\tau = \mathrm{Pt}(B)$ » est le théorème de Diaconescu vu depuis les
modèles.}

\emph{Les morphismes.} Un morphisme $T' \to T$ $\longleftrightarrow$ un
$\lambda$-foncteur $R \to R'$ $\longleftrightarrow$ un $\lambda^{\circ}$-foncteur
$S \to S'$ $\longleftrightarrow$ un foncteur $\tau \to \tau'$ commutant aux
limites inductives quelconques et envoyant $S$ dans $S'$ (ou son adjoint
à droite $\tau' \to \tau$) $\longleftrightarrow$ un morphisme de topos $B' \to
B$ dont l'image inverse transforme objets ponctuels (représentables) en
objets ponctuels, déduit par localisation de $B$ par rapport à un ensemble
de flèches de $R \subset B$.

\emph{Pleinement fidèle} $T' \to T$ $\longleftrightarrow$ localisation $R
\to R'$ $\longleftrightarrow$ localisation $S \to S'$ $\longleftrightarrow$
$\tau' \to \tau$ pleinement fidèle $\longleftrightarrow$ $B' \to B$ un
plongement, i.e. $f_{*}$ pleinement fidèle.

\emph{Sous-théorie pleine} $\longleftrightarrow$ ensemble $M$ de flèches de
$R$ stable par composition et par $\lambda$-limites $\longleftrightarrow$
ensemble de flèches de $S$ stable par composition et $\lambda^{\circ}$-limites
$\longleftrightarrow$ sous-catégorie pleine $\tau'$ de $\tau$ stable par
limites projectives, formée des $X$ tels que $\mathrm{Hom}(u, X)$ soit
bijectif pour les $u$ que les objets de $\tau'$ rendent bijectifs
$\longleftrightarrow$ sous-topos $B'$ de $B$ dont le plongement $f$ a une
image inverse préservant les objets ponctuels.\footnote{C'est le
\emph{Problème 2} de la page 17 dans toutes ses colonnes, et, colonne
$\tau$, la description des sous-catégories réflexives par
\emph{orthogonalité} : $\tau'$ est l'ensemble des objets orthogonaux à
une classe de flèches, ce qui est la forme moderne des « axiomes ».}

\emph{Fidèle} $T' \to T$ $\longleftrightarrow$ $R \to R'$ $\lambda$-engendre
$R'$ $\longleftrightarrow$ $S \to S'$ $\lambda^{\circ}$-engendre $S'$
$\longleftrightarrow$ $\tau' \to \tau$ fidèle $\longleftrightarrow$ $f^{*} :
B \to B'$ fidèle.

\emph{Limites projectives} de $\lambda$-types $\longleftrightarrow$ limites
inductives de $\lambda$-catégories $\longleftrightarrow$ limites inductives
de $\lambda^{\circ}$-catégories $\longleftrightarrow$ limites projectives des
$\tau_{i}$ $\longleftrightarrow$ limites projectives des topos.\footnote{Le
dictionnaire ne démontre rien ; il enregistre les pages 62 à 78 sous forme
de table, et il est correct dans toutes ses cases à deux réserves près :
la colonne $\tau$ de la ligne « fidèle » porte un « (ou plutôt ??) »
qui est le sien, et la ligne « limites projectives des $\tau_{i}$ » n'est
vraie, comme il le note, que « si $\lambda_{1}$ = tous les diagrammes
finis », i.e. quand la limite des catégories LFP est LFP, ce qui est un
théorème (Makkai--Paré) et non une évidence.}
\pagerange{92}{95}
\section*{Le corps universel (pages 92 à 95)}

Un brouillon rapide, sur deux colonnes, prend l'exemple qui met la
distinction « limites projectives / limites inductives » à l'épreuve. Soit
$T$ l'espèce de structure « anneau » et $T^{\circ}$ celle de « corps ». La
seconde n'est pas définissable par limites projectives : la catégorie des
corps n'est pas stable par produits. Les sous-types \emph{pleins} d'un
type définissable par limites projectives le sont encore, et sont donnés
par une sous-catégorie pleine $\tau'$ de $\tau = T(\mathrm{Ens})$ décrite par
orthogonalité — $C \in \tau'$ si et seulement si, pour tout homomorphisme
$A \to B$ que les objets de $\tau'$ rendent bijectif, $\mathrm{Hom}(B, C)
\to \mathrm{Hom}(A, C)$ est bijectif ; mais « corps » n'est pas un sous-type
plein de « anneau ».\footnote{Le début de la page 92 est le plus abîmé du
dossier ; la ligne sur les « pseudo-corps » et l'homomorphisme
$\mathbb{Z}[H] \to \mathbb{Z}[t, t^{-1}] \times \mathbb{Z}$, dont la moitié
est illisible, prépare ce qui suit : le plus petit sous-type plein de
« anneau » contenant les corps est celui des anneaux \emph{réduits} (page
41, article 11), et la sous-catégorie pleine des pseudo-corps est
l'orthogonale décrite ici. Nous ne reconstruisons pas la ligne.} Peut-on
définir le sous-objet $A^{*}$ des éléments inversibles d'un anneau $A$ par
limites projectives et inductives ? Oui : dans $R = (\mathrm{Ann})^{\circ} =
(\mathrm{Sch\,aff})$, l'anneau universel est l'objet $\mathbb{Z}[t]$, et
$A^{*}$ correspond au sous-objet $\mathbb{Z}[t, t^{-1}] = \mathbb{Z}[x,y]/
(xy - 1)$, qui est une limite projective finie. Pour $A$ un objet anneau
d'une catégorie à limites finies, on a donc un morphisme $e \sqcup A^{*}
\to A$ — la section nulle et l'inclusion des inversibles —, et \emph{on
dira que $A$ est un corps si ce morphisme est un isomorphisme}.

Plaçons-nous dans le topos universel $\widehat{R}$, avec l'objet anneau $O
= \mathrm{Sp}(\mathbb{Z}[t])$. Il faut rendre inversible la flèche $\varphi
: e \sqcup O^{*} \to O$ — qui n'est pas une flèche de $R$, mais de
$\widehat{R}$ — ; elle se factorise

\begin{tikzcd}
e \sqcup O^{*} \arrow[r, "\varphi"] \arrow[d, "j"'] & O \\
D \arrow[ur, "i"'] &
\end{tikzcd}

où $i : D \to O$ est dans $R$, et rendre $\varphi$ inversible revient à
rendre $i$ et $j$ inversibles. On rend d'abord inversibles les morphismes
$i$ (qui sont dans $R$) : $R'$ est la catégorie de fractions de $R$
correspondante, et $\widehat{R'}$ classifie les pseudo-corps. Il faut
ensuite rendre $e \sqcup O^{*} \to O$ un monomorphisme, i.e. $e \cap O^{*}$
vide, ce qui rend le préfaisceau $\varepsilon$ ($\varepsilon(X) = \{e\}$ si
$X = \emptyset$, vide sinon) égal à $\emptyset$. On trouve : \emph{si $C$
est une catégorie à limites projectives finies et limites inductives
quelconques, les corps dans $C$ correspondent aux foncteurs $F : R' \to C$
exacts à gauche et commutant aux sommes} ; \emph{la catégorie $C$
universelle pour cela est la catégorie des faisceaux sur $R'$ pour la
topologie de Zariski — c'est un topos —, et le corps universel est le
faisceau structural sur ce topos}.\footnote{C'est la description
du \emph{topos classifiant des corps}, qui n'est pas un topos de
préfaisceaux : la disjonction « $x = 0$ ou $x$ inversible » se traduit en
demandant que $\mathrm{Spec}\,A$ soit recouvert par $V(f)$ et $D(f)$ pour
tout $f$, ce qui est la topologie de la page 41 (article 11, « sections
sur chaque fibre »), et la condition « $0 \neq 1$ » en demandant que le
schéma vide soit couvert par la famille vide. Le manuscrit passe par les
pseudo-corps (rendre $D \to O$ inversible, i.e. $A \simeq A_{f} \times
A/f$ pour tout $f$), ce qui est correct : un corps est un anneau
absolument plat local, ou, comme la page 41 le dit, « pseudo-corps
local ». Le mot « Zariski » pour la topologie sur $R'$ est le sien ; sur
la catégorie $R'$ des schémas absolument plats les recouvrements de
Zariski sont les recouvrements $\{V(f), D(f)\}$, et c'est bien ce qu'il
faut. « Le corps universel est le faisceau structural » : l'anneau
générique du topos classifiant est un corps au sens interne, énoncé qui
est aujourd'hui le premier exemple de la théorie (Johnstone, Wraith).}

La page 95 pose ensuite, pour les anneaux intègres, la question du
sous-type $O' \subset O$ des éléments « qui opèrent régulièrement sur $O$ »
(les non-diviseurs de zéro), avec $O' \supset O^{*}$ : est-il définissable
par limites projectives et inductives ? Dans le cas universel, $O'(X)$ est
l'ensemble des $f \in A = O(X)$ qui sont réguliers après toute extension
$A \to A'$, ce qui force $f \notin M$ pour tout idéal maximal $M$, donc
$O' = O^{*}$ ; et sur $\widehat{R}$ l'anneau universel est localement un
corps, donc $O'_{p} = O^{*}_{p}$, i.e. $A' = A^{*}$.\footnote{La conclusion
est exacte pour la raison qu'il indique : un élément régulier
\emph{universellement} — après tout changement de base — est inversible,
car son image dans le corps résiduel d'un idéal maximal doit être non
nulle. Le sous-type des non-diviseurs de zéro tout court n'est pas
géométrique (il n'est pas stable par changement de base), et c'est la
réponse à la question posée.}

\pagerange{97}{106}
\section*{$\lambda$-types : la copie au net (pages 97 à 106)}

Dix feuillets paginés de ① à ⑩, d'une écriture posée, sont l'état le plus
achevé du dossier : la définition des théories relatives à un couple de
types de limites, la représentabilité, la définition récurrente des
théories algébriques, et, pour les théories définies par limites
projectives, la reconstruction de la théorie à partir de ses modèles. Les
alinéas sont numérotés [1] à [8] et les formules (1.1) à (8.2.17) ; nous
gardons ces numéros.

\pagerange{97}{97}
\subsection*{[1] $\lambda$-catégories et $\lambda$-types}

Soit $\lambda = (\overleftarrow{\lambda}, \overrightarrow{\lambda})$ un
couple d'ensembles de diagrammes de l'univers fixé $\mathcal{U}$ (1.1). La
2-catégorie $(\mathrm{Cat}_{\lambda})$ (1.2) a pour objets les
\emph{$\lambda$-catégories} — celles qui ont les limites projectives de
type $\overleftarrow{\lambda}$ et les limites inductives de type
$\overrightarrow{\lambda}$ —, pour 1-flèches les \emph{$\lambda$-foncteurs},
qui y commutent, et pour 2-flèches les transformations naturelles ; ainsi
$\mathrm{Cat}_{\lambda} \hookrightarrow \mathrm{Cat}$ est pleinement fidèle
sur les 2-flèches. Un \emph{$\lambda$-type} $T$ est une espèce de structure
qui, à toute $\lambda$-catégorie $C$, associe la catégorie $T(C)$ des
structures de type $T$ dans $C$, fonctoriellement en $C$ : une catégorie
cofibrée $T_{\lambda}$ sur $\mathrm{Cat}_{\lambda}$ (1.3). Les morphismes
sont ceux des catégories cofibrées,
\begin{equation}
\mathrm{Hom}_{\lambda}(T, T') = \mathrm{Hom}_{\mathrm{cofib}/(\mathrm{Cat}_{
\lambda})}(T_{\lambda}, T'_{\lambda}), \tag{1.4}
\end{equation}
et les $\lambda$-types forment une 2-catégorie $(\lambda\text{-types})
\hookrightarrow \text{2-Cofib}(\mathrm{Cat}_{\lambda})$ (1.5), 2-pleinement
fidèle.\footnote{$\mathrm{Cat}_{\lambda}$ est la $\Delta$ des pages
précédentes avec $d = \lambda$, et $T_{\lambda}$ est la « théorie sur
$\Delta$ » de la page 83 ; la copie au net remplace $\Delta$ par
$\mathrm{Cat}_{\lambda}$ et « théorie » par « type » — nous suivons la
copie dans cette section et renvoyons, pour les mêmes objets, aux pages 62
à 87 par leurs numéros. Le « 1-fidèle » de (1.5) et de « $\mathrm{Cat}_{
\lambda} \hookrightarrow \mathrm{Cat}$ » désigne un 2-foncteur pleinement
fidèle sur les 2-flèches (une sous-2-catégorie « 2-pleine », mais non
pleine sur les 1-flèches) ; cette lecture est celle que les usages
ultérieurs — « 2-fidèle » pour la 2-catégorie des $\lambda$-types dans
les cofibrées, page 97 — imposent.}

\pagerange{97}{98}
\subsection*{[2] Objets de base ; $\lambda$-$I$-types}

Si $T$ est « une structure sur $I$ objets de base », on introduit le type
$\tau^{I}$ de la structure vide sur $I$ objets, $\tau^{I}(C) = C^{I}$
(2.1), et le foncteur canonique « système des objets de base »
\begin{equation}
b^{C,T}_{\lambda} : T_{\lambda}(C) \longrightarrow C^{I}, \tag{2.2}
\end{equation}
fidèle, définissant pour $C$ variable un morphisme fidèle $b^{T}_{\lambda}
: T_{\lambda} \to \tau^{I}$ (2.3).\footnote{« Foncteur 0-fidèle » dans le
manuscrit : fidèle sur les 1-flèches, sans rien demander des objets. Que
la structure détermine \emph{fidèlement} les objets de base est ce que
« objets de base » veut dire ; c'est la Remarque 1 de la page 86.} Pour
$T$, $T'$ deux types relatifs à $I$, on pose
\begin{equation}
\mathrm{Hom}_{\lambda}(I ; T, T') = \text{catégorie des couples } (u, \alpha),
\tag{2.4}
\end{equation}
$u : T_{\lambda} \to T'_{\lambda}$ un morphisme et $\alpha$ un isomorphisme
$b^{T'} \circ u \simeq b^{T}$ :

\begin{tikzcd}
T_{\lambda} \arrow[rr, "u"] \arrow[dr, "b^{T}_{\lambda}"'] & & T'_{\lambda} \arrow[dl, "b^{T'}_{\lambda}"] \\
& \tau^{I} &
\end{tikzcd}

d'où la 2-catégorie $(\lambda\text{-}I\text{-types}) \hookrightarrow
\text{2-Cofib}(\mathrm{Cat}_{\lambda})/\tau^{I}$ (2.5) et le 2-foncteur d'oubli
$(\lambda\text{-}I\text{-types}) \to (\lambda\text{-types})$ (2.6), qui est
fidèle.

\pagerange{98}{99}
\subsection*{[3] Constructions ; [4] représentabilité}

Si $I$ est réduit à un point, $\tau^{I} = \tau$, $\tau(C) = C$ (3.1), et
l'on pose
\begin{equation}
R^{\lambda}_{T} = \mathrm{Hom}_{\lambda}(T, \tau), \tag{3.2}
\end{equation}
la \emph{catégorie des $\lambda$-constructions} pour $T$. On ne sait pas a
priori si c'est une $\mathcal{U}$-catégorie ; mais c'est une
$\lambda$-catégorie, les $\lambda$-limites s'y calculant « argument par
argument », et pour toute $C$ et tout $\xi \in T(C)$ le foncteur
d'évaluation
\begin{equation}
\xi^{\lambda} : R^{\lambda}_{T} \longrightarrow C, \qquad \rho \mapsto \rho(\xi),
\tag{3.3}
\end{equation}
est un $\lambda$-foncteur ; on a donc un accouplement $R^{\lambda}_{T} \times
T(C) \to C$ (3.4), i.e. un $\lambda$-foncteur $R^{\lambda}_{T} \to \underline{
\mathrm{Hom}}(T(C), C)$ (3.5). Si $T$ est donné avec $I$, on a de plus $I
\to \mathrm{Ob}\,R^{\lambda}_{T}$, fonctoriel en $T$ (3.6) : chaque objet de
base est une construction.\footnote{Le manuscrit numérote (3.5) deux
formules, la seconde étant notre (3.6) ; nous renumérotons. La catégorie
des constructions est ce que la page 81 note $B_{T}$ ; une construction
est un « terme » de la théorie, un moyen de fabriquer un objet de $C$ à
partir de toute $T$-structure sur $C$, et $\xi^{\lambda}$ est son
évaluation. La marge gauche de la page 98 annonce des formules (3.7) et
(3.8) — le foncteur $T(C) \to \underline{\mathrm{Hom}}_{\lambda}(R^{\lambda}_{T},
C)$ transposé de (3.4) — qui sont celles de [4].}

\textbf{Définition.} \emph{$T$ est $\lambda$-représentable si la catégorie
cofibrée $T_{\lambda}$ est 2-représentable ; l'objet $R_{T}$ qui la
représente est la \emph{catégorie $\lambda$-modelante} de $T$ :}
\begin{equation}
T(C) \cong \mathrm{Hom}_{\lambda}(R_{T}, C) \tag{4.1}
\end{equation}
\emph{via une $T$-structure $\lambda$-universelle $\xi^{\lambda}_{T} \in
T(R_{T})$ (4.2).}\footnote{Le manuscrit écrit « catégorie
$\lambda$-modelaire » ; la page 111 dira « topos modulaire » et
« catégorie modulaire ». Nous disons \emph{modelante} ici et \emph{topos
classifiant} plus loin, en signalant chaque fois le mot du manuscrit ;
c'est la catégorie que la littérature appelle \emph{syntaxique} ou
\emph{classifiante}.}

$R_{T}$ est défini à équivalence près, unique à isomorphisme unique près,
et pour toute catégorie cofibrée $\mathcal{F}$ sur $\mathrm{Cat}_{\lambda}$,
\begin{equation}
\mathrm{Hom}_{\mathrm{cofib}}(T_{\lambda}, \mathcal{F}) \cong \mathcal{F}(R_{T})
\tag{4.3}
\end{equation}
(Yoneda) ; d'où $\mathrm{Hom}_{\lambda}(T, T') \cong T'(R_{T})$ (4.4), et
pour $T' = \tau$ :
\begin{equation}
R^{\lambda}_{T} = \mathrm{Hom}_{\lambda}(T, \tau) \cong \tau(R_{T}) = R_{T}.
\tag{4.5}
\end{equation}
\emph{Donc, pour un type représentable, la catégorie des constructions est
la catégorie modelante, et c'est une $\mathcal{U}$-catégorie.} Pour $C =
(\mathrm{Ens})$ :
\begin{equation}
T(\mathrm{Ens}) \cong \mathrm{Hom}_{\lambda}(R^{\lambda}_{T}, \mathrm{Ens}).
\tag{4.6}
\end{equation}

\pagerange{99}{101}
\subsection*{[5] Les règles de formation, et les questions}

On veut caractériser les catégories cofibrées sur $\mathrm{Cat}_{\lambda}$
(munies d'un morphisme dans $\tau^{I}$) qui proviennent de « théories de
types de structures », i.e. donner une définition de $(\lambda\text{-types})$
autre que (1.5). On introduit $S^{\lambda}_{T,I} \subset R^{\lambda}_{T}$, la
sous-catégorie pleine engendrée par les objets de base sous les
$\lambda$-limites, et les règles :
\begin{enumerate}
\item[a)] (structure vide) $\tau^{I} \in \mathrm{Ob}(\lambda\text{-}I\text{-types})$
pour $I \in \mathcal{U}$ ;
\item[b)] (ajouter des axiomes) si $T$ est un $\lambda$-$I$-type et $A
\subset \mathrm{Fl}(S^{\lambda}_{T,I})$, la sous-catégorie cofibrée pleine
$T'_{A}$ des $\xi \in T(C)$ tels que $\xi^{\lambda}$ (3.3) transforme les
flèches de $A$ en isomorphismes est un $\lambda$-$I$-type ;
\item[c)] (ajouter des données) si $T$ est un $\lambda$-$I$-type et $\alpha
: B \to \mathrm{Ob}\,S^{\lambda}_{T,I} \times \mathrm{Ob}\,S^{\lambda}_{T,I}$,
la catégorie cofibrée $T_{B,\alpha}$ des couples $(\xi, (u_{\beta})_{\beta
\in B})$, $\xi \in T(C)$ et $u_{\beta} \in \mathrm{Hom}_{C}(\rho(\xi),
\sigma(\xi))$ pour $\alpha(\beta) = (\rho, \sigma)$, est un
$\lambda$-$I$-type ;
\item[d)] (composition) si $(T_{j})_{j \in J}$ est une petite famille de
$\lambda$-$I$-types, le produit 2-fibré des $T_{j}(C)$ sur $C^{I}$ est un
$\lambda$-$I$-type.
\end{enumerate}
Un \emph{N.B.} demande s'il faut aussi une règle d'égalisation de couples
de flèches.

\textbf{Définition 5.1.} \emph{La 2-catégorie des $\lambda$-$I$-types est la
plus petite sous-2-catégorie 2-pleine de $\text{2-Cofib}(\mathrm{Cat}_{
\lambda})/\tau^{I}$ stable par les opérations a) b) c) d). On peut en donner
une construction transfinie.}\footnote{C'est la définition « par
présentation » des théories : sortes, opérations, axiomes, produits
fibrés ; b) code les axiomes par l'inversibilité d'une flèche entre
constructions, ce qui, avec le dédoublement des sommets (page 31),
couvre les égalités et les conditions d'exactitude — la question du
\emph{N.B.} a donc une réponse positive, que la page 84 avait déjà donnée
en marge. L'ensemble a)--d) est ce que la théorie des esquisses formalise
en une fois : un $\lambda$-$I$-type est une esquisse petite à limites de
types $\lambda$, et ses modèles dans $C$ sont $T(C)$.}

\textbf{Question 5.1.} \emph{Les $\lambda$-$I$-types sont-ils exactement les
catégories cofibrées $\lambda$-représentables par un $R^{\lambda}_{T} \in
\mathrm{Ob}\,\mathrm{Cat}_{\lambda}$ muni d'un $I \to \mathrm{Ob}\,R^{\lambda}_{T}$
tel que $\mathrm{Ob}\,R$ soit la plus petite partie stable par les
$\lambda$-limites et contenant $I$ ?}

Le paragraphe 5.4 décompose la question en quatre problèmes de
constructions universelles dans $\mathrm{Cat}_{\lambda}$, chacun avec une
question de $\lambda$-engendrement :
\begin{itemize}
\item[a')] $\tau^{I}$ est-il représentable par une $\lambda$-catégorie
libre $R_{I}$ sur $I$ objets, $\lambda$-engendrée par l'image de $I$ ?
\item[b')] pour $R \in \mathrm{Cat}_{\lambda}$ et $A \subset \mathrm{Fl}\,R$,
existe-t-il $R_{A}$ universelle rendant inversibles les flèches de $A$, et
est-elle $\lambda$-engendrée par l'image de $R$ ?
\item[c')] pour $\alpha : B \to \mathrm{Ob}\,R \times \mathrm{Ob}\,R$,
existe-t-il $R_{B,\alpha}$ universelle munie de $i : R \to R_{B,\alpha}$ et
de flèches indexées par $B$ entre les $i(\rho)$, $i(\sigma)$, et est-elle
$\lambda$-engendrée par l'image de $i$ ?
\item[d')] pour une petite famille $(R_{j})$ munie de $I \to \mathrm{Ob}\,
R_{j}$, existe-t-il une somme amalgamée $R$ sous $I$ dans $\mathrm{Cat}_{
\lambda}$, avec un système transitif d'isomorphismes entre les images de
$I$, $\lambda$-engendrée par la réunion des images des $R_{j}$ ?
\end{itemize}
Les existences sont acquises par le théorème d'enveloppe des pages 3 et
27, et seuls les $\lambda$-engendrements sont en question — c'est
l'« essentiel » de 5.4.\footnote{Ce sont, dans l'ordre, les problèmes 1,
2, « ajouter des flèches » et « composition » des pages 82--85 ; la
question de l'engendrement est celle de savoir si la clôture transfinie
sous les $\lambda$-limites de l'image de $I$ est toute la catégorie
construite, ce qui est vrai par construction lorsque $\lambda$ est petit
(page 76, \emph{Question}).} Le paragraphe 5.5 propose un critère de
$\lambda$-engendrement par fidélité : si $S \subset \mathrm{Ob}\,R$ est tel
que, pour toute $\lambda$-catégorie $R'$ et tous $\lambda$-foncteurs $F, G :
R \to R'$, $\mathrm{Hom}(F, G) \to \prod_{s \in S} \mathrm{Hom}(F(s), G(s))$
soit injectif, $R$ est-elle $\lambda$-engendrée par $S$ ? La réciproque
est évidente.\footnote{La réponse est négative en général : « $S$
détecte les 2-flèches » est strictement plus faible que « $S$ engendre
par $\lambda$-limites » — la première propriété ne voit pas les
facteurs directs, par exemple —, et 5.6 en tire la conséquence qu'il
faut « peut-être renforcer la notion de sous-catégorie
$\lambda$-engendrée, en tenant compte des flèches, pas seulement des
objets ».} Les paragraphes 5.6 et 5.7 posent la réciproque de la question
5.1 : si $R$ est $\lambda$-engendrée par $I$, la catégorie cofibrée qu'elle
définit, munie de $T \to \tau^{I}$, est-elle un $\lambda$-$I$-type ? Et,
$I$ omis, un $\lambda$-type est un $\lambda$-$I$-type pour $I$ et $\varphi$
convenables, i.e. une catégorie cofibrée $\lambda$-représentable par un $R$
dont les objets admettent une petite partie $\lambda$-génératrice « en un
sens convenable à dégager ».

\pagerange{102}{103}
\subsection*{[6] $\lambda$ variable ; [7] dualité}

Pour $\lambda \subset \lambda'$ (plus de limites exigées), on a un
2-foncteur de restriction
\begin{equation}
\eta_{\lambda\lambda'} : (\lambda\text{-types}) \longrightarrow
(\lambda'\text{-types}), \qquad \eta(T_{\lambda})(C) = T_{\lambda}(C)
\ \text{pour } C \in \mathrm{Ob}\,\mathrm{Cat}_{\lambda'} \subset \mathrm{Ob}\,
\mathrm{Cat}_{\lambda}, \tag{6.1}
\end{equation}
qui correspond, sur les catégories modelantes, à l'enveloppe
$\mathrm{Cat}_{\lambda} \to \mathrm{Cat}_{\lambda'}$ (6.2), la
« $\lambda'$-catégorie enveloppe ». Exemples : pour $\lambda$ = limites
projectives finies et $\lambda'$ = limites projectives quelconques,
l'enveloppe est $C \mapsto \mathrm{Pro}(C)$ ; pour les limites inductives,
$C \mapsto \mathrm{Ind}(C)$. Il faudrait des conditions sous lesquelles
(6.1) est pleinement fidèle : $\mathrm{Hom}_{\lambda}(T_{1}, T_{2}) \to
\mathrm{Hom}_{\lambda'}(T_{1}, T_{2})$ est-il une équivalence ? « Sans doute
faux » en général, « mais pleinement fidèle si $\overrightarrow{\lambda} =
\emptyset$ ».\footnote{Pour $\lambda = $ limites finies, $\lambda' = $
toutes les limites, c'est le second corollaire de la page 78 :
pleinement fidèle, d'image les morphismes cohérents. La réserve
« $\overrightarrow{\lambda} = \emptyset$ » est justifiée : dès qu'il y a
des limites inductives dans $\lambda$, l'enveloppe $\mathrm{Cat}_{\lambda}
\to \mathrm{Cat}_{\lambda'}$ ne conserve plus les catégories modelantes
« à générateurs près », et la restriction perd des morphismes. Le
manuscrit nomme les $\lambda'$-catégories enveloppes « analysantes »,
mot que nous lisons sous réserve.}

\emph{Dualité.} L'opposé d'un type $\lambda$ échange $\overleftarrow{\lambda}$
et $\overrightarrow{\lambda}$ ; à tout $\lambda$-type $T$ on associe le
$\lambda^{\circ}$-type $T^{\circ}$,
\begin{equation}
T^{\circ}(C) = T(C^{\circ})^{\circ}, \tag{7.1}
\end{equation}
i.e. $T^{\circ}_{\lambda^{\circ}} \cong (T_{\lambda})^{\circ}$ composé avec
le passage à l'opposé $\mathrm{Cat}_{\lambda^{\circ}} \simeq \mathrm{Cat}_{
\lambda}$ ; d'où des 2-équivalences $(\lambda\text{-types}) \cong
(\lambda^{\circ}\text{-types})$ (7.3), de même avec $I$ (7.4), compatibles
avec $C \mapsto C^{\circ}$ sur les catégories modelantes
(7.6).\footnote{C'est pourquoi tout ce qui suit sur $\overrightarrow{\lambda}
= \emptyset$ vaut aussi pour $\overleftarrow{\lambda} = \emptyset$, « bien
sûr, on a un énoncé dual pour les $\varinjlim$ » (page 104). La
numérotation saute (7.2) et (7.5), biffés sur la page.}

\pagerange{103}{104}
\subsection*{[8] Structures définies par limites projectives : limites, pleine fidélité, carrés 2-cartésiens}

On suppose désormais $\overrightarrow{\lambda} = \emptyset$, et $T$
désigne un $\lambda$-type.

\textbf{8.1.} \emph{Pour toute $\lambda$-catégorie $C$, les limites
projectives de type $\lambda$ existent dans $T(C)$, et plus généralement
les limites projectives de tout type $d$ qui existent dans $C$ ; les
foncteurs construction $T(C) \to C$ y commutent.} Écrire $T(C) \cong
\mathrm{Hom}_{\lambda}(R_{T}, C)$ ; les limites de type $d$ de foncteurs $R_{T}
\to C$ se calculent argument par argument dans $\mathrm{Hom}(R_{T}, C)$, et
les $\lambda$-foncteurs y sont stables parce que les limites commutent aux
limites. \textbf{8.1.1.} \emph{Plus généralement, pour $\lambda$ quelconque :
si $C$ admet les limites projectives de type $d$ et si celles-ci
commutent dans $C$ aux limites inductives de type $\overrightarrow{\lambda}$,
alors $T(C)$ admet les limites projectives de type $d$, calculées argument
par argument, et toute construction $\rho : T(C) \to C$ y commute ; de
même $T(C) \to T(C')$ pour $C \to C'$ commutant à ces limites, et $T(C) \to
T'(C)$ pour $T \to T'$.} C'est la proposition de la page 65, au net.

\textbf{8.2.} \emph{Pour tout $\lambda$-type $T$ et tout $\lambda$-foncteur
pleinement fidèle (resp. fidèle) $i : C \to C'$ (8.2.1), $T(i) : T(C) \to
T(C')$ (8.2.2) est pleinement fidèle (resp. fidèle). Si de plus $T$ est
muni d'une base fidèle $b : T \to \tau^{I}$, i.e. si $R = R^{\lambda}_{T}$ est
muni de $\varphi : I \to \mathrm{Ob}\,R$ qui le $\lambda$-engendre, alors,
dans le cas pleinement fidèle, le carré}

\begin{tikzcd}
T(C) \arrow[r] \arrow[d] & T(C') \arrow[d] \\
C^{I} \arrow[r, hook] & C'^{I}
\end{tikzcd}

\emph{(8.2.3) est 2-cartésien.} En termes de $R$ : la sous-catégorie pleine
$\mathrm{Hom}_{\lambda}(R, C)$ de $\mathrm{Hom}_{\lambda}(R, C')$ est formée des
$\xi : R \to C'$ tels que $\xi(\varphi(i)) \in C$ pour tout $i \in I$ — car
$I$ $\lambda$-engendre $R$ et $\xi$ commute aux $\lambda$-limites, donc
envoie tout $R$ dans l'image essentielle de $C$, qui est stable par ces
limites. « En termes imagés » : pour une famille $(E_{i})$ d'objets de $C$,
il revient au même de se donner dessus une structure d'espèce $T$ dans $C$
ou sur leurs images dans $C'$.

Appliquons ceci au plongement de Yoneda $i : C \hookrightarrow \widehat{C}$
(8.2.4), qui est un $\lambda$-foncteur puisque $\overrightarrow{\lambda} =
\emptyset$ : le carré

\begin{tikzcd}
T(C) \arrow[r] \arrow[d] & T(\widehat{C}) \arrow[d] \\
C^{I} \arrow[r] & \widehat{C}^{I}
\end{tikzcd}

\emph{(8.2.5) est 2-cartésien.} Puis, le foncteur contravariant canonique de
$\widehat{C}$ dans $\mathrm{Hom}(\widehat{C}, \mathrm{Ens})$ restreint à $C$
donne $C^{\circ} \to \mathrm{Hom}_{\lambda}(\widehat{C}, \mathrm{Ens})$, $X
\mapsto (F \mapsto F(X))$ (8.2.6), à valeurs dans les foncteurs commutant à
tout type « petit » de limites ; d'où, en composant avec $\mathrm{Hom}_{
\lambda}(\widehat{C}, \mathrm{Ens}) \to \mathrm{Hom}(T(\widehat{C}),
T(\mathrm{Ens}))$, un foncteur canonique $C^{\circ} \to \mathrm{Hom}(T(\widehat{C}),
T(\mathrm{Ens}))$ (8.2.7), ou, ce qui revient au même,
\begin{equation}
\alpha : T(\widehat{C}) \xrightarrow{\ \approx\ } \mathrm{Hom}(C^{\circ}, S_{T}),
\qquad S_{T} = T(\mathrm{Ens}). \tag{8.2.8}
\end{equation}
\emph{Le foncteur $\alpha$ est toujours une équivalence} : $\mathrm{Hom}_{
\lambda}(R, \widehat{C}) \cong \mathrm{Hom}(R \times C^{\circ}, \mathrm{Ens})
\cong \mathrm{Hom}(C^{\circ}, \mathrm{Hom}_{\lambda}(R, \mathrm{Ens}))$, les
limites de type $\lambda$ dans $\widehat{C}$ se calculant argument par
argument (formule encadrée en marge de la page 105).\footnote{Le
manuscrit fait précéder l'assertion de deux hypothèses de taille —
« l'hypothèse « $\widehat{C}$ une $\mathcal{U}$-catégorie » est
essentielle » (marge de la page 104), « pour tout $C$ tel que
$\widehat{C}$ soit une $\mathcal{U}$-catégorie » (page 105) —, et d'un
alinéa biffé qui les discutait ; elles sont satisfaites pour $C$
essentiellement petite, ce que nous supposons. L'équivalence (8.2.8) est
la formule $T_{R}(\widehat{J}) \simeq \mathrm{Hom}(J^{\circ}, T_{R}(\mathrm{Ens}))$
de la page 65.} Si $T$ a un ensemble $I$ d'objets de base, on a un carré
commutatif (8.2.9)

\begin{tikzcd}[column sep=large]
T(\widehat{C}) \arrow[r, "\alpha"] \arrow[d, "b_{\widehat{C}}"'] & \mathrm{Hom}(C^{\circ}, S_{T}) \arrow[d, "{\mathrm{Hom}(C^{\circ},\, b_{\mathrm{Ens}})}"] \\
\widehat{C}^{I} \arrow[r, "\cong"'] & \mathrm{Hom}(C^{\circ}, \mathrm{Ens})^{I}
\end{tikzcd}

et par (8.2.5) \emph{le carré}

\begin{tikzcd}[column sep=large]
T(C) \arrow[r, hook] \arrow[d, "b_{C}"'] & \mathrm{Hom}(C^{\circ}, S_{T}) \arrow[d, "{\varphi \mapsto (\beta_{i} \circ \varphi)_{i \in I}}"] \\
C^{I} \arrow[r, hook] & \widehat{C}^{I}
\end{tikzcd}

\emph{(8.2.10) est 2-cartésien}, $\beta_{i} = b_{i,\mathrm{Ens}} : S_{T} \to
\mathrm{Ens}$ étant les foncteurs « $i$-ème objet de base ». En termes
imagés : se donner une structure d'espèce $T$ sur $(E_{i}) \in C^{I}$
revient à se donner, fonctoriellement en $Z \in \mathrm{Ob}\,C$, une
structure d'espèce $T$ sur la famille d'ensembles $(\mathrm{Hom}(Z,
E_{i}))_{i \in I}$.\footnote{C'est la « définition par les points » d'une
structure algébrique dans une catégorie : un groupe dans $C$ est un objet
$G$ avec une structure de groupe sur chaque $\mathrm{Hom}(Z, G)$,
naturelle en $Z$. L'énoncé (8.2.10) dit exactement quand une telle donnée
provient de $C$ : quand les objets sous-jacents sont représentables ; et
sa forme générale, pour toute théorie à limites projectives, est ce que la
page 67 avait déjà établi. Le manuscrit écrit l'indice « $i \in Z$ » au
lieu de « $i \in I$ » dans la première famille.} \emph{Donc, lorsque
$\alpha$ est une équivalence, $T(C)$ est connu à équivalence près quand on
connaît}
\begin{equation}
S = S_{T} = T(\mathrm{Ens}) \quad \text{avec les foncteurs} \quad
\beta_{i} : S_{T} \to \mathrm{Ens}, \tag{8.2.11}
\end{equation}
\emph{comme la sous-catégorie pleine de $\mathrm{Hom}(C^{\circ}, S_{T})$ des
$\varphi$ tels que les $\beta_{i} \circ \varphi$ soient représentables.}

\pagerange{105}{106}
\subsection*{8.2.12 : reconstruire $R$ à partir de $S$}

Intrinsèquement, c'est la sous-catégorie pleine des $\varphi : C^{\circ} \to
S_{T}$ tels que $\beta_{\mathrm{Ens}} \circ \varphi$ soit représentable pour
\emph{toute} construction $\beta \in R^{\lambda}_{T}$. Le foncteur
\begin{equation}
\beta \longmapsto \beta_{\mathrm{Ens}} : R = R^{\lambda}_{T} \longrightarrow
\mathrm{Hom}(S_{T}, \mathrm{Ens}) = \widehat{S^{\circ}_{T}} \tag{8.2.13}
\end{equation}
est le foncteur évident $R \to \mathrm{Hom}(\mathrm{Hom}_{\lambda}(R,
\mathrm{Ens}), \mathrm{Ens})$ (8.2.14). Or $S = \mathrm{Hom}_{\lambda}(R,
\mathrm{Ens})$, sous-catégorie pleine de $\widehat{R^{\circ}}$, \emph{contient
$R^{\circ}$}, identifié aux foncteurs représentables de $R$ dans $\mathrm{Ens}$,
qui commutent aux limites :
\begin{equation}
R^{\circ} \subset S \subset \widehat{R^{\circ}}, \tag{8.2.15}
\end{equation}
et l'accouplement $\pi : R \times S \to \mathrm{Ens}$ est
\begin{equation}
\pi(c, \xi) = \xi(c) = \mathrm{Hom}_{\widehat{R^{\circ}}}(c, \xi) =
\mathrm{Hom}_{S}(c, \xi), \tag{8.2.16}
\end{equation}
de sorte que (8.2.13) est le composé de deux foncteurs pleinement fidèles
\begin{equation}
R \hookrightarrow S^{\circ} \hookrightarrow
\widehat{S^{\circ}} = \mathrm{Hom}(S, \mathrm{Ens}), \tag{8.2.17}
\end{equation}
le premier étant (8.2.15), le second Yoneda ; donc (8.2.13) est pleinement
fidèle. \emph{Ainsi $R$ est connu, à équivalence près, quand on connaît $S$
et l'image de $R$ dans $\widehat{S^{\circ}}$.}

\textbf{Pour résumer.} \emph{Si $T$ est un $\lambda$-type avec
$\overrightarrow{\lambda} = \emptyset$, alors $S = T(\mathrm{Ens})$ est muni
d'une sous-catégorie pleine $\Sigma \subset S$, formée des $\varphi \in
\mathrm{Ob}\,S$ tels que pour toute $\lambda$-catégorie $C$ et tout $\xi \in
T(C)$, définissant $\tilde\xi : C^{\circ} \to S$, $x \mapsto \mathrm{Hom}(x,
\xi)$, le foncteur composé $x \mapsto \mathrm{Hom}_{S}(\varphi, \tilde\xi(x))$
soit représentable ; $\Sigma$ est stable dans $S$ par limites inductives
de type $\lambda$, parce que $C$ est stable dans $\widehat{C}$ par limites
projectives ; et pour toute $C$ on a des $\lambda$-foncteurs $\Sigma^{\circ}
\to \underline{\mathrm{Hom}}(T(C), C)$ et $T(C) \to \mathrm{Hom}_{\lambda}(
\Sigma^{\circ}, C)$.} La phrase suivante — « on voit alors que l'on trouve
ainsi une \ldots » — s'arrête là, et la copie au net avec
elle.\footnote{Le manuscrit écrit « tout $\xi \in C_{\lambda}$ » là où le
sens exige $\xi \in T(C)$, ce que la page 68 écrit ; nous avons rétabli.
Ce que la phrase interrompue allait dire est la page 68 : une équivalence
$\Sigma^{\circ} \simeq R$, i.e. \emph{la catégorie modelante se retrouve
dans les modèles ensemblistes comme la sous-catégorie des objets
« représentables en toute construction »}. Pour $\lambda$ = limites
finies c'est $\Sigma = S_{\mathrm{pf}}$ et la dualité de Gabriel--Ulmer
(page 77) ; en général $R^{\circ} \subset \Sigma$ est clair (page 68, note)
et l'égalité est ce que la copie n'a pas eu le temps d'écrire. La
stabilité de $\Sigma$ par limites inductives de type $\lambda$ est celle de
la page 68, avec la même démonstration.}

\pagerange{108}{111}
\section*{Topos classifiant, structure universelle, topos modulaire (pages 108 à 111)}

\pagerange{108}{109}
\subsection*{L'image inverse de la structure universelle}

Un feuillet au crayon reprend, pour les théories à limites finies, la
question de la page 17 depuis le topos. Soit $C$ une catégorie
essentiellement petite à limites projectives finies, $T$ l'espèce de
structures que $C$ définit, $\widehat{C} = B_{T} = B$ le topos classifiant,
$O_{B}$ la structure universelle dans $B$ — le foncteur d'inclusion $C \to
\widehat{C}$, vu comme objet de $T(B) = \mathrm{Hom}_{\mathrm{lex}}(C, B)$ —, et
$E$ un topos. Une $T$-structure $O_{E} \in T(E) \simeq \mathrm{Hom}_{
\mathrm{lex}}(C, E)$, $O_{E} = f_{0}$, correspond à un morphisme de topos $f :
E \to B_{T}$, et
\[
O_{E} \simeq f^{*}(O_{B}), \qquad f_{0} = f^{*} \mid C.
\]
Pour $U \in \mathrm{Ob}\,E$ et $X \in \mathrm{Ob}\,B$, on définit une
application
\begin{equation}
\mathrm{Hom}_{E}(U, f^{*}X) \longrightarrow \mathrm{Hom}_{T(\mathrm{Ens})}
(O_{B}(X), O_{E}(U)), \tag{1}
\end{equation}
où $O_{E}(U) \in T(\mathrm{Ens})$ est l'image de $O_{E}$ par le foncteur exact
à gauche $Z \mapsto \mathrm{Hom}_{E}(U, Z)$, et $O_{B}(X)$ de même dans
$\widehat{C}$ : on a $\mathrm{Hom}_{E}(U, f^{*}X) \to \mathrm{Hom}(O_{E}(f^{*}X),
O_{E}(U)) \to \mathrm{Hom}(O_{B}(X), O_{E}(U))$, la seconde flèche par le
morphisme $O_{B}(X) \to O_{E}(f^{*}X)$ de $T(\mathrm{Ens})$ ; ici $T(\mathrm{Ens})
= \mathrm{Ind}(C^{\circ})$. \emph{Pour $X \in \mathrm{Ob}\,C$, (1) est une
bijection} : le second membre est $\mathrm{Hom}_{\mathrm{Ind}(C^{\circ})}(X,
O_{E}(U))$, et
\begin{gather*}
O_{B}(X) = \mathrm{Hom}_{C}(X, -), \qquad O_{E}(U) = \bigl(Y \mapsto
\mathrm{Hom}_{E}(U, f_{0}(Y))\bigr), \\
\mathrm{Hom}(O_{B}(X), O_{E}(U)) = O_{E}(U)(X) = \mathrm{Hom}_{E}(U, f_{0}(X)),
\end{gather*}
par Yoneda, ce qui est bien $\mathrm{Hom}_{E}(U, f^{*}X)$.\footnote{C'est la
description explicite de l'image inverse par un morphisme de topos vers
un topos de préfaisceaux : $f^{*}(X)$ pour $X$ représentable est $f_{0}(X)$,
et pour $X$ quelconque une limite inductive de tels. L'énoncé « Prouver que
si $X \in \mathrm{Ob}\,C$, l'application précédente est une bijection » est
posé, et la ligne qui suit le prouve. La formule $T(\mathrm{Ens}) =
\mathrm{Ind}(C^{\circ})$, ajoutée à l'encre bleue, et $\mathrm{Pro}(C) \simeq
T(\mathrm{Ens})^{\circ}$ plus bas, sont le cas des limites finies de (8.2.15)
et de la page 77.}

On en profite pour calculer $f^{*}(\mathcal{F})$ pour $\mathcal{F}$ un champ
sur $B$. Il revient au même de se donner $\mathcal{F}$ ou sa restriction
à $C$, une catégorie fibrée sur $C$ (avec des données de descente que la
page laisse en blanc). Si $B' \subset B$ est un sous-topos, correspondant
à une topologie $\tau$ sur $C$ et à une sous-espèce $T'$ de $T$, et si
$O_{E}$ est de cette espèce, i.e. $f$ se factorise par $B'$, i.e. $f_{0}$
est continu pour $\tau$, alors un champ $\mathcal{F}'$ sur $B'$ est un
champ $\mathcal{F}_{0}$ sur $B$ satisfaisant une condition relative à
$\tau$, et $f'^{*}(\mathcal{F}')$ se calcule en termes de $f_{0}$ et
$\mathcal{F}'$. On commence par prolonger $\mathcal{F}_{0}$, catégorie fibrée
sur $C$, à $\mathrm{Pro}(C) \simeq T(\mathrm{Ens})^{\circ}$ ; alors
$\mathcal{F}_{0}(O_{E}(U))$ est défini, covariant en l'objet $O_{E}(U) \in
T(\mathrm{Ens})$ ; pour $U$ variable, $U \mapsto \mathcal{F}_{0}(O_{E}(U))$ est
une catégorie fibrée sur $E$, image inverse de $\mathcal{F}_{0}$ par $O_{E} :
E^{\circ} \to T(\mathrm{Ens}) \cong \mathrm{Pro}(C)^{\circ}$, et $f^{*}
(\mathcal{F})$ est le champ associé.\footnote{Le feuillet est au crayon,
très pâle, et la moitié des mots de liaison est perdue ; ce paragraphe est
la reconstruction de son argument, et il est correct : l'image inverse
d'un champ sur un topos de préfaisceaux se calcule en prolongeant sa
restriction à $C$ aux pro-objets (par limites projectives, ce que
« covariant en l'objet de $T(\mathrm{Ens})$ » désigne), en composant avec
le foncteur $U \mapsto O_{E}(U)$, et en prenant le champ associé. C'est le
cas $\mathrm{Cat}$-valué de (8.2.10).}

\pagerange{111}{111}
\subsection*{Topos modulaire pour des structures sur objets d'un topos fixé}

Un feuillet quadrillé, sous le titre « Topos modulaire pour des structures
algébriques sur objets variés d'un topos fixé $X$ », pose le problème
relatif. On se donne des objets $E_{i} \in X$, $i \in I$, et une espèce de
structure $T$ de topos classifiant $B$ muni de ses objets de base $b_{i}
\in B$. On s'intéresse, pour tout topos $X'$ sur $X$, aux structures
d'espèce $T$ dans $X'$ — les morphismes $f : X' \to B$ — dont les objets de
base sont les images inverses $E'_{i}$ des $E_{i}$ : $f^{*}(b_{i}) \simeq
E'_{i}$. La catégorie de ces objets pour $X'$ fixé, $X' \to X$ et $X' \to B$
variables, est le produit 2-fibré de $\underline{\mathrm{Hom}}_{\mathrm{Top}}(X',
X)$ et $\underline{\mathrm{Hom}}_{\mathrm{Top}}(X', B)$ au-dessus de $\underline{
\mathrm{Hom}}_{\mathrm{Top}}(X', B_{I})$ ; \emph{le topos modulaire cherché est
donc le produit 2-fibré $B \times_{B_{I}} X$ dans la 2-catégorie des
topos},

\begin{tikzcd}
B \arrow[dr, "b"'] & & X \arrow[dl, "(E_{i})"] \\
& B_{I} = \widehat{(\mathrm{Ens}_{f})^{\circ}}{}^{I} &
\end{tikzcd}

$B_{I}$ étant le classifiant de $I$ objets, « on n'a pas encore prouvé que
cela marche ».\footnote{Cela marche : les produits fibrés de topos
existent (page 21, SGA 4 IV 8.2), et $B \times_{B_{I}} X$ classifie
exactement les $T$-structures sur $X'$ dont les objets de base sont les
$E'_{i}$ — c'est l'exemple type de ce que le produit fibré de classifiants
classifie. Le mot « modulaire » du titre, et « catégorie modulaire » à la
ligne suivante, est celui que la page 98 écrit « modelaire » : le topos
des modèles.} Dans le cas où $T$ est définie par limites finies, de
catégorie modelante $R$ avec $B = \widehat{R}$ et $I \to \mathrm{Ob}\,R$, on
s'intéresse aux foncteurs exacts à gauche $R \to X'$ transformant les
$r_{i}$ en les $E'_{i}$ ; soit $\mathcal{S}$ le champ sur $X$ des structures
$(X', (E'_{i}), f : R \to X')$ avec $f(r_{i}) \simeq E'_{i}$, et $\underline{I}
_{X}$ le champ dont c'est la fibre — la page s'arrête sur cette
définition, avec en marge, barré, « champ en catégories associé aux
catégories \ldots ».\footnote{Le programme est celui des \emph{topos
modulaires} — au sens de : topos des modules d'une structure — relatifs à
une base, qui est la forme que prendra, chez Hakim, la construction des
schémas relatifs à un topos annelé, et, en logique catégorique, la
notion de topos classifiant \emph{relatif}. Le dossier s'arrête au moment
où la question est posée dans les bons termes.}

\end{document}
